% Les mécanismes fondamentaux de la reproduction sexuée chez les Mammifères — SVTEEHB Tle TI — Cours signé OnBuch+
% Programme officiel MINESEC. Compiler avec : tectonic N03-reproduction-sexuee-mammiferes.tex
\newcommand{\DOCMATIERE}{SVTEEHB}
\newcommand{\DOCNIVEAU}{Tle TI}
\newcommand{\DOCMODULE}{Module I : Le Monde Vivant}
\newcommand{\DOCLECON}{N03}
\newcommand{\DOCTITRE}{Les mécanismes fondamentaux de la reproduction sexuée chez les Mammifères}
% OnBuch+ — préambule « cours signés » ENRICHI (compilé avec tectonic / XeLaTeX)
% Système visuel « Pop » d'OnBuch+ : fond crème (jamais blanc), bordures encre
% épaisses, ombres « dures » décalées, titres Archivo Black en capitales.
% Version MATHÉMATIQUES : graphes (pgfplots/tikz), boîtes pédagogiques
% (définition, propriété, méthode, attention, le savais-tu, exercice résolu,
% exercice, TP, bilan), page de garde et signature OnBuch+ ancrée en bas.
%
% Macros attendues dans main.tex AVANT \input{preamble} :
%   \DOCMATIERE  \DOCNIVEAU  \DOCMODULE  \DOCLECON (numéro)  \DOCTITRE
\documentclass[11pt]{article}

\usepackage{fontspec}
\usepackage[french]{babel}
\usepackage[a4paper,margin=2cm,top=3.4cm,bottom=2.8cm,headheight=1.4cm,footskip=1.3cm]{geometry}
\usepackage{amsmath,amssymb,amsfonts}
\usepackage{mathtools} % \xmapsto, \coloneqq... souvent utilisés par les modèles
\usepackage{cancel} % simplifications barrées (bilans de forces)
% \abs{z} (module, valeur absolue) et \norm{u} (norme), utilisés par les modèles
\providecommand{\abs}{}\let\abs\relax\DeclarePairedDelimiter{\abs}{\lvert}{\rvert}
\providecommand{\norm}{}\let\norm\relax\DeclarePairedDelimiter{\norm}{\lVert}{\rVert}
\usepackage[table]{xcolor} % [table] : fournit aussi \rowcolors (alternance de lignes)
\usepackage{pagecolor}
\usepackage{tikz}
\usetikzlibrary{arrows.meta,positioning,calc,shapes.geometric,shapes.misc,shapes.arrows,decorations.pathmorphing,decorations.pathreplacing,decorations.markings,patterns,shadows,mindmap,trees,fit,backgrounds,matrix,intersections,angles,quotes}
\usepackage{pgfplots}
\usepackage[version=4]{mhchem} % équations nucléaires \ce{^{235}_{92}U}
\usepackage[european,siunitx]{circuitikz} % schémas électriques
\pgfplotsset{compat=1.18}
\usepgfplotslibrary{fillbetween,groupplots} % aires entre courbes (calcul intégral)
\usepackage{enumitem}
\usepackage{fancyhdr}
% [most] charge toutes les bibliothèques tcolorbox (listings, minted, raster,
% documentation...) et gonfle inutilement la mémoire TeX sur un document long
% (~300 boîtes) : on ne charge que ce qui est réellement utilisé.
\usepackage[skins,breakable]{tcolorbox}
\usepackage{graphicx}
\usepackage{colortbl}
\usepackage{booktabs}
\usepackage{tabularx}
\usepackage{multirow}
\usepackage{diagbox} % case d'en-tête diagonale des tableaux à double entrée
% Colonnes extensibles centrée / gauche / droite (C, L, R), souvent utilisées par les modèles
\newcolumntype{C}{>{\centering\arraybackslash}X}
\newcolumntype{L}{>{\raggedright\arraybackslash}X}
\newcolumntype{R}{>{\raggedleft\arraybackslash}X}
\usepackage{array}
\usepackage{multicol}
\usepackage{siunitx}
% parse-numbers=false : accepte aussi \SI{2,5 \times 10^{-3}}{...}
\sisetup{locale=FR,output-decimal-marker={,},group-minimum-digits=4,parse-numbers=false}
\DeclareSIUnit\ohm{\ensuremath{\symup{\Omega}}} % U+2126 absent de la police
\DeclareSIUnit\division{div} % réglages d'oscilloscope (ms/div, V/div)
\DeclareSIUnit\tr{tr} % tours (vitesse de rotation, ex. \SI{1500}{\tr\per\minute})
\DeclareSIUnit\molar{M} % concentration molaire (ex. \SI{3}{\molar}, \SI{1}{\micro\molar})
\DeclareSIUnit\an{an} % année (le modèle confond parfois avec \year, un compteur TeX) — ex. \SI{5}{\kilo\meter\per\an}
\DeclareSIUnit\FCFA{FCFA} % franc CFA (ex. \SI{500}{\FCFA\per\kilo\gram})
\DeclareSIUnit\degreeBrix{\textdegree Brix} % teneur en sucre d'un sirop/jus (réfractométrie)
\DeclareSIUnit\month{mois} % le modèle utilise parfois \month, un compteur TeX, comme unité de durée
\DeclareSIUnit\microliter{\micro\liter} % le modèle écrit parfois \microliter en un seul token
\DeclareSIUnit\million{millions} % dénombrement (ex. \SI{15}{\million\per\mL} spermatozoïdes)
\usepackage{qrcode}
% \degree : commande fréquemment utilisée par les modèles pour un angle ou une
% température en dehors de \SI{}{} (non fournie par les paquets chargés).
\providecommand{\degree}{^{\circ}}
% \qed (fin de démonstration), utilisé par les modèles sans amsthm
\providecommand{\qed}{\hfill$\square$}
% \barre : double barre des valeurs interdites dans un tableau de variations (tkz-tab)
\providecommand{\barre}{\ensuremath{\|}}
% environnement proof (amsthm non chargé) utilisé par les modèles
\ifdefined\proof\else\newenvironment{proof}[1][Démonstration]{\par\noindent\textit{#1.}\ }{\qed\par}\fi
\usepackage{lastpage}
\usepackage{hyperref}
\hypersetup{hidelinks}

% ============================================================
% Palette « Pop » OnBuch+ — mêmes hex que la classe P (plus_ui.dart)
% ============================================================
\definecolor{poporange}{HTML}{F59321}
\definecolor{popdark}{HTML}{E07A0C}
\definecolor{popcream}{HTML}{FFF6E8}
\definecolor{popink}{HTML}{1C1714}
\definecolor{poppurple}{HTML}{7A5AE0}
\definecolor{popgreen}{HTML}{2E9E6A}
\definecolor{poppink}{HTML}{E0457B}
\definecolor{popblue}{HTML}{2F7DE1}
\definecolor{popgold}{HTML}{C98A12}
\definecolor{popmuted}{HTML}{8A7F76}
\definecolor{popline}{HTML}{EADFD2}
\definecolor{popgray}{HTML}{8A7F76}
\colorlet{popgrey}{popgray}
% Alias tolérants (le modèle nomme parfois les couleurs différemment)
\colorlet{popwhite}{white}
\colorlet{popblack}{popink}
\colorlet{popred}{poppink}
\colorlet{popyellow}{popgold}
% Teintes claires (fonds de tableaux/encadrés) — le modèle nomme parfois
% les couleurs avec un suffixe L (ex. popblueL) sans les avoir définies.
\colorlet{poporangeL}{poporange!20}
\colorlet{popdarkL}{popdark!20}
\colorlet{poppurpleL}{poppurple!20}
\colorlet{popgreenL}{popgreen!20}
\colorlet{poppinkL}{poppink!20}
\colorlet{popblueL}{popblue!20}
\colorlet{popgoldL}{popgold!20}
\colorlet{popredL}{popred!20}
\colorlet{popgrayL}{popgray!20}
% Teintes claires pour fonds de tableaux / zones
\colorlet{popblueL}{popblue!10!white}
\colorlet{poporangeL}{poporange!14!white}
\colorlet{popgreenL}{popgreen!12!white}
\colorlet{poppurpleL}{poppurple!12!white}
\colorlet{poppinkL}{poppink!12!white}
\colorlet{popgoldL}{popgold!14!white}

\pagecolor{popcream}

% ============================================================
% Typographie
% ============================================================
\defaultfontfeatures{Ligatures=TeX}
\setmainfont{PlusJakartaSans-Medium.ttf}[Path=../../../fonts/,BoldFont=PlusJakartaSans-Bold.ttf]
% Maths en sans-serif (Fira Math) avec chiffres et lettres droites de Plus
% Jakarta, pour que les formules se fondent dans le texte.
\usepackage[math-style=ISO,bold-style=ISO]{unicode-math}
\setmathfont{FiraMath-Regular.otf}[Path=../../../fonts/]
\setmathfont{PlusJakartaSans-Medium.ttf}[Path=../../../fonts/,range={up/num,up/{Latin,latin},\mathup}]
\usepackage{icomma} % virgule décimale sans espace en mode math
% Caractères mathématiques Unicode absents des polices de texte (Plus Jakarta,
% Archivo Black) : rendus par leur équivalent mathématique, en texte comme en
% formule — sinon ils s'affichent en carré vide (ex. « Arithmétique dans ℤ »).
\usepackage{newunicodechar}
\newunicodechar{ℝ}{\ensuremath{\mathbb{R}}}
\newunicodechar{ℤ}{\ensuremath{\mathbb{Z}}}
\newunicodechar{ℂ}{\ensuremath{\mathbb{C}}}
\newunicodechar{ℕ}{\ensuremath{\mathbb{N}}}
\newunicodechar{ℚ}{\ensuremath{\mathbb{Q}}}
\newunicodechar{𝔻}{\ensuremath{\mathbb{D}}}
\newunicodechar{α}{\ensuremath{\alpha}}
\newunicodechar{β}{\ensuremath{\beta}}
\newunicodechar{γ}{\ensuremath{\gamma}}
\newunicodechar{δ}{\ensuremath{\delta}}
\newunicodechar{ε}{\ensuremath{\varepsilon}}
\newunicodechar{θ}{\ensuremath{\theta}}
\newunicodechar{λ}{\ensuremath{\lambda}}
\newunicodechar{μ}{\ensuremath{\mu}}
\newunicodechar{σ}{\ensuremath{\sigma}}
\newunicodechar{τ}{\ensuremath{\tau}}
\newunicodechar{φ}{\ensuremath{\varphi}}
\newunicodechar{ψ}{\ensuremath{\psi}}
\newunicodechar{ω}{\ensuremath{\omega}}
\newunicodechar{Σ}{\ensuremath{\Sigma}}
% majuscules produites par \MakeUppercase dans les titres (α-aminés -> Α)
\newunicodechar{Α}{\ensuremath{\alpha}}
\newunicodechar{Β}{\ensuremath{\beta}}
\newunicodechar{Γ}{\ensuremath{\Gamma}}
\newunicodechar{Δ}{\ensuremath{\Delta}}
\newunicodechar{Ε}{\ensuremath{\varepsilon}}
\newunicodechar{Θ}{\ensuremath{\Theta}}
\newunicodechar{Λ}{\ensuremath{\Lambda}}
\newunicodechar{Μ}{\ensuremath{\mu}}
\newunicodechar{Τ}{\ensuremath{\tau}}
\newunicodechar{Φ}{\ensuremath{\Phi}}
\newunicodechar{Ψ}{\ensuremath{\Psi}}
\newunicodechar{Ω}{\ensuremath{\Omega}}
\newunicodechar{↦}{\ensuremath{\mapsto}}
\newunicodechar{∈}{\ensuremath{\in}}
\newunicodechar{∉}{\ensuremath{\notin}}
\newunicodechar{∧}{\ensuremath{\wedge}}
\newunicodechar{⇔}{\ensuremath{\Leftrightarrow}}
\newunicodechar{⟺}{\ensuremath{\Longleftrightarrow}}
\newunicodechar{⇒}{\ensuremath{\Rightarrow}}
\newunicodechar{⟹}{\ensuremath{\Longrightarrow}}
\newunicodechar{∘}{\ensuremath{\circ}}
\newunicodechar{∪}{\ensuremath{\cup}}
\newunicodechar{∩}{\ensuremath{\cap}}
\newunicodechar{≡}{\ensuremath{\equiv}}
\newunicodechar{⊂}{\ensuremath{\subset}}
\newunicodechar{⊥}{\ensuremath{\perp}}
\newunicodechar{∃}{\ensuremath{\exists}}
\newunicodechar{∀}{\ensuremath{\forall}}
\newunicodechar{∛}{\ensuremath{\sqrt[3]{\,}}}
\newunicodechar{∜}{\ensuremath{\sqrt[4]{\,}}}
\newunicodechar{⃗}{\ensuremath{{}^{\to}}}
\newunicodechar{ⁿ}{\ifmmode^{n}\else\textsuperscript{\ensuremath{n}}\fi}
\newunicodechar{ˣ}{\ifmmode^{x}\else\textsuperscript{\ensuremath{x}}\fi}
\newunicodechar{ᵇ}{\ifmmode^{b}\else\textsuperscript{\ensuremath{b}}\fi}
\newunicodechar{ʳ}{\ifmmode^{r}\else\textsuperscript{\ensuremath{r}}\fi}
\newunicodechar{ᵘ}{\ifmmode^{u}\else\textsuperscript{\ensuremath{u}}\fi}
\newunicodechar{ʸ}{\ifmmode^{y}\else\textsuperscript{\ensuremath{y}}\fi}
\newunicodechar{ᵃ}{\ifmmode^{a}\else\textsuperscript{\ensuremath{a}}\fi}
\newunicodechar{ᵉ}{\ifmmode^{e}\else\textsuperscript{\ensuremath{e}}\fi}
\newunicodechar{ᵏ}{\ifmmode^{k}\else\textsuperscript{\ensuremath{k}}\fi}
\newunicodechar{ᵖ}{\ifmmode^{p}\else\textsuperscript{\ensuremath{p}}\fi}
\newunicodechar{ᴺ}{\ifmmode^{N}\else\textsuperscript{\ensuremath{N}}\fi}
\newunicodechar{ᶜ}{\ifmmode^{c}\else\textsuperscript{\ensuremath{c}}\fi}
\newunicodechar{ʰ}{\ifmmode^{h}\else\textsuperscript{\ensuremath{h}}\fi}
\newunicodechar{ᵠ}{\ifmmode^{q}\else\textsuperscript{\ensuremath{q}}\fi}
\newunicodechar{ₙ}{\ifmmode_{n}\else\textsubscript{\ensuremath{n}}\fi}
\newunicodechar{ₐ}{\ifmmode_{a}\else\textsubscript{\ensuremath{a}}\fi}
\newunicodechar{ᵢ}{\ifmmode_{i}\else\textsubscript{\ensuremath{i}}\fi}
\newunicodechar{ᵣ}{\ifmmode_{r}\else\textsubscript{\ensuremath{r}}\fi}
\newunicodechar{ₖ}{\ifmmode_{k}\else\textsubscript{\ensuremath{k}}\fi}
\newunicodechar{ᵦ}{\ifmmode_{\beta}\else\textsubscript{\ensuremath{\beta}}\fi}
\newunicodechar{ₓ}{\ifmmode_{x}\else\textsubscript{\ensuremath{x}}\fi}
\newfontfamily\popdisplay{ArchivoBlack-Regular.ttf}[Path=../../../fonts/]
\newfontfamily\poplogo{SpaceGrotesk-Bold.ttf}[Path=../../../fonts/]
\newfontfamily\popsemi{PlusJakartaSans-SemiBold.ttf}[Path=../../../fonts/]
\newfontfamily\popxbold{PlusJakartaSans-ExtraBold.ttf}[Path=../../../fonts/]
\newfontfamily\popxboldls{PlusJakartaSans-ExtraBold.ttf}[Path=../../../fonts/,LetterSpace=3.0]

\setlist[enumerate]{leftmargin=1.7em,itemsep=5pt,topsep=4pt}
\setlist[itemize]{leftmargin=1.7em,itemsep=4pt,topsep=4pt,label={\color{poporange}\textbullet}}
\setlength{\parindent}{0pt}
\setlength{\parskip}{5pt}
\renewcommand{\arraystretch}{1.25}

% pgfplots : style Pop par défaut
\pgfplotsset{
  every axis/.append style={axis line style={popink,line width=1pt},tick style={popink},
    grid style={popline,line width=0.5pt},label style={font=\small},tick label style={font=\footnotesize}},
  popcurve/.style={poporange,line width=1.8pt,no markers,smooth},
}
% Même style disponible dans un \draw TikZ simple (hors pgfplots)
\tikzset{popcurve/.style={poporange,line width=1.8pt}}
\tikzset{popgrid/.style={popline,very thin}}
\colorlet{popcurve}{poporange} % le nom du style est parfois utilisé comme couleur

% ============================================================
% Pill / squiggle
% ============================================================
\newcommand{\poppill}[3]{%
  \tikz[baseline=-0.55ex]\node[draw=popink,line width=1.2pt,rounded corners=2.6pt,%
    fill=#2,inner xsep=6.5pt,inner ysep=3pt]{%
    \popxboldls\fontsize{7}{7}\selectfont\color{#3}\MakeUppercase{#1}};%
}
\newcommand{\popsquiggle}[1]{%
  \tikz[baseline=-0.4ex]{
    \draw[#1,line width=2.6pt,line cap=round]
      plot [smooth] coordinates {(0,0.09) (0.34,0.01) (0.68,0.21) (1.02,0.01) (1.36,0.10)};
  }%
}
% Mot-clé surligné (marqueur orange sous le texte)
\newcommand{\cle}[1]{{\popxbold\color{popdark}#1}}

% ============================================================
% En-tête / pied
% ============================================================
\pagestyle{fancy}
\fancyhf{}
\renewcommand{\headrulewidth}{0pt}
\renewcommand{\footrulewidth}{0pt}
\lhead{\raisebox{-0.15ex}{\poplogo\fontsize{13}{13}\selectfont\color{popink}OnBuch\color{poporange}+}}
\rhead{\poppill{\DOCMATIERE\ \textbullet\ \DOCNIVEAU\ \textbullet\ Leçon \DOCLECON}{white}{popink}}
\lfoot{\popxbold\fontsize{7.6}{7.6}\selectfont\color{popmuted}\MakeUppercase{Cours signé OnBuch+}\ \textbullet\ {\popsemi\DOCTITRE}}
\rfoot{\tikz[baseline=-0.6ex]\node[rounded corners=8.5pt,fill=popink,minimum height=17pt,minimum width=17pt,inner xsep=5pt,inner sep=0]{\popxbold\fontsize{8}{8}\selectfont\color{popcream}\thepage\,/\,\pageref*{LastPage}};}

% ============================================================
% Titres
% ============================================================
\newcommand{\chaptitre}[2]{%
  \begin{tcolorbox}[enhanced,colback=poporange,colframe=popink,boxrule=2.6pt,arc=11pt,
      drop shadow=popink,left=15pt,right=15pt,top=12pt,bottom=11pt,before skip=3pt,after skip=18pt]
    \poppill{#2}{popink}{popcream}\par\vspace{9pt}
    {\raggedright\hyphenpenalty=10000\exhyphenpenalty=10000\popdisplay\fontsize{22}{24}\selectfont\color{popcream}\MakeUppercase{#1}\par}
  \end{tcolorbox}
}
% Section numérotée : pastille orange avec le numéro + titre Archivo
\newcounter{popsec}
\newcommand{\coursec}[1]{%
  \stepcounter{popsec}\setcounter{popsub}{0}%
  \par\vspace{16pt}\noindent
  \begin{minipage}{\linewidth}
  \tikz[baseline=-0.7ex]\node[draw=popink,line width=1.6pt,fill=poporange,rounded corners=4pt,minimum size=22pt,inner sep=0]{\popdisplay\fontsize{12}{12}\selectfont\color{popcream}\thepopsec};%
  \hspace{8pt}{\popdisplay\fontsize{15}{17}\selectfont\color{popink}\MakeUppercase{#1}}\par
  \vspace{4pt}\noindent{\color{poporange}\rule{36pt}{3.6pt}}
  \end{minipage}\par\nopagebreak\vspace{8pt}
}
\newcounter{popsub}
\newcommand{\courssub}[1]{%
  \stepcounter{popsub}%
  \par\vspace{9pt}\noindent{\popxbold\fontsize{11.5}{13}\selectfont\color{popink}{\color{poporange}\thepopsec.\thepopsub}\hspace{6pt}#1}\par\nopagebreak\vspace{4pt}
}

% ============================================================
% Boîtes pédagogiques — même recette : fond blanc, bordure épaisse colorée,
% ombre dure encre, badge plein. Titre optionnel [#1].
% ============================================================
\tcbset{popbase/.style={breakable,enhanced,colback=white,boxrule=2.3pt,arc=9pt,
  shadow={3pt}{-3pt}{0pt}{popink},left=14pt,right=14pt,top=11pt,bottom=11pt,before skip=11pt,after skip=15pt,
  fontupper=\popsemi\fontsize{10.6}{14.5}\selectfont\color{popink},
  fontlower=\popsemi\fontsize{10.6}{14.5}\selectfont\color{popink}}}
% Coche dessinée (le glyphe ✓ n'existe pas dans les polices du document)
\newcommand{\popcheck}[1]{\tikz[baseline=-0.5ex]\draw[#1,line width=1.6pt,line cap=round,line join=round](-0.13,0.0)--(-0.04,-0.09)--(0.13,0.1);}
\providecommand{\coche}{\popcheck{popgreen}}
\providecommand{\resultat}[1]{\par\smallskip\noindent\fcolorbox{popgreen}{popgreenL}{\parbox{\dimexpr\linewidth-2\fboxsep-2\fboxrule\relax}{\textbf{Résultat.} #1}}\par\smallskip}
\newcommand{\popboxhead}[3]{\poppill{#1}{#2}{white}\ifx\relax#3\relax\else\hspace{6pt}{\popxbold\fontsize{10.6}{12}\selectfont\color{popink}#3}\fi\par\vspace{7pt}}

\newtcolorbox{aretenir}[1][]{popbase,colframe=popgreen,before upper={\popboxhead{À retenir}{popgreen}{#1}}}
\newtcolorbox{exemplebox}[1][]{popbase,colframe=poporange,before upper={\popboxhead{Exemple}{poporange}{#1}}}
\newtcolorbox{definition}[1][]{popbase,colframe=popblue,before upper={\popboxhead{Définition}{popblue}{#1}}}
\newtcolorbox{propriete}[1][]{popbase,colframe=poppurple,before upper={\popboxhead{Propriété}{poppurple}{#1}}}
\newtcolorbox{methode}[1][]{popbase,colframe=popgold,before upper={\popboxhead{Méthode}{popgold}{#1}}}
\newtcolorbox{attention}[1][]{popbase,colframe=poppink,before upper={\popboxhead{Attention}{poppink}{#1}}}
\newtcolorbox{experience}[1][]{popbase,colframe=popblue,colback=popblueL,before upper={\popboxhead{Expérience}{popblue}{#1}}}
% « Le savais-tu ? » : carte violette pleine (contexte, culture, Cameroun)
\newtcolorbox{savaistu}[1][]{popbase,colback=poppurple,colframe=popink,
  fontupper=\popsemi\fontsize{10.4}{14.2}\selectfont\color{white},
  before upper={\poppill{Le savais-tu\,?}{white}{poppurple}\ifx\relax#1\relax\else\hspace{6pt}{\popxbold\color{white}#1}\fi\par\vspace{7pt}}}
% Exercice résolu : énoncé en haut, \tcblower puis la solution sur fond crème
\newcounter{popexo}\newcounter{popexr}
\newtcolorbox{exoresolu}[1][]{popbase,colframe=poporange,colbacklower=popcream,
  segmentation style={popink,line width=1.2pt,dashed},
  before upper={\stepcounter{popexr}\popboxhead{Exercice résolu \thepopexr}{poporange}{#1}},
  before lower={\poppill{Solution}{popink}{popcream}\par\vspace{6pt}}}
% Exercice (non résolu en place) — la correction est dans la partie Corrigés
\newtcolorbox{exercice}[1][]{popbase,colframe=popink,
  before upper={\stepcounter{popexo}\popboxhead{Exercice \thepopexo}{popink}{#1}}}
% Corrigé
\newtcolorbox{corrige}[1][]{popbase,colframe=popgreen,colback=popgreenL,
  before upper={\popboxhead{Corrigé}{popgreen}{#1}}}
% Objectifs / prérequis / bilan
\newtcolorbox{objectifs}{popbase,colframe=popink,colback=poporangeL,
  before upper={\popboxhead{Objectifs de la leçon}{popink}{}}}
\newtcolorbox{prerequis}{popbase,colframe=popink,colback=popblueL,
  before upper={\popboxhead{Prérequis}{popblue}{}}}
\newtcolorbox{bilan}{popbase,colframe=popink,colback=white,boxrule=2.8pt,
  before upper={\popboxhead{Fiche bilan}{poporange}{L'essentiel à connaître pour l'examen}}}
% Figure encadrée avec légende
\newtcolorbox{popfigure}[1][]{enhanced,breakable=false,colback=white,colframe=popline,boxrule=1.4pt,arc=8pt,
  left=10pt,right=10pt,top=10pt,bottom=8pt,before skip=10pt,after skip=12pt,halign=center,
  lower separated=false,halign lower=center,
  fontlower=\popsemi\fontsize{9}{11.5}\selectfont\color{popmuted},
  #1}
\newcommand{\legende}[1]{\par\vspace{6pt}{\popsemi\fontsize{9}{11.5}\selectfont\color{popmuted}#1}}

% ============================================================
% Signature OnBuch+
% ============================================================
\newcommand{\poplogoinline}[1]{{\poplogo\fontsize{#1}{#1}\selectfont\color{popink}OnBuch\color{poporange}+}}
\newcommand{\signatureonbuch}{%
  \begin{tcolorbox}[enhanced,colback=popink,colframe=popink,boxrule=2.2pt,arc=10pt,
      drop shadow=poporange,left=14pt,right=14pt,top=10pt,bottom=10pt,before skip=0pt,after skip=0pt]
    \tikz[baseline=-0.7ex]\node[circle,fill=popgreen,draw=popcream,line width=1.3pt,minimum size=22pt,inner sep=0]{\popcheck{white}};%
    \hspace{9pt}\begin{minipage}[c]{0.62\linewidth}
      {\popxbold\fontsize{12}{13}\selectfont\color{popcream}Signé\ \textbullet\ {\poplogo OnBuch\color{poporange}+}}\par\vspace{2pt}
      {\raggedright\popsemi\fontsize{8}{10.5}\selectfont\color{popline}\MakeUppercase{Équipe pédagogique OnBuch+}\\\MakeUppercase{Conforme au programme officiel MINESEC}\par}
    \end{minipage}\hfill
    {\poplogo\fontsize{22}{22}\selectfont\color{popcream}OnBuch\color{poporange}+}
  \end{tcolorbox}%
}

% ============================================================
% Page de garde
% #1 titre · #2 matière · #3 niveau · #4 module · #5 n° leçon · #6 durée ·
% #7 description / objectifs (texte ou itemize)
% ============================================================
\newcommand{\pagegarde}[7]{%
  \thispagestyle{empty}
  \newgeometry{a4paper,margin=1.7cm,top=1.6cm,bottom=1.4cm}%
  \begin{tikzpicture}[remember picture,overlay]
    \fill[poporange!18] ([xshift=-3.2cm,yshift=-3.6cm]current page.north east) circle (4.6cm);
    \fill[poppurple!14] ([xshift=2.4cm,yshift=7.6cm]current page.south west) circle (3.8cm);
  \end{tikzpicture}%
  {\poplogo\fontsize{30}{30}\selectfont\color{popink}OnBuch\color{poporange}+}\hfill
  \raisebox{0.6ex}{\poppill{Cours signé \textbullet\ Premium}{popink}{popcream}}\par
  \vspace{1pt}\popsquiggle{poppurple}\par
  \vspace{8pt}
  \begin{tcolorbox}[enhanced,colback=poporange,colframe=popink,boxrule=2.8pt,arc=13pt,
      drop shadow=popink,left=18pt,right=18pt,top=12pt,bottom=13pt,after skip=10pt]
    \poppill{#2 \textbullet\ #3}{popink}{popcream}\hspace{5pt}\poppill{#4}{white}{popink}\par\vspace{8pt}
    {\popdisplay\fontsize{14}{14}\selectfont\color{popink}LEÇON #5}\par\vspace{4pt}
    {\raggedright\hyphenpenalty=10000\exhyphenpenalty=10000\popdisplay\fontsize{25}{27}\selectfont\color{popcream}\MakeUppercase{#1}\par}
    \vspace{8pt}
    {\popxbold\fontsize{9.5}{9.5}\selectfont\color{popink}\MakeUppercase{Durée conseillée : #6}}
  \end{tcolorbox}
  \begin{tcolorbox}[enhanced,colback=white,colframe=popink,boxrule=2.2pt,arc=10pt,
      drop shadow=popink,left=16pt,right=16pt,top=10pt,bottom=10pt,after skip=8pt]
    \poppill{Dans cette leçon}{poppurple}{white}\par\vspace{5pt}
    \popsemi\fontsize{9.3}{12.6}\selectfont\color{popink}#7
  \end{tcolorbox}
  \noindent\begin{minipage}[t]{0.487\linewidth}
    \begin{tcolorbox}[enhanced,colback=popgreen,colframe=popink,boxrule=2.2pt,arc=10pt,
        drop shadow=popink,left=11pt,right=11pt,top=10pt,bottom=10pt,height=3.1cm,valign=center]
      \tcbox[colback=white,colframe=popink,boxrule=1.3pt,arc=5pt,boxsep=0pt,nobeforeafter,
        left=3pt,right=3pt,top=3pt,bottom=3pt,valign=center]{\qrcode[height=1.9cm]{https://play.google.com/store/apps/details?id=com.luvvix.onbuch&hl=fr}}%
      \hspace{8pt}\begin{minipage}[c]{0.5\linewidth}
        {\popdisplay\fontsize{11}{12.5}\selectfont\color{white}\MakeUppercase{Télécharge OnBuch}}\par\vspace{4pt}
        {\raggedright\popsemi\fontsize{8.4}{11}\selectfont\color{white}Gratuit \textbullet\ Google Play\\Tuteur IA, annales, concours, résultats\par}
      \end{minipage}
    \end{tcolorbox}
  \end{minipage}\hfill
  \begin{minipage}[t]{0.487\linewidth}
    \begin{tcolorbox}[enhanced,colback=poppurple,colframe=popink,boxrule=2.2pt,arc=10pt,
        drop shadow=popink,left=11pt,right=11pt,top=10pt,bottom=10pt,height=3.1cm,valign=center]
      \tikz[baseline=-0.7ex]\node[circle,fill=white,draw=popink,line width=1.3pt,minimum size=1.6cm,inner sep=0]{\popdisplay\fontsize{24}{24}\selectfont\color{poppurple}+};%
      \hspace{8pt}\begin{minipage}[c]{0.6\linewidth}
        {\popdisplay\fontsize{11}{12.5}\selectfont\color{white}\MakeUppercase{Passe à OnBuch+}}\par\vspace{4pt}
        {\raggedright\popsemi\fontsize{8.4}{11}\selectfont\color{white}Tous les cours signés, Tuteur IA illimité, fiches PDF — onglet OnBuch+ de l'app\par}
      \end{minipage}
    \end{tcolorbox}
  \end{minipage}
  \vspace{8pt}
  \popcontenu
  \vfill
  \signatureonbuch
  \restoregeometry
  \newpage
}

% Rangée « Ce cours contient » (page de garde)
\newcommand{\poptile}[3]{% #1 couleur #2 gros texte #3 légende
  \begin{tcolorbox}[enhanced,colback=#1,colframe=popink,boxrule=2pt,arc=9pt,shadow={2.5pt}{-2.5pt}{0pt}{popink},
      left=6pt,right=6pt,top=9pt,bottom=9pt,halign=center,valign=center,height=1.75cm,width=\linewidth,nobeforeafter]
    {\popdisplay\fontsize{12.5}{13}\selectfont\color{white}\MakeUppercase{#2}}\par\vspace{3pt}
    {\centering\popsemi\fontsize{7.8}{9.5}\selectfont\color{white}#3\par}
  \end{tcolorbox}}
\newcommand{\popcontenu}{%
  \par\noindent{\popxboldls\fontsize{7.5}{7.5}\selectfont\color{popmuted}CE COURS SIGNÉ CONTIENT}\par\vspace{7pt}
  \noindent\begin{minipage}[t]{0.235\linewidth}\poptile{popblue}{Cours}{complet, illustré, pas à pas}\end{minipage}\hfill
  \begin{minipage}[t]{0.235\linewidth}\poptile{poporange}{Méthodes}{et exercices résolus}\end{minipage}\hfill
  \begin{minipage}[t]{0.235\linewidth}\poptile{poppink}{Exercices}{type examen + corrigés}\end{minipage}\hfill
  \begin{minipage}[t]{0.235\linewidth}\poptile{popgreen}{Fiche bilan}{l'essentiel à retenir}\end{minipage}\par}

% Sommaire
\newtcolorbox{sommaire}{enhanced,colback=white,colframe=popink,boxrule=2.2pt,arc=10pt,
  drop shadow=popink,left=18pt,right=18pt,top=13pt,bottom=13pt,before skip=6pt,after skip=20pt,
  before upper={\poppill{Sommaire}{poppurple}{white}\par\vspace{9pt}\popsemi\fontsize{10.6}{14.5}\selectfont\color{popink}}}

% Signature de fin de cours — poussée tout en bas de la dernière page
\newcommand{\findecours}{%
  \par\vfill\nopagebreak
  \begin{center}\popsquiggle{poporange}\end{center}\vspace{4pt}
  \signatureonbuch
}
\AtBeginDocument{\def\llbracket{\mathopen{[\mkern-3.5mu[}}\def\rrbracket{\mathclose{]\mkern-3.5mu]}}} % intervalles d'entiers (glyphe ⟦ absent de la police)
% Glyphes absents de Fira Math : redéfinis à partir de glyphes disponibles
\AtBeginDocument{%
  \def\setminus{\mathbin{\backslash}}\let\smallsetminus\setminus
  \def\vdots{\mathord{\vbox{\baselineskip3.2pt\lineskiplimit0pt\kern4pt\hbox{.}\hbox{.}\hbox{.}}}}%
  \def\ddots{\mathinner{\mkern1mu\raise6.5pt\hbox{.}\mkern2mu\raise3.6pt\hbox{.}\mkern2mu\raise0.7pt\hbox{.}\mkern1mu}}%
  \def\triangle{\mathord{\scalebox{0.9}{$\symup{\Delta}$}}}%
  \def\nabla{\mathord{\raisebox{\depth}{\scalebox{1}[-1]{$\symup{\Delta}$}}}}%
  \def\checkmark{\popcheck{popdark}}%
  \def\star{\mathbin{\ast}}\def\bigstar{\mathord{\ast}}%
  \def\diamond{\mathbin{\scalebox{0.6}{$\Diamond$}}}%
  \def\bigcup{\mathop{\mathchoice{\raisebox{-0.3em}{\scalebox{1.7}{$\cup$}}}{\scalebox{1.25}{$\cup$}}{\cup}{\cup}}\displaylimits}%
  \def\bigcap{\mathop{\mathchoice{\raisebox{-0.3em}{\scalebox{1.7}{$\cap$}}}{\scalebox{1.25}{$\cap$}}{\cap}{\cap}}\displaylimits}%
  \def\lesseqqgtr{\mathrel{\lessgtr}}\def\bigoplus{\mathop{\oplus}\displaylimits}%
}
\providecommand{\ssi}{si et seulement si} % abréviation fréquente
\AtBeginDocument{\providecommand{\wideparen}{\overparen}} % arc de cercle

%% FIGFIX-BEGIN v3 — correctifs globaux des figures (docs/onbuch-generation/tools/figfix)
\usepackage{adjustbox}\usepackage{etoolbox}
% toute figure TikZ/pgfplots plus large que la ligne est réduite à la largeur disponible
\BeforeBeginEnvironment{tikzpicture}{\begin{adjustbox}{max width=\linewidth}}
\AfterEndEnvironment{tikzpicture}{\end{adjustbox}}
% légendes pgfplots lisibles par-dessus les courbes
\pgfplotsset{every axis/.append style={legend style={fill=white,fill opacity=0.92,text opacity=1,draw=black!15,inner sep=3pt}}}
% étiquettes d'arêtes (syntaxe edge["..."]) avec fond blanc
\tikzset{every edge quotes/.append style={fill=white,inner sep=1.2pt}}
% bulles de cartes mentales : texte à la ligne dans la bulle, bulle un peu plus grande
\tikzset{every concept/.append style={text width=2.0cm,align=center,minimum size=2.3cm,inner sep=2pt}}
%% FIGFIX-END
\begin{document}

\pagegarde{Les mécanismes fondamentaux de la reproduction sexuée chez les Mammifères}{SVTEEHB}{Tle TI}{Module I : Le Monde Vivant}{N03}{3 h}{Cette leçon étudie les mécanismes fondamentaux de la reproduction sexuée chez les Mammifères : structure et rôle des gonades, méiose, spermatogenèse et ovogenèse, puis fécondation. Elle montre, à partir de faits expérimentaux et de courbes d'ADN, comment la méiose et la fécondation assurent ensemble la constance du nombre de chromosomes et le brassage de l'information génétique.
\vspace{4pt}
\begin{multicols}{2}\small
\begin{itemize}
\item À la fin de cette leçon, je sais décrire la structure et les rôles des gonades mâle et femelle chez les Mammifères.
\item À la fin de cette leçon, je sais ordonner et décrire les différentes phases de la méiose.
\item À la fin de cette leçon, je sais expliquer la nécessité de la méiose dans la formation des gamètes.
\item À la fin de cette leçon, je sais comparer la spermatogenèse et l'ovogenèse (lieux, étapes, rendement, chronologie).
\item À la fin de cette leçon, je sais tracer et interpréter la courbe de variation de la quantité d'ADN au cours de la spermatogenèse.
\item À la fin de cette leçon, je sais identifier sur des documents les étapes de la fécondation.
\end{itemize}\end{multicols}}

\begin{sommaire}
\begin{multicols}{2}
\begin{enumerate}
\item Les gonades : structures et rôles
\item La méiose : une division en deux étapes
\item La gamétogenèse : spermatogenèse et ovogenèse
\item La fécondation
\item Bilan : méiose et fécondation, un couple indissociable
\item Activité d'intégration
\item Méthodes et exercices résolus
\item Exercices
\item Corrigés des exercices
\item Fiche bilan
\end{enumerate}
\end{multicols}
\end{sommaire}

\begin{objectifs}
\begin{itemize}
\item À la fin de cette leçon, je sais décrire la structure et les rôles des gonades mâle et femelle chez les Mammifères.
\item À la fin de cette leçon, je sais ordonner et décrire les différentes phases de la méiose.
\item À la fin de cette leçon, je sais expliquer la nécessité de la méiose dans la formation des gamètes.
\item À la fin de cette leçon, je sais comparer la spermatogenèse et l'ovogenèse (lieux, étapes, rendement, chronologie).
\item À la fin de cette leçon, je sais tracer et interpréter la courbe de variation de la quantité d'ADN au cours de la spermatogenèse.
\item À la fin de cette leçon, je sais identifier sur des documents les étapes de la fécondation.
\item À la fin de cette leçon, je sais expliquer la nécessité de la fécondation dans la reproduction sexuée.
\item À la fin de cette leçon, je sais expliquer comment méiose et fécondation assurent la constance du nombre de chromosomes et la diversité génétique.
\end{itemize}
\end{objectifs}

\begin{prerequis}
\begin{itemize}
\item Notion de chromosome, d'ADN et de caryotype (classe de Première)
\item La mitose et le cycle cellulaire (classe de Première)
\item Notion de cellules haploïdes et diploïdes (début d'année de Terminale)
\item Anatomie générale des appareils reproducteurs humains (classe de Seconde/Première)
\item Lecture et exploitation de courbes et de documents scientifiques
\end{itemize}
\end{prerequis}

\newpage

\coursec{Les gonades : structures et rôles}

\textbf{Situation d'accroche.} Au marché de Mokolo à Yaoundé, une vendeuse de poulets observe chaque semaine un fait banal en apparence : les poussins d'une même couvée, pourtant issus du même coq et de la même poule, ne se ressemblent jamais parfaitement. Quelques rues plus loin, à l'hôpital, un médecin explique à un couple que leur enfant a reçu exactement la moitié de son patrimoine génétique de chaque parent. Or l'espèce humaine possède $46$ chromosomes par cellule. Comment la cellule-œuf peut-elle en contenir $46$ si le spermatozoïde et l'ovule en apportaient chacun $46$ ? Il faudrait donc qu'un mécanisme \emph{divise par deux} le stock d'ADN avant la rencontre des gamètes, puis que la rencontre elle-même \emph{restaure} le nombre exact propre à l'espèce. Où ce mécanisme a-t-il lieu ? Dans quels organes ? C'est toute la problématique de cette leçon : \textbf{remonter des gonades, organes de fabrication des gamètes, jusqu'à la fécondation}, pour comprendre comment la reproduction sexuée transmet l'information génétique tout en brassant les caractères. Dans cette première partie, nous identifions les organes concernés, leur structure fine et leur double rôle.

\courssub{Gonades, gamètes, gamétogenèse : le vocabulaire de base}

\begin{definition}[Gonade, gamète, gamétogenèse]
Une \cle{gonade} est un organe reproducteur qui produit les \cle{gamètes}, cellules sexuelles spécialisées dans la reproduction. Chez les Mammifères, la gonade mâle est le \cle{testicule}, qui produit les \cle{spermatozoïdes} ; la gonade femelle est l'\cle{ovaire}, qui produit les \cle{ovules} (plus précisément des ovocytes II au moment de l'ovulation). La \cle{gamétogenèse} est l'ensemble des divisions et différenciations cellulaires qui conduisent à la formation des gamètes : on parle de \cle{spermatogenèse} chez le mâle et d'\cle{ovogenèse} chez la femelle. Les gamètes sont des cellules \cle{haploïdes} : ils ne contiennent qu'un seul exemplaire de chaque chromosome ($n = 23$ chromosomes chez l'Homme), alors que les cellules ordinaires de l'organisme sont \cle{diploïdes} ($2n = 46$ chromosomes).
\end{definition}

Retiens l'idée directrice : si les gamètes étaient diploïdes, chaque génération doublerait son nombre de chromosomes ($46 + 46 = 92$, puis $184$\ldots), ce qui est incompatible avec la stabilité de l'espèce. La gamétogenèse doit donc comporter une étape de \emph{réduction} du nombre de chromosomes : c'est la méiose, que nous étudierons dans la section suivante. Mais d'abord, où cela se passe-t-il ?

\courssub{Le testicule : structure et rôles}

\textbf{Observation.} Chez l'Homme, les deux testicules ne sont pas logés dans l'abdomen mais dans une pousse cutanée externe, le \cle{scrotum} (ou bourses). Chaque testicule, de forme ovoïde (environ \SI{4}{cm} de long, \SI{25}{g} chez l'adulte), est constitué d'un enchevêtrement de longs tubes repliés sur eux-mêmes : les \cle{tubes séminifères} (environ \SI{250}{m} de tubes au total pour les deux testicules). Ces tubes débouchent sur l'\cle{épididyme}, lieu de maturation et de stockage des spermatozoïdes.

\textbf{Analyse d'une coupe histologique.} À la loupe binoculaire puis au microscope, une coupe transversale de testicule montre de nombreuses sections circulaires de tubes séminifères. En observant un tube de la paroi vers le centre (la lumière), on constate un fait remarquable : les cellules sont d'autant plus \emph{matures} qu'elles sont proches du centre. À la périphérie, accolées à la membrane basale, se trouvent de petites cellules rondes, les \cle{spermatogonies} (cellules souches diploïdes) ; plus près du centre apparaissent des cellules de plus en plus grosses (spermatocytes), puis de petites cellules rondes (spermatides), et enfin, au bord de la lumière, des cellules flagellées : les spermatozoïdes, qui seront libérés dans le tube. Cette organisation en gradient de maturation est la signature spatiale de la spermatogenèse.

\begin{popfigure}
\begin{tikzpicture}[scale=1.0]
% Paroi du tube
\draw[fill=popcream, thick, popdark] (0,0) circle (4.2);
\draw[fill=white, thick, popdark] (0,0) circle (1.1);
% Membrane basale
\draw[popblue, very thick] (0,0) circle (4.0);
% Spermatogonies (périphérie)
\foreach \a in {0,30,...,330}{
  \draw[fill=popblueL, draw=popblue] (\a:3.7) circle (0.22);}
% Spermatocytes (intermédiaire)
\foreach \a in {15,60,105,150,195,240,285,330}{
  \draw[fill=popgreenL, draw=popgreen] (\a:2.9) circle (0.35);}
% Spermatides
\foreach \a in {0,40,80,120,160,200,240,280,320}{
  \draw[fill=poporangeL, draw=poporange] (\a:2.0) circle (0.18);}
% Spermatozoïdes (tête + flagelle vers la lumière)
\foreach \a in {20,65,110,155,200,245,290,335}{
  \draw[fill=poppink, draw=poppink] (\a:1.45) circle (0.13);
  \draw[poppink, thick] (\a:1.32) -- (\a:0.75);}
% Cellules de Sertoli (triangles allongés)
\foreach \a in {45,135,225,315}{
  \draw[fill=popgoldL, draw=popgold, thick] (\a:3.95) -- (\a+8:1.2) -- (\a-8:1.2) -- cycle;}
% Lumière
\node[popmuted] at (0,0) {\small lumière};
% Légendes fléchées
\draw[->, popblue] (3.7,0.6) -- (5.6,1.4) node[right, align=left, font=\small]{\textbf{1} Spermatogonies ($2n$)\\ (cellules souches)};
\draw[->, popgreen] (2.5,1.6) -- (5.0,2.9) node[right, align=left, font=\small]{\textbf{2} Spermatocytes};
\draw[->, poporange] (1.5,1.35) -- (4.6,4.1) node[right, align=left, font=\small]{\textbf{3} Spermatides};
\draw[->, poppink] (1.1,0.95) -- (3.4,5.0) node[right, align=left, font=\small]{\textbf{4} Spermatozoïdes\\ (flagelle vers la lumière)};
\draw[->, popgold] (-3.3,2.3) -- (-5.6,3.4) node[left, align=right, font=\small]{\textbf{5} Cellule de Sertoli\\ (cellule nourricière)};
\draw[->, popblue] (-4.0,-0.5) -- (-5.6,-1.2) node[left, align=right, font=\small]{\textbf{6} Membrane basale\\ (paroi du tube)};
\end{tikzpicture}
\legende{Coupe transversale schématique d'un tube séminifère. La maturation des cellules sexuelles mâles progresse de la paroi (membrane basale, 6) vers la lumière du tube : spermatogonies (1), spermatocytes (2), spermatides (3), spermatozoïdes (4). Les cellules de Sertoli (5), étirées de la paroi à la lumière, entourent et nourrissent les cellules en cours de maturation.}
\end{popfigure}

\textbf{Deux populations cellulaires aux rôles distincts.} Entre les tubes séminifères, dans le tissu interstitiel, se trouvent les \cle{cellules de Leydig} : elles sécrètent la \cle{testostérone}, hormone mâle responsable de l'apparition et du maintien des caractères sexuels secondaires (muqueuse des voies génitales, pilosité, mue de la voix, développement musculaire) et indispensable au déroulement normal de la spermatogenèse. À l'intérieur des tubes, les \cle{cellules de Sertoli}, grandes cellules étirées de la membrane basale jusqu'à la lumière, jouent un rôle \emph{nourricier} et protecteur : elles apportent aux cellules en maturation les nutriments nécessaires, phagocytent les déchets et constituent la \cle{barrière hémato-testiculaire} qui isole les futurs gamètes du sang.

\courssub{L'ovaire : structure et rôles}

Chez la femme, les deux ovaires, logés dans la cavité pelvienne, ont la taille d'une amande (environ \SI{3}{cm} de long). Une coupe d'ovaire montre, à la périphérie, des structures sphériques de tailles variées : les \cle{follicules ovariens}, chacun contenant une cellule sexuelle femelle en cours de maturation (l'ovocyte). On distingue des follicules primordiaux (petits, au repos), des follicules en croissance, puis un follicule mûr (follicule de De Graaf), dont la paroi finit par se rompre : c'est l'\cle{ovulation}, qui libère l'ovocyte. Après l'ovulation, le reste du follicule se transforme en une structure jaunâtre, le \cle{corps jaune}, qui sécrète des hormones puis régresse si aucune grossesse ne débute.

\begin{popfigure}
\begin{tikzpicture}[scale=1.0]
% Ovaire
\draw[fill=popcream, thick, popdark] (0,0) ellipse (5.2 and 3.6);
% Follicules primordiaux
\foreach \p in {(-4.3,1.8),(-4.5,0.6),(-4.2,-0.7),(-3.9,-1.9)}{
  \draw[fill=popblueL, draw=popblue] \p circle (0.28);
  \draw[fill=popblue] \p circle (0.09);}
% Follicule en croissance
\draw[fill=popgreenL, draw=popgreen] (-2.3,1.6) circle (0.75);
\draw[fill=popgreen] (-2.3,1.6) circle (0.18);
% Follicule mûr (De Graaf) à la surface
\draw[fill=poporangeL, draw=poporange, thick] (3.6,1.4) circle (1.35);
\draw[fill=white, draw=poporange] (3.6,1.4) circle (0.9);
\draw[fill=poporange] (3.0,2.0) circle (0.2);
% Corps jaune
\draw[fill=popgoldL, draw=popgold, thick, decorate, decoration={zigzag,segment length=4,amplitude=1.5}] (0.6,-1.5) circle (1.0);
% Ovulation : flèche de sortie
\draw[->, poppink, very thick] (4.9,1.6) -- (6.3,2.2) node[right, align=left, font=\small]{\textbf{4} Ovulation : libération\\ de l'ovocyte};
% Légendes
\draw[->, popblue] (-4.3,0.6) -- (-6.4,0.4) node[left, align=right, font=\small]{\textbf{1} Follicules\\ primordiaux};
\draw[->, popgreen] (-2.3,1.6) -- (-3.2,3.4) node[above, align=center, font=\small]{\textbf{2} Follicule\\ en croissance};
\draw[->, poporange] (3.6,2.75) -- (2.6,4.0) node[above, align=center, font=\small]{\textbf{3} Follicule mûr\\ (de De Graaf)};
\draw[->, popgold] (0.6,-2.5) -- (0.6,-4.3) node[below, align=center, font=\small]{\textbf{5} Corps jaune\\ (sécrète la progestérone)};
\end{tikzpicture}
\legende{Coupe schématique d'un ovaire de Mammifère montrant les différents stades folliculaires (en réalité non simultanés à cette échelle) : follicules primordiaux (1), follicule en croissance (2), follicule mûr (3) dont la rupture libère l'ovocyte lors de l'ovulation (4) ; le follicule rompu se transforme en corps jaune (5).}
\end{popfigure}

L'ovaire est aussi une glande endocrine : les follicules en croissance sécrètent les \cle{œstrogènes} (responsables des caractères sexuels secondaires féminins et de la première phase du cycle utérin), et le corps jaune sécrète la \cle{progestérone} (hormone de la seconde phase du cycle et de la grossesse).

\courssub{La preuve expérimentale du double rôle des gonades : castration et greffe}

Comment savoir que les gonades ne servent pas seulement à produire des gamètes, mais fabriquent aussi des substances agissant sur tout le corps ? C'est la démarche expérimentale classique d'\emph{ablation--greffe}.

\begin{experience}[Les expériences de Berthold sur le coq (1849)]
\textbf{Matériel et protocole.} Arnold Berthold utilise de jeunes coqs, dont le caractère sexuel secondaire le plus visible est la crête (avec la queue, les ergots et le chant). Il réalise trois lots :
\begin{enumerate}
\item \textbf{Lot témoin} : coqs intacts, non opérés.
\item \textbf{Lot castré} : ablation des deux testicules.
\item \textbf{Lot castré puis greffé} : ablation des testicules, puis greffe d'un testicule dans la cavité abdominale (sans reconnexion nerveuse).
\end{enumerate}
\textbf{Observations.} Le lot témoin développe une crête volumineuse et rouge, chante et se comporte en mâle. Le lot castré garde une crête atrophiée et pâle, ne chante pas : les caractères sexuels secondaires n'apparaissent pas. Le lot greffé, bien que le testicule ne soit relié ni aux nerfs ni aux voies génitales, développe normalement sa crête et ses comportements mâles.
\textbf{Interprétation.} Puisque la greffe agit sans connexion nerveuse, le testicule libère dans le \emph{sang} une substance chimique — une \cle{hormone}, la testostérone — responsable des caractères sexuels mâles. La castration supprime la source, la greffe la restaure.
\textbf{Conclusion.} Le testicule possède un double rôle : production de gamètes (fonction exocrine, dans les tubes séminifères) et sécrétion d'hormones (fonction endocrine, cellules de Leydig). Des expériences analogues d'ovariectomie et de greffe d'ovaire établissent le même double rôle chez la femelle.
\end{experience}

\begin{propriete}[Double fonction des gonades]
Les gonades sont des \cle{glandes mixtes} (amphocrines) :
\begin{itemize}
\item une \textbf{fonction exocrine} (cytogène) : production des gamètes haploïdes par gamétogenèse — spermatozoïdes dans les tubes séminifères, ovocytes dans les follicules ovariens ;
\item une \textbf{fonction endocrine} : sécrétion d'hormones sexuelles déversées dans le sang — testostérone (cellules de Leydig) chez le mâle ; œstrogènes (follicules) et progestérone (corps jaune) chez la femelle — qui contrôlent les caractères sexuels secondaires et les cycles reproducteurs.
\end{itemize}
Ce double rôle est démontré par les expériences de castration (suppression des deux fonctions) et de greffe (restauration de la fonction endocrine).
\end{propriete}

\begin{methode}[Interpréter une expérience de castration / greffe]
Face à un document de ce type au Baccalauréat, procède toujours ainsi :
\begin{enumerate}
\item \textbf{Identifie le facteur testé} : ici, la présence ou l'absence de la gonade (c'est la variable expérimentale).
\item \textbf{Repère le lot témoin} (animal intact) et le ou les lots expérimentaux (castré, castré-greffé) ; vérifie que tout le reste est identique (âge, alimentation, conditions d'élevage).
\item \textbf{Compare les résultats} : ce qui disparaît après ablation et réapparaît après greffe dépend de la gonade.
\item \textbf{Conclus en distinguant les deux fonctions} : si la greffe agit à distance sans connexion nerveuse ni canalaire, l'effet est \emph{sanguin}, donc hormonal (fonction endocrine) ; la production de gamètes (fonction exocrine) n'est restaurée que si le greffon contient le tissu germinal fonctionnel.
\end{enumerate}
\end{methode}

\begin{exoresolu}[Castration et injection d'extrait testiculaire]
On réalise trois lots de jeunes rats mâles. Lot A : rats intacts. Lot B : rats castrés. Lot C : rats castrés recevant chaque jour une injection d'extrait de testicule broyé. Après deux mois, on mesure la masse des vésicules séminales (glandes annexes dont le développement dépend des hormones mâles) : A : \SI{450}{mg} ; B : \SI{80}{mg} ; C : \SI{430}{mg}. Que peut-on conclure ?
\tcblower
\textbf{Solution détaillée.} Le facteur testé est la présence du testicule (lot B versus témoin A) et la présence de substances testiculaires (lot C). La castration (B) provoque une forte atrophie des vésicules séminales (\SI{80}{mg} contre \SI{450}{mg}) : leur développement dépend donc du testicule. L'injection d'extrait testiculaire (C) restaure presque totalement la masse (\SI{430}{mg}) alors que les rats n'ont plus de testicules : une \emph{substance chimique} fabriquée par le testicule et véhiculée par le sang (la testostérone) suffit à maintenir les caractères sexuels. \textbf{Conclusion :} le testicule exerce, en plus de sa fonction de production de gamètes, une fonction endocrine de sécrétion d'hormone mâle.
\end{exoresolu}

\begin{attention}[Gonade ou glande annexe : ne pas confondre !]
Erreur classique au Baccalauréat : attribuer la production des gamètes aux \emph{glandes annexes}. La \cle{prostate} et les \cle{vésicules séminales} ne produisent \emph{aucun gamète} : elles sécrètent le liquide séminal qui nourrit et transporte les spermatozoïdes. De même, l'\cle{épididyme} stocke et fait mûrir les spermatozoïdes mais ne les fabrique pas. Seules les gonades (testicules, ovaires) réalisent la gamétogenèse. Autre piège : dire que la greffe de Berthold prouve un contrôle \emph{nerveux} — c'est exactement l'inverse, puisque le greffon agit sans aucune connexion nerveuse.
\end{attention}

\begin{savaistu}[Pourquoi les testicules sont-ils hors de l'abdomen ?]
La spermatogenèse est thermosensible : elle exige une température d'environ \SI{35}{\celsius}, soit \SI{2}{\celsius} de moins que la température interne du corps (\SI{37}{\celsius}). C'est pourquoi les testicules sont logés dans le scrotum, à l'extérieur de la cavité abdominale, et qu'un muscle (le crémaster) les rapproche ou les éloigne du corps selon la température ambiante. Une exposition prolongée à la chaleur (fièvre prolongée, \cle{cryptorchidie} — testicule non descendu, fréquente chez le nouveau-né et traitée dans les hôpitaux camerounais par chirurgie) peut altérer durablement la fertilité.
\end{savaistu}

\begin{aretenir}[Savoir admis : un contrôle hormonal des gonades]
L'activité des gonades (gamétogenèse et sécrétions hormonales) n'est pas autonome : elle est contrôlée par deux hormones hypophysaires, la \cle{FSH} et la \cle{LH}, elles-mêmes sous dépendance de l'\cle{hypothalamus}. Les mécanismes fins de cette régulation seront étudiés dans la leçon consacrée aux cycles sexuels ; pour l'instant, retiens simplement qu'\emph{il existe} un contrôle hormonal des gonades.
\end{aretenir}

\textbf{Bilan de la section.} Les gonades — testicules et ovaires — sont des glandes mixtes : elles produisent les gamètes (fonction exocrine) et les hormones sexuelles (fonction endocrine), comme le prouvent les expériences de castration et de greffe. Leur structure interne (tubes séminifères, follicules ovariens) révèle des cellules sexuelles à tous les stades de maturation : les gonades sont bien le \emph{siège de la gamétogenèse}. Or cette gamétogenèse doit transformer des cellules diploïdes ($2n = 46$ chromosomes) en gamètes haploïdes ($n = 23$) : c'est précisément le rôle de la méiose, que nous étudions maintenant.

\coursec{La méiose : une division en deux étapes}

Dans la section précédente, nous avons vu que les gonades — testicules et ovaires — sont les organes où naissent les gamètes. Mais une question fondamentale reste posée : \emph{comment} une cellule diploïde de la gonade peut-elle produire des gamètes haploïdes ? La réponse tient en un mot : la \cle{méiose}. C'est ce mécanisme remarquable que nous allons maintenant étudier en détail.

\courssub{Le problème posé : comment éviter le doublement du nombre de chromosomes ?}

\begin{experience}[Une expérience de pensée pour comprendre la nécessité de la méiose]
Raisonnons à partir d'un fait connu : le caryotype humain compte $2n = 46$ chromosomes, organisés en 23 paires. Chaque individu naît de la fusion de deux gamètes : un spermatozoïde et un ovule.

\textbf{Hypothèse à tester :} supposons que les gamètes soient produits par \emph{mitose}, c'est-à-dire par une division qui conserve le nombre de chromosomes. Chaque gamète porterait alors 46 chromosomes.

\textbf{Conséquence :} à la fécondation, le zygote posséderait $46 + 46 = 92$ chromosomes ; à la génération suivante, 184 ; puis 368... Le nombre de chromosomes doublerait à chaque génération, ce qui est absurde : le caryotype humain est \emph{constant} d'une génération à l'autre.

\textbf{Conclusion :} il doit exister, quelque part dans la lignée germinale, une division \emph{réductionnelle} qui divise par deux le nombre de chromosomes : de $2n = 46$ à $n = 23$. Cette division, c'est la méiose.
\end{experience}

L'observation au microscope le confirme. Sur des coupes de testicules (chez le criquet, matériel historique classique, ou chez les Mammifères), on observe dans les tubes séminifères des figures de division très différentes de la mitose : les chromosomes s'y \emph{apparient deux à deux}, formant des structures doubles jamais vues en mitose. Cette observation a conduit les cytologistes de la fin du XIX\textsuperscript{e} siècle (van Beneden, 1883, chez l'ascaride) à décrire la méiose.

\begin{definition}[Méiose]
La \cle{méiose} est une succession de deux divisions cellulaires, précédées d'une \emph{unique} réplication de l'ADN, qui transforme une cellule diploïde ($2n$ chromosomes) en quatre cellules haploïdes ($n$ chromosomes). Elle ne se produit que dans les lignées germinales des gonades, au cours de la gamétogenèse.
\end{definition}

\begin{definition}[Chromosomes homologues]
Deux chromosomes sont dits \cle{homologues} lorsqu'ils ont la même taille, la même forme, la même position de centromère et portent les mêmes gènes aux mêmes emplacements (locus), mais éventuellement sous des versions différentes (allèles). Dans chaque paire, un chromosome vient du père (chromosome \emph{paternel}), l'autre de la mère (chromosome \emph{maternel}).
\end{definition}

\courssub{Vue d'ensemble : une réplication, deux divisions, quatre cellules}

Le principe général de la méiose tient en une phrase à retenir absolument :

\begin{aretenir}[Le schéma général de la méiose]
Une cellule-mère diploïde ($2n$ chromosomes) réplique son ADN \textbf{une seule fois} (phase S de l'interphase), puis subit \textbf{deux divisions successives} :
\begin{itemize}
\item la \cle{méiose I} ou \textbf{division réductionnelle} : elle sépare les chromosomes \emph{homologues} et fait passer la cellule de $2n$ à $n$ chromosomes ;
\item la \cle{méiose II} ou \textbf{division équationnelle} : elle sépare les \emph{chromatides sœurs}, comme une mitose ordinaire, sans changer la ploïdie.
\end{itemize}
Bilan : \textbf{1 cellule diploïde $\rightarrow$ 4 cellules haploïdes}, génétiquement différentes entre elles.
\end{aretenir}

\begin{propriete}[Pourquoi la méiose est-elle réductionnelle ?]
La réduction du nombre de chromosomes tient à un fait précis : la réplication de l'ADN (qui double la quantité d'ADN, par exemple de $2Q$ à $4Q$ dans une cellule diploïde, $Q$ étant la quantité d'ADN d'un gamète) est suivie de \emph{deux} séparations successives sans réplication intermédiaire : la séparation des homologues en anaphase I, puis la séparation des chromatides en anaphase II. Le rapport « une réplication / deux divisions » explique à lui seul le passage de $2n$ à $n$.
\end{propriete}

\courssub{La méiose I : la division réductionnelle}

\textbf{Prophase I : l'appariement des homologues et le crossing-over.} C'est la phase la plus longue et la plus riche en événements. Les chromosomes, déjà répliqués (chacun est formé de \emph{deux chromatides sœurs} jointes au centromère), se condensent. Puis se produit un phénomène absent de la mitose : chaque chromosome recherche son homologue et s'apparie avec lui, point par point, gène contre gène. Ce phénomène s'appelle la \emph{synapsis}.

\begin{definition}[Bivalent]
Un \cle{bivalent} (ou tétrade) est l'ensemble formé par deux chromosomes homologues appariés, chacun constitué de deux chromatides. Un bivalent comporte donc \textbf{quatre chromatides}. Une cellule humaine en prophase I contient 23 bivalents.
\end{definition}

Au sein des bivalents, un événement capital a lieu : les chromatides \emph{non sœurs} (l'une d'origine paternelle, l'autre d'origine maternelle) peuvent se croiser, se rompre au même point et échanger des segments réciproques.

\begin{definition}[Crossing-over]
Le \cle{crossing-over} (ou enjambement) est l'échange réciproque de segments de chromatides entre deux chromosomes homologues au cours de la prophase I. Les points de contact visibles au microscope sont appelés \emph{chiasmas}. Il en résulte des chromatides \emph{recombinées}, portant des allèles d'origine paternelle et maternelle sur un même chromosome.
\end{definition}

\begin{popfigure}
\begin{center}
\begin{tikzpicture}[scale=0.85]
% Bivalent avec crossing-over
% Chromosome paternel (bleu) : deux chromatides
\draw[line width=4pt, popblue] (0,0) -- (0,4);
\draw[line width=4pt, popblue] (0.35,0) -- (0.35,4);
% Chromosome maternel (orange) : deux chromatides
\draw[line width=4pt, poporange] (1.2,0) -- (1.2,4);
\draw[line width=4pt, poporange] (1.55,0) -- (1.55,4);
% Centromères
\fill[popdark] (0.175,2) circle (0.14);
\fill[popdark] (1.375,2) circle (0.14);
% Chiasma : croisement entre chromatides non soeurs
\draw[line width=2.5pt, popblue] (0.35,2.9) .. controls (0.7,2.6) .. (1.2,2.2);
\draw[line width=2.5pt, poporange] (1.2,2.9) .. controls (0.85,2.6) .. (0.35,2.2);
\fill[poppink] (0.775,2.55) circle (0.09);
\node[poppink, font=\small\bfseries] at (2.6,2.55) {chiasma};
\draw[->, poppink] (2.35,2.55) -- (0.95,2.55);
% Légendes
\node[popblue, font=\small] at (0.175,-0.5) {homologue paternel};
\node[poporange, font=\small] at (1.375,-0.9) {homologue maternel};
\node[font=\small] at (0.175,4.4) {2 chromatides sœurs};
% Flèche vers résultat
\draw[->, very thick, popdark] (3.4,2) -- (4.4,2);
% Après crossing-over
\begin{scope}[xshift=5cm]
\draw[line width=4pt, popblue] (0,2.6) -- (0,4);
\draw[line width=4pt, poporange] (0,0) -- (0,2.6);
\draw[line width=4pt, popblue] (0.35,0) -- (0.35,4);
\draw[line width=4pt, poporange] (1.2,2.6) -- (1.2,4);
\draw[line width=4pt, popblue] (1.2,0) -- (1.2,2.6);
\draw[line width=4pt, poporange] (1.55,0) -- (1.55,4);
\fill[popdark] (0.175,2) circle (0.14);
\fill[popdark] (1.375,2) circle (0.14);
\node[font=\small, align=center] at (0.775,-0.7) {chromatides\\ recombinées};
\end{scope}
\end{tikzpicture}
\end{center}
\legende{Schéma d'un bivalent en prophase I. À gauche : les deux homologues appariés (bleu : paternel ; orange : maternel), chacun formé de deux chromatides sœurs ; un chiasma se forme entre deux chromatides \emph{non sœurs}. À droite : après crossing-over, deux chromatides sont recombinées (bicolores), deux restent parentales.}
\end{popfigure}

\textbf{Métaphase I.} Les bivalents se placent sur la plaque équatoriale du fuseau, \emph{par paires} : les deux centromères d'un même bivalent sont accrochés à des pôles opposés. Point capital : l'orientation de chaque bivalent est \emph{aléatoire} — le chromosome paternel peut être tourné vers un pôle ou vers l'autre, indépendamment des autres paires. C'est la base du brassage interchromosomique.

\textbf{Anaphase I.} Les chromosomes \emph{homologues} se séparent et migrent vers les pôles opposés. Chaque chromosome migrant est encore formé de \emph{deux chromatides} : \textbf{le centromère ne se divise pas}. C'est ici, et seulement ici, que s'opère la réduction : chaque pôle ne reçoit qu'un chromosome de chaque paire, donc $n$ chromosomes.

\textbf{Télophase I.} Les chromosomes arrivent aux pôles, deux noyaux haploïdes se reforment (mais avec des chromosomes à deux chromatides), puis le cytoplasme se divise : on obtient \textbf{deux cellules haploïdes} ($n$ chromosomes à deux chromatides).

\begin{attention}[Le piège classique de l'anaphase I]
En \textbf{anaphase I}, ce sont les \emph{chromosomes homologues} qui se séparent, \textbf{pas} les chromatides : le centromère reste entier et chaque chromosome migrant compte encore deux chromatides. La division du centromère et la séparation des chromatides sœurs n'ont lieu qu'en \textbf{anaphase II}. Confondre les deux est l'erreur la plus fréquente au Baccalauréat — retiens : \emph{anaphase I = séparation des homologues ; anaphase II = séparation des chromatides}.
\end{attention}

\courssub{L'interphase brève et la méiose II : la division équationnelle}

Entre les deux divisions s'intercale une \textbf{interphase très brève} (parfois absente), appelée intercinèse. Son trait essentiel : il n'y a \textbf{aucune réplication de l'ADN}. La quantité d'ADN ne change pas.

La méiose II déroule ensuite, dans chacune des deux cellules, un programme comparable à une mitose :

\begin{itemize}
\item \textbf{Prophase II :} les chromosomes (toujours à deux chromatides) se recondensent ; le fuseau se forme. Il n'y a ni appariement, ni crossing-over : il n'y a plus d'homologues dans la cellule.
\item \textbf{Métaphase II :} les chromosomes s'alignent \emph{individuellement} sur la plaque équatoriale, comme en mitose.
\item \textbf{Anaphase II :} les centromères se divisent enfin ; les \emph{chromatides sœurs} se séparent et migrent vers les pôles opposés. Chaque chromatide devient un chromosome à une chromatide.
\item \textbf{Télophase II :} quatre noyaux haploïdes se reforment ; la division cytoplasmique aboutit à \textbf{quatre cellules haploïdes} à chromosomes simples.
\end{itemize}

\begin{popfigure}
\begin{center}
\begin{tikzpicture}[scale=0.62, every node/.style={font=\scriptsize}]
% ===== Cellule mère 2n=4 =====
\begin{scope}
\draw[popline] (0,0) ellipse (1.5 and 1.9);
\draw[line width=2.2pt, popblue] (-0.7,-0.7) -- (-0.7,0.7);
\draw[line width=2.2pt, popblue!55] (-0.35,-0.7) -- (-0.35,0.7);
\draw[line width=2.2pt, poporange] (0.35,-0.5) -- (0.35,0.5);
\draw[line width=2.2pt, poporange!55] (0.7,-0.5) -- (0.7,0.5);
\node[align=center] at (0,-2.5) {cellule-mère\\ $2n = 4$};
\end{scope}
\draw[->, very thick, popdark] (1.8,0) -- (2.8,0) node[midway, above]{réplication};
% ===== Prophase I : bivalents =====
\begin{scope}[xshift=4.9cm]
\draw[popline] (0,0) ellipse (1.5 and 1.9);
\draw[line width=2.2pt, popblue] (-0.85,-0.7) -- (-0.85,0.7);
\draw[line width=2.2pt, popblue] (-0.6,-0.7) -- (-0.6,0.7);
\draw[line width=2.2pt, popblue!55] (-0.35,-0.7) -- (-0.35,0.7);
\draw[line width=2.2pt, popblue!55] (-0.1,-0.7) -- (-0.1,0.7);
\draw[line width=2.2pt, poporange] (0.15,-0.5) -- (0.15,0.5);
\draw[line width=2.2pt, poporange] (0.4,-0.5) -- (0.4,0.5);
\draw[line width=2.2pt, poporange!55] (0.65,-0.5) -- (0.65,0.5);
\draw[line width=2.2pt, poporange!55] (0.9,-0.5) -- (0.9,0.5);
\node[align=center] at (0,-2.5) {prophase I\\ bivalents, crossing-over};
\end{scope}
\draw[->, very thick, popdark] (6.7,0) -- (7.7,0);
% ===== Anaphase I =====
\begin{scope}[xshift=9.8cm]
\draw[popline] (0,0) ellipse (1.5 and 1.9);
\draw[line width=2.2pt, popblue] (-0.85,0.9) -- (-0.85,1.5);
\draw[line width=2.2pt, popblue] (-0.6,0.9) -- (-0.6,1.5);
\draw[line width=2.2pt, poporange] (0.15,0.9) -- (0.15,1.5);
\draw[line width=2.2pt, poporange] (0.4,0.9) -- (0.4,1.5);
\draw[line width=2.2pt, popblue!55] (-0.85,-1.5) -- (-0.85,-0.9);
\draw[line width=2.2pt, popblue!55] (-0.6,-1.5) -- (-0.6,-0.9);
\draw[line width=2.2pt, poporange!55] (0.15,-1.5) -- (0.15,-0.9);
\draw[line width=2.2pt, poporange!55] (0.4,-1.5) -- (0.4,-0.9);
\node[align=center] at (0,-2.5) {anaphase I\\ séparation des homologues};
\end{scope}
\draw[->, very thick, popdark] (11.6,0) -- (12.6,0);
% ===== Deux cellules n=2 =====
\begin{scope}[xshift=14.6cm]
\draw[popline] (0,1) ellipse (1 and 0.8);
\draw[line width=2pt, popblue] (-0.5,0.8) -- (-0.5,1.4);
\draw[line width=2pt, popblue] (-0.25,0.8) -- (-0.25,1.4);
\draw[line width=2pt, poporange] (0.3,0.8) -- (0.3,1.4);
\draw[line width=2pt, poporange] (0.55,0.8) -- (0.55,1.4);
\draw[popline] (0,-1) ellipse (1 and 0.8);
\draw[line width=2pt, popblue!55] (-0.5,-1.2) -- (-0.5,-0.6);
\draw[line width=2pt, popblue!55] (-0.25,-1.2) -- (-0.25,-0.6);
\draw[line width=2pt, poporange!55] (0.3,-1.2) -- (0.3,-0.6);
\draw[line width=2pt, poporange!55] (0.55,-1.2) -- (0.55,-0.6);
\node[align=center] at (0,-2.5) {2 cellules\\ $n = 2$ (2 chromatides)};
\end{scope}
% ===== Anaphase II =====
\draw[->, very thick, popdark] (14.6,-2.9) -- (14.6,-3.7) node[midway, right, align=left]{méiose II};
\begin{scope}[xshift=9.8cm, yshift=-5.6cm]
\draw[popline] (0,0) ellipse (1.5 and 1.6);
\draw[line width=2pt, popblue] (-0.7,0.7) -- (-0.7,1.2);
\draw[line width=2pt, popblue] (-0.7,-1.2) -- (-0.7,-0.7);
\draw[line width=2pt, poporange] (0.5,0.7) -- (0.5,1.2);
\draw[line width=2pt, poporange] (0.5,-1.2) -- (0.5,-0.7);
\node[align=center] at (0,-2.2) {anaphase II\\ séparation des chromatides};
\end{scope}
\draw[->, very thick, popdark] (8.1,-5.6) -- (7.1,-5.6);
% ===== 4 cellules haploïdes =====
\begin{scope}[xshift=4.9cm, yshift=-5.6cm]
\draw[popline] (-0.75,0.75) circle (0.55);
\draw[popline] (0.75,0.75) circle (0.55);
\draw[popline] (-0.75,-0.75) circle (0.55);
\draw[popline] (0.75,-0.75) circle (0.55);
\draw[line width=1.6pt, popblue] (-0.9,0.55) -- (-0.9,0.95);
\draw[line width=1.6pt, poporange] (-0.55,0.55) -- (-0.55,0.95);
\draw[line width=1.6pt, popblue] (0.6,0.55) -- (0.6,0.95);
\draw[line width=1.6pt, poporange] (0.95,0.55) -- (0.95,0.95);
\draw[line width=1.6pt, popblue!55] (-0.9,-0.95) -- (-0.9,-0.55);
\draw[line width=1.6pt, poporange!55] (-0.55,-0.95) -- (-0.55,-0.55);
\draw[line width=1.6pt, popblue!55] (0.6,-0.95) -- (0.6,-0.55);
\draw[line width=1.6pt, poporange!55] (0.95,-0.95) -- (0.95,-0.55);
\node[align=center] at (0,-2.2) {4 cellules haploïdes\\ $n = 2$ (1 chromatide)};
\end{scope}
\end{tikzpicture}
\end{center}
\legende{Les deux divisions de la méiose avec $2n = 4$ (deux paires d'homologues ; couleurs vives : origine paternelle, couleurs claires : origine maternelle). La méiose I sépare les homologues (division réductionnelle) ; la méiose II sépare les chromatides sœurs (division équationnelle). Bilan : quatre cellules haploïdes à chromosomes simples.}
\end{popfigure}

\courssub{Bilan quantitatif : chromosomes et ADN}

Suivons, pour une cellule humaine, le nombre de chromosomes et la quantité d'ADN (en notant $Q$ la quantité d'ADN d'une cellule haploïde à chromosomes simples) :

\begin{center}
\begin{tabularx}{\linewidth}{|l|X|X|X|}
\hline
\rowcolor{popblueL} \textbf{Étape} & \textbf{Chromosomes} & \textbf{Chromatides/chromosome} & \textbf{Quantité d'ADN} \\
\hline
Cellule-mère (avant réplication) & $2n = 46$ & 1 & $2Q$ \\
\hline
Après réplication (début méiose I) & $2n = 46$ & 2 & $4Q$ \\
\hline
Fin de la méiose I (2 cellules) & $n = 23$ & 2 & $2Q$ par cellule \\
\hline
Fin de la méiose II (4 cellules) & $n = 23$ & 1 & $Q$ par cellule \\
\hline
\end{tabularx}
\end{center}

\begin{exemplebox}[Le cas de l'Homme]
Chez l'Homme, une spermatogonie à $2n = 46$ chromosomes réplique son ADN (passage de $2Q$ à $4Q$). La méiose I produit deux cellules à $n = 23$ chromosomes \emph{à deux chromatides} ($2Q$ chacune) ; la méiose II produit quatre spermatides à $n = 23$ chromosomes \emph{à une chromatide} ($Q$ chacune). La fécondation ($23 + 23$) restaure le nombre diploïde $2n = 46$.
\end{exemplebox}

\courssub{Conséquence génétique : la méiose, source de la diversité}

La méiose ne réduit pas seulement le nombre de chromosomes : elle \emph{brasse} les allèles par deux mécanismes complémentaires.

\begin{itemize}
\item \textbf{Brassage interchromosomique :} en anaphase I, la répartition des homologues paternels et maternels vers les pôles est \emph{indépendante et aléatoire} d'une paire à l'autre. Pour $n$ paires, il existe $2^{n}$ combinaisons possibles de gamètes.
\item \textbf{Brassage intrachromosomique :} les crossing-overs de la prophase I créent des chromatides recombinées, mêlant allèles paternels et maternels \emph{au sein d'un même chromosome}. Le nombre de combinaisons devient alors pratiquement infini.
\end{itemize}

\begin{savaistu}[Huit millions de gamètes possibles... avant même les crossing-overs]
Avec ses 23 paires de chromosomes, un être humain peut produire, par le seul brassage interchromosomique, $2^{23} \approx \num{8388608}$ types de gamètes différents, soit environ \textbf{8 millions} de combinaisons. Un couple peut donc, sans compter les crossing-overs, engendrer $2^{23} \times 2^{23} = 2^{46} \approx 70\,000$ milliards de zygotes génétiquement distincts. Voilà pourquoi, hors vrais jumeaux, deux enfants d'une même famille ne sont jamais génétiquement identiques — une réalité que chacun peut observer dans sa propre fratrie !
\end{savaistu}

\courssub{Méthode et comparaison avec la mitose}

\begin{methode}[Ordonner les phases d'une méiose sur des microphotographies]
Face à une série de clichés à classer (exercice classique du Bac), applique ces repères dans l'ordre :
\begin{enumerate}
\item \textbf{Des bivalents visibles ?} (chromosomes associés deux à deux, éventuellement des chiasmas) $\rightarrow$ c'est la \textbf{méiose I} : bivalents dispersés = prophase I ; bivalents alignés sur l'équateur = métaphase I.
\item \textbf{Des chromosomes à deux chromatides qui migrent vers les pôles} (centromères non fendus) $\rightarrow$ \textbf{anaphase I}.
\item \textbf{Des chromosomes à deux chromatides alignés seuls} (sans homologue) sur l'équateur $\rightarrow$ \textbf{métaphase II}.
\item \textbf{Des chromatides (chromosomes simples) qui migrent} vers les pôles $\rightarrow$ \textbf{anaphase II}.
\item \textbf{Quatre noyaux ou quatre cellules} en formation $\rightarrow$ \textbf{télophase II}.
\end{enumerate}
Le critère décisif reste : \emph{présence de bivalents = méiose I ; centromères fendus = anaphase II}.
\end{methode}

\begin{center}
\begin{tabularx}{\linewidth}{|l|X|X|}
\hline
\rowcolor{popblueL} \textbf{Critère} & \textbf{Mitose} & \textbf{Méiose} \\
\hline
Nombre de divisions & 1 & 2 (méiose I puis II) \\
\hline
Réplications de l'ADN & 1 & 1 \\
\hline
Appariement des homologues & Non & Oui (bivalents, prophase I) \\
\hline
Crossing-over & Non & Oui (prophase I) \\
\hline
Nombre de cellules filles & 2 & 4 \\
\hline
Ploïdie des cellules filles & $2n$ (diploïdes) & $n$ (haploïdes) \\
\hline
Identité génétique & Filles identiques à la mère & Filles toutes différentes \\
\hline
Rôle & Croissance, renouvellement cellulaire & Production des gamètes, diversité génétique \\
\hline
\end{tabularx}
\end{center}

\begin{exoresolu}[Raisonner sur les quantités d'ADN]
\textbf{Énoncé.} Chez le criquet, $2n = 24$ chromosomes. Une spermatogonie contient une quantité d'ADN égale à $2Q$. 1) Quelle est la quantité d'ADN de cette cellule au début de la prophase I ? 2) Combien de bivalents observe-t-on en métaphase I ? 3) Quelles sont le nombre de chromosomes et la quantité d'ADN de chaque cellule à la fin de la méiose I, puis à la fin de la méiose II ?
\tcblower
\textbf{Solution.}
\begin{enumerate}
\item Avant la prophase I, l'ADN a été répliqué une fois : la quantité double, donc $2 \times 2Q = 4Q$.
\item Les homologues s'apparient deux à deux : $24 / 2 = 12$ \textbf{bivalents}.
\item Fin de la méiose I : les homologues se sont séparés, chaque cellule reçoit $n = 12$ chromosomes \emph{à deux chromatides}, soit $2Q$ d'ADN. Fin de la méiose II : les chromatides se sont séparées, chacune des quatre cellules reçoit $n = 12$ chromosomes \emph{à une chromatide}, soit $Q$ d'ADN.
\end{enumerate}
\end{exoresolu}

\begin{aretenir}[L'essentiel de la section]
La méiose est une division \emph{réductionnelle} : une seule réplication de l'ADN suivie de deux divisions produit quatre cellules haploïdes à partir d'une cellule diploïde. La méiose I sépare les homologues (anaphase I, centromères intacts) ; la méiose II sépare les chromatides (anaphase II, centromères fendus). Grâce au brassage interchromosomique (répartition aléatoire des homologues) et intrachromosomique (crossing-over), la méiose est la grande source de la diversité génétique des gamètes — diversité que la fécondation, que nous étudierons bientôt, amplifiera encore en restaurant la diploïdie.
\end{aretenir}

\coursec{La gamétogenèse : spermatogenèse et ovogenèse}

Après avoir analysé le mécanisme universel de la méiose, nous comprenons maintenant \emph{comment} le nombre de chromosomes est divisé par deux. Reste à voir \emph{où}, \emph{quand} et \emph{combien de fois} ce processus s'insère dans la fabrication réelle des gamètes. C'est l'objet de la \cle{gamétogenèse}. Observons d'abord une coupe histologique de testicule et d'ovaire : les cellules germinales ne subissent pas la méiose immédiatement après leur formation. Elles traversent d'abord une phase de multiplication, puis une phase de croissance où elles accumulent des réserves, avant d'entrer en méiose (maturation), pour finir par se différencier sans division. Cette chronologie, commune aux deux sexes, masque toutefois des différences spectaculaires de rythme, de rendement et de devenir cellulaire.

\courssub{Définition générale et phases communes}

\begin{definition}[Gamétogenèse]
La \cle{gamétogenèse} est l'ensemble des processus cellulaires (mitoses, croissance, méiose, différenciation) qui transforment une cellule germinale diploïde (2n) en gamètes haploïdes (n) morphologiquement et fonctionnellement matures. Elle se déroule dans les gonades (testicules chez le mâle, ovaires chez la femelle) et comprend quatre phases successives :
\begin{enumerate}
    \item \textbf{Phase de multiplication} : mitoses successives des cellules germinales souches pour constituer un stock de cellules précurseurs.
    \item \textbf{Phase de croissance} : augmentation considérable du volume cellulaire, synthèse d'ARN et de protéines, constitution de réserves nutritives (vitellus, surtout chez la femelle) ; la \textbf{réplication de l'ADN} (phase S) a lieu en fin de croissance.
    \item \textbf{Phase de maturation} : c'est la \cle{méiose} (deux divisions successives sans réplication intermédiaire) qui ramène la formule chromosomique de 2n (4Q d'ADN) à n (Q d'ADN).
    \item \textbf{Phase de différenciation} : transformations morphologiques et fonctionnelles \textbf{sans division cellulaire} pour acquérir la structure spécialisée du gamète mature (flagelle, acrosome, enveloppes protectrices, etc.).
\end{enumerate}
\end{definition}

\begin{attention}[Piège terminologique]
Ne confondez pas \emph{gamétogenèse} (processus global) et \emph{méiose} (seule la phase de maturation). La méiose est une \emph{partie} de la gamétogenèse. Au Bac, si l'on vous demande « Décrivez la gamétogenèse », attendez-vous à décrire les 4 phases, pas seulement la méiose.
\end{attention}

\courssub{La spermatogenèse : une production continue et massive}

\begin{experience}[Observation d'une coupe de tube séminifère (document microscopique)]
\textbf{Matériel} : Coupe histologique de testicule de mammifère (rat, souris ou homme) colorée à l'hématoxyline-éosine. Grossissement $\times 400$.
\textbf{Observations} : Les tubes séminifères sont bordés d'un épithélium germinatif stratifié. Au contact de la membrane basale : petites cellules à noyau rond, peu coloré = \cle{spermatogonies} (2n). Vers la lumière : cellules plus grosses, à noyau volumineux et chromatine en filaments = \cle{spermatocytes I} (2n, 4Q). Plus rares (car leur durée de vie est brève) : \cle{spermatocytes II} (n, 2Q). Vers le centre : petites cellules allongées à noyau très condensé = \cle{spermatides} (n, Q), puis \cle{spermatozoïdes} (n, Q) dont les têtes sont enfoncées dans les cellules de Sertoli et les flagelles baignent dans la lumière du tube.
\textbf{Interprétation} : L'organisation spatiale reflète la chronologie : les divisions mitotiques se font en périphérie, puis la méiose et la différenciation poussent les cellules vers la lumière. Les cellules de Sertoli (noyau triangulaire clair) nourrissent, soutiennent et phagocytent les débris cytoplasmiques des cellules germinales.
\end{experience}

\begin{propriete}[Spermatogenèse : localisation, étapes et rendement]
\begin{itemize}
    \item \textbf{Localisation} : tubes séminifères des testicules (chez l'homme, plusieurs centaines de mètres de tubes au total).
    \item \textbf{Chronologie cellulaire} :
    \[
    \text{Spermatogonie (2n, 2Q)} \xrightarrow{\text{mitoses}} \text{Spermatocyte I (2n, 4Q)} \xrightarrow{\text{Méiose I}} \text{Spermatocyte II (n, 2Q)} \xrightarrow{\text{Méiose II}} \text{Spermatide (n, Q)} \xrightarrow{\text{Spermiogenèse}} \text{Spermatozoïde (n, Q)}
    \]
    \item \textbf{Rendement} : une spermatogonie entrée en maturation donne \textbf{4 spermatozoïdes fonctionnels} (divisions symétriques).
    \item \textbf{Kinétique} : continue à partir de la puberté (vers 13--14 ans chez l'homme) jusqu'à la mort. Durée totale $\approx 74$ jours chez l'homme. Production de l'ordre de $10^{8}$ spermatozoïdes par jour.
\end{itemize}
\end{propriete}

\begin{definition}[Spermiogenèse]
La \cle{spermiogenèse} est la phase de différenciation de la spermatogenèse au cours de laquelle une spermatide (cellule sphérique à noyau centré) se transforme en spermatozoïde \textbf{sans division cellulaire}. Elle implique : condensation extrême de la chromatine (remplacement des histones par des protamines), formation de l'\cle{acrosome} (vésicule issue de l'appareil de Golgi, contenant des enzymes hydrolytiques : hyaluronidase, acrosine), regroupement des mitochondries en hélice dans la pièce intermédiaire, élongation et formation du flagelle (axonème de structure 9 + 2), élimination du cytoplasme résiduel (phagocyté par les cellules de Sertoli).
\end{definition}

\begin{popfigure}
\begin{tikzpicture}[scale=0.9, every node/.style={font=\small}]
\def\headlen{2.5}
\def\midlen{1.5}
\def\taillen{6}
% Tête
\draw[popblue, fill=popblue!20, thick] (0,0.3) -- (\headlen,0.3) -- (\headlen,-0.3) -- (0,-0.3) -- cycle;
\draw[poppurple, fill=poppurple!30] (1.5,0) ellipse (0.8 and 0.22); % Noyau
\draw[poporange, fill=poporange!40] (0,0.3) -- (0.7,0.3) -- (0.7,-0.3) -- (0,-0.3) -- cycle; % Acrosome
% Pièce intermédiaire
\draw[popgreen, fill=popgreen!20, thick] (\headlen,0.25) -- (\headlen+\midlen,0.2) -- (\headlen+\midlen,-0.2) -- (\headlen,-0.25) -- cycle;
\foreach \x in {2.8,3.1,3.4,3.7} {
    \draw[popdark, fill=popdark!50] (\x,0) ellipse (0.1 and 0.15); % Mitochondries
}
% Flagelle
\draw[popdark, thick] (\headlen+\midlen,0) -- (\headlen+\midlen+\taillen,0);
% Annotations
\node[popblue, anchor=east] at (-0.3,0.6) {\textbf{Tête}};
\node[poppurple, anchor=east, align=right] at (-0.3,0.1) {Noyau (ADN haploïde,\\ chromatine hypercondensée)};
\node[poporange, anchor=east, align=right] at (-0.3,-0.5) {Acrosome (enzymes :\\ hyaluronidase, acrosine)};
\node[popgreen, anchor=south west, align=left] at (\headlen+\midlen+0.3,0.5) {\textbf{Pièce intermédiaire} : mitochondries (ATP pour la mobilité)};
\node[popdark, anchor=west, align=left] at (\headlen+\midlen+\taillen+0.3,0) {\textbf{Flagelle} : axonème 9+2, propulsion};
% Échelle
\draw[|-|, thick] (0,-1.2) -- (\headlen+\midlen+\taillen,-1.2) node[midway, below] {$\approx 60\,\mu\text{m}$ au total};
\end{tikzpicture}
\legende{Structure du spermatozoïde humain mature. La tête contient le noyau haploïde (n) coiffé de l'acrosome. La pièce intermédiaire est riche en mitochondries qui fournissent l'ATP nécessaire au battement du flagelle (structure axonémale 9+2), organite de la motilité.}
\end{popfigure}

\begin{savaistu}[Contexte camerounais : infertilité masculine et analyse de sperme]
Au Centre Hospitalier Universitaire de Yaoundé comme à l'Hôpital Laquintinie de Douala, le \emph{spermogramme} (analyse du sperme) est l'examen de première intention face à une infertilité de couple. Il évalue la concentration (norme OMS 2021 : $\ge 16 \times 10^{6}$ spermatozoïdes/mL), la mobilité (spermatozoïdes progressifs $\ge 30\,\%$), la morphologie (formes normales $\ge 4\,\%$) et la vitalité. Une azoospermie (absence de spermatozoïdes) peut témoigner d'un blocage de la spermatogenèse (varicocèle, exposition à des pesticides agricoles, chimiothérapie) ou d'une obstruction des canaux excréteurs. La compréhension fine des étapes de la spermatogenèse guide le diagnostic étiologique et la prise en charge (chirurgie, FIV/ICSI).
\end{savaistu}

\courssub{L'ovogenèse : une production discontinue, inégale et limitée dans le temps}

\begin{experience}[Observation d'une coupe d'ovaire (document microscopique)]
\textbf{Matériel} : Coupe d'ovaire de mammifère (lapine, rate, femme) colorée. Grossissements $\times 100$ et $\times 400$.
\textbf{Observations} : Contrairement au testicule, l'ovaire présente une zone corticale et une zone médullaire. Le cortex contient des \cle{follicules} à divers stades :
\begin{itemize}
    \item \textbf{Follicule primordial} : une cellule volumineuse (ovocyte I, $\approx 50\,\mu\text{m}$) à noyau clair (vésicule germinative), bloquée en \textbf{prophase I} (stade dictyotène), entourée d'une couche de cellules folliculaires aplaties.
    \item \textbf{Follicule primaire puis secondaire} : l'ovocyte I a grossi (jusqu'à $\approx 100\,\mu\text{m}$), les cellules folliculaires se sont multipliées (assise stratifiée), apparition de la \cle{zone pellucide} (matrice extracellulaire riche en glycoprotéines, dont ZP3).
    \item \textbf{Follicule tertiaire (antral, ou de De Graaf)} : cavité (antrum) remplie de liquide folliculaire riche en œstrogènes. L'ovocyte I est excentré dans le \emph{cumulus oophorus}. Juste avant l'ovulation, l'ovocyte I achève la méiose I $\rightarrow$ \cle{ovocyte II} (n, 2Q) + \cle{premier globule polaire} (n, 2Q, très peu de cytoplasme). L'ovocyte II commence la méiose II et \textbf{se bloque en métaphase II} ; c'est lui qui est expulsé à l'ovulation.
    \item \textbf{Corps jaune} : vestige du follicule rompu après l'ovulation, sécrétant la progestérone.
\end{itemize}
\textbf{Interprétation} : Le stock d'ovocytes I est constitué \textbf{avant la naissance} (vers la 20$^{\text{e}}$ semaine de vie fœtale : $\approx 6$ à $7 \times 10^{6}$ cellules germinales, puis apoptose massive $\rightarrow$ $\approx 1$ à $2 \times 10^{6}$ à la naissance $\rightarrow$ $\approx 3$ à $4 \times 10^{5}$ à la puberté). Il n'y a \textbf{aucun renouvellement post-natal}. À chaque cycle, une cohorte de follicules entre en croissance, mais un seul (en général) aboutit à l'ovulation. La méiose II ne s'achève \textbf{que si la fécondation a lieu}.
\end{experience}

\begin{propriete}[Ovogenèse : localisation, étapes, rendement et chronologie]
\begin{itemize}
    \item \textbf{Localisation} : cortex des ovaires.
    \item \textbf{Chronologie cellulaire} :
    \begin{align*}
    \text{Ovogonie (2n, 2Q)} &\xrightarrow{\text{mitoses (vie fœtale)}} \text{Ovocyte I (2n, 4Q), bloqué en prophase I (dictyotène)} \\
    &\xrightarrow{\text{reprise à chaque cycle après la puberté}} \text{Ovocyte II (n, 2Q)} + \text{Globule polaire I (n, 2Q)} \\
    &\xrightarrow{\text{blocage en métaphase II ; ovulation}} \xrightarrow{\text{si fécondation}} \text{Ovule (n, Q)} + \text{Globule polaire II (n, Q)}
    \end{align*}
    \item \textbf{Rendement} : une ovogonie $\rightarrow$ \textbf{1 seul ovule fonctionnel} + 2 à 3 globules polaires (qui dégénèrent par apoptose). Les divisions de la méiose sont \textbf{inégales} (cytokinèse asymétrique) : l'ovocyte conserve presque tout le cytoplasme, les organites, les ARN maternels et le vitellus ; les globules polaires ne reçoivent qu'un noyau haploïde.
    \item \textbf{Kinétique} : discontinue, cyclique (cycle ovarien $\approx 28$ jours), limitée dans le temps (ménopause vers 50 ans). Nombre total d'ovulations $\approx 400$ au cours de la vie.
\end{itemize}
\end{propriete}

\begin{definition}[Globule polaire]
Un \cle{globule polaire} est une petite cellule haploïde (n) issue des divisions inégales de la méiose lors de l'ovogenèse. Il contient un noyau haploïde mais très peu de cytoplasme et d'organites ; il est voué à la dégénérescence (apoptose). Le premier globule polaire (GP I) se forme à la fin de la méiose I (juste avant l'ovulation) ; le second (GP II) se forme à la fin de la méiose II (seulement après l'entrée du spermatozoïde). Leur rôle est d'éliminer l'excédent de matériel génétique tout en préservant les réserves cytoplasmiques de l'ovule.
\end{definition}

\begin{attention}[Erreur classique au Bac : « L'ovogenèse produit 4 ovules »]
\textbf{FAUX.} La méiose produit bien 4 noyaux haploïdes, mais la cytokinèse est inégale. Une seule cellule reçoit le cytoplasme et devient l'ovule fonctionnel ; les autres noyaux sont exclus dans les globules polaires (GP I, GP II, et parfois division du GP I). \textbf{Rédigez toujours :} « L'ovogenèse aboutit à la formation d'un ovule et de 2 à 3 globules polaires qui dégénèrent. » Justifiez l'inégalité : « Elle permet de concentrer les réserves nutritives (vitellus), les mitochondries, les ARN maternels et le matériel cytoplasmique dans une seule cellule, pour assurer les premiers stades du développement embryonnaire avant l'implantation. »
\end{attention}

\begin{aretenir}[Comparaison spermatogenèse / ovogenèse]
\begin{center}
\begin{tabularx}{\linewidth}{|l|X|X|}\hline
\rowcolor{popblueL}
\textbf{Critère} & \textbf{Spermatogenèse (mâle)} & \textbf{Ovogenèse (femelle)} \\ \hline
Localisation & Tubes séminifères (testicules) & Follicules ovariens (ovaires) \\ \hline
Début & Puberté ($\approx$ 13--14 ans) & Vie fœtale (20$^{\text{e}}$ SA) : constitution du stock \\ \hline
Rythme & Continu, toute la vie & Discontinu, cyclique (1 ovulation par cycle), arrêt à la ménopause \\ \hline
Durée totale & $\approx$ 74 jours (homme) & Des années (blocage en prophase I) + $\approx$ 3 mois de croissance folliculaire \\ \hline
Divisions & Symétriques (égales) & Asymétriques (inégales) \\ \hline
Rendement / cellule entrée en méiose & 4 spermatozoïdes fonctionnels & 1 ovule fonctionnel + globules polaires \\ \hline
Stock de précurseurs & Renouvelé en permanence (cellules souches spermatogoniales) & Fixé avant la naissance, non renouvelé, s'épuise \\ \hline
Blocages méiotiques & Aucun (enchaînement continu) & 2 blocages : prophase I (long, années) et métaphase II (court, heures) \\ \hline
Différenciation & Spermiogenèse (flagelle, acrosome) & Achèvement de la méiose II + modifications corticales à la fécondation \\ \hline
\end{tabularx}
\end{center}
\end{aretenir}

\courssub{Variation de la quantité d'ADN au cours de la spermatogenèse : méthode et tracé}

C'est un savoir-faire exigé au Bac : « Tracer et interpréter la courbe de variation de la quantité d'ADN en fonction du temps au cours de la spermatogenèse ». La quantité d'ADN est exprimée en \textbf{Q} (quantité d'ADN contenue dans un gamète haploïde mature). Une cellule diploïde en G1 possède 2Q ; après réplication (G2), elle possède 4Q.

\begin{methode}[Tracer et interpréter la courbe Q = f(t) pour la spermatogenèse]
\begin{enumerate}
    \item \textbf{Repère de base} : définir Q = quantité d'ADN d'un spermatozoïde (n). La spermatogonie en G1 est 2n $\rightarrow$ 2Q.
    \item \textbf{Phase de multiplication (mitoses)} : palier à \textbf{2Q}. Les mitoses conservent la quantité d'ADN par cellule (2Q $\rightarrow$ 2Q après chaque division, la réplication étant compensée par la division).
    \item \textbf{Phase de croissance (phase S)} : \textbf{montée 2Q $\rightarrow$ 4Q} (réplication de l'ADN). La cellule devient spermatocyte I (2n, 4Q). Palier à 4Q (prophase I, métaphase I).
    \item \textbf{Méiose I (division réductionnelle)} : \textbf{chute 4Q $\rightarrow$ 2Q} (séparation des chromosomes homologues). Formation de 2 spermatocytes II (n, 2Q). Palier bref à 2Q.
    \item \textbf{Méiose II (division équationnelle)} : \textbf{chute 2Q $\rightarrow$ Q} (séparation des chromatides sœurs). Formation de 4 spermatides (n, Q). Palier à Q.
    \item \textbf{Spermiogenèse} : palier à \textbf{Q} (ni division, ni réplication).
    \item \textbf{Légende} : axe des abscisses = temps (ou étapes successives) ; axe des ordonnées = quantité d'ADN (Q). Indiquer les paliers (2Q, 4Q, 2Q, Q) et les types cellulaires correspondants.
\end{enumerate}
\end{methode}

\begin{popfigure}
\begin{tikzpicture}
\begin{axis}[
    width=\linewidth,
    height=8cm,
    axis lines=middle,
    xlabel={Temps / Étapes successives},
    ylabel={Quantité d'ADN (Q)},
    xtick={1,2,3,4,5,6,7},
    xticklabels={Spermatogonie (G1), Phase S, Spermatocyte I, Méiose I, Spermatocyte II, Méiose II, Spermatide / Spermatozoïde},
    ytick={1,2,4},
    yticklabels={Q, 2Q, 4Q},
    ymin=0, ymax=4.5,
    xmin=0.5, xmax=7.5,
    tick label style={font=\footnotesize, align=center},
    x tick label style={rotate=30, anchor=east},
    grid=major,
    grid style={popline},
]
% Palier 2Q (Spermatogonie)
\addplot[popcurve, thick, domain=1:1.5, samples=2] {2};
% Montée Phase S (réplication)
\addplot[popcurve, thick, domain=1.5:2, samples=50] ({x}, {2 + 2*(x-1.5)/0.5});
% Palier 4Q (Spermatocyte I)
\addplot[popcurve, thick, domain=2:3, samples=2] {4};
% Chute Méiose I
\addplot[popcurve, thick, domain=3:3.5, samples=50] ({x}, {4 - 2*(x-3)/0.5});
% Palier 2Q (Spermatocyte II)
\addplot[popcurve, thick, domain=3.5:4.5, samples=2] {2};
% Chute Méiose II
\addplot[popcurve, thick, domain=4.5:5, samples=50] ({x}, {2 - 1*(x-4.5)/0.5});
% Palier Q (Spermatide/Spermatozoïde)
\addplot[popcurve, thick, domain=5:7, samples=2] {1};
% Annotations
\node[popblue, anchor=south, font=\small] at (axis cs:1.25,2.1) {2Q};
\node[poporange, anchor=south, font=\small] at (axis cs:2.5,4.1) {4Q};
\node[popgreen, anchor=south, font=\small] at (axis cs:4,2.1) {2Q};
\node[poppurple, anchor=south, font=\small] at (axis cs:6,1.1) {Q};
\node[popdark, anchor=north west, font=\small, align=left] at (axis cs:0.7,4.4) {
    \textbf{Repères :}\\
    $\nearrow$ = Réplication (phase S)\\
    $\searrow$ = Méiose I (réduction)\\
    $\searrow$ = Méiose II (équation)
};
\end{axis}
\end{tikzpicture}
\legende{Variation de la quantité d'ADN (en Q) au cours de la spermatogenèse. La réplication (phase S) fait passer la quantité d'ADN de 2Q à 4Q. La méiose I la ramène à 2Q (séparation des chromosomes homologues). La méiose II la ramène à Q (séparation des chromatides sœurs). La spermiogenèse se déroule à quantité d'ADN constante (Q).}
\end{popfigure}

\begin{exoresolu}[Analyse d'un histogramme d'ADN et identification des populations cellulaires]
\textbf{Énoncé :} On présente à un élève l'histogramme ci-dessous, obtenu par cytométrie de flux sur une population de cellules testiculaires marquées par un fluorochrome quantifiant l'ADN. La courbe montre trois pics principaux à Q, 2Q et 4Q, et une zone étalée entre 2Q et 4Q.
\begin{center}
\begin{tikzpicture}
\begin{axis}[
    width=0.7\linewidth, height=5cm,
    axis lines=middle, xlabel={Quantité d'ADN (Q)}, ylabel={Nombre de cellules},
    xmin=0, xmax=5, ymin=0, ymax=100,
    xtick={1,2,3,4}, xticklabels={Q, 2Q, 3Q, 4Q},
    ytick=\empty,
]
\addplot[popblue, thick, domain=0.5:1.5, samples=100] {80*exp(-(x-1)^2/0.05)};
\addplot[poporange, thick, domain=1.5:2.5, samples=100] {60*exp(-(x-2)^2/0.05)};
\addplot[poppurple, thick, domain=2.2:3.8, samples=100] {10*exp(-(x-3)^2/0.4)};
\addplot[popgreen, thick, domain=3.5:4.5, samples=100] {40*exp(-(x-4)^2/0.05)};
\end{axis}
\end{tikzpicture}
\end{center}
\textbf{Questions :}
\begin{enumerate}
    \item Identifier les types cellulaires correspondant aux pics à Q, 2Q et 4Q.
    \item Interpréter la zone étalée entre 2Q et 4Q.
\end{enumerate}
\tcblower
\textbf{Solution détaillée :}
\begin{enumerate}
    \item \textbf{Pic à 2Q} : spermatogonies (en G1) \textbf{et} spermatocytes II (après méiose I). Ces deux populations sont confondues car elles ont la même quantité d'ADN (2Q) mais des formules chromosomiques différentes (2n pour les spermatogonies, n pour les spermatocytes II).
    \textbf{Pic à 4Q} : spermatocytes I (après réplication, en prophase I / métaphase I). Formule 2n, 4Q.
    \textbf{Pic à Q} : spermatides et spermatozoïdes (après méiose II). Formule n, Q.
    \item \textbf{Zone étalée entre 2Q et 4Q} : elle correspond aux cellules en \textbf{phase S (réplication de l'ADN)}. La quantité d'ADN augmente continûment de 2Q vers 4Q pendant la phase S ; les cellules prélevées au hasard se trouvent à des stades intermédiaires de réplication, créant une distribution large entre les deux pics. C'est la preuve que la réplication est un processus continu, et non un saut instantané.
\end{enumerate}
\textbf{Point méthode :} au Bac, face à un histogramme de cytométrie de flux, identifiez les pics nets (G1, G2/M, gamètes) et la zone étalée (phase S). Ne dites jamais « méiose » pour les valeurs intermédiaires entre 2Q et 4Q.
\end{exoresolu}

\courssub{Exemple chiffré et perspective évolutive}

\begin{exemplebox}[Comparaison quantitative : mâle / femelle chez l'Homme]
\begin{itemize}
    \item \textbf{Homme (20--50 ans)} : $\approx 10^{8}$ spermatozoïdes produits par jour, soit de l'ordre de $\mathbf{10^{12}}$ spermatozoïdes au cours de la vie. Seuls quelques centaines atteignent le site de la fécondation ; un seul féconde l'ovocyte II.
    \item \textbf{Femme (vie fœtale $\rightarrow$ ménopause)} : stock initial $\approx 6$ à $7 \times 10^{6}$ cellules germinales (20$^{\text{e}}$ SA) $\rightarrow$ $\approx 1$ à $2 \times 10^{6}$ à la naissance $\rightarrow$ $\approx 3$ à $4 \times 10^{5}$ à la puberté $\rightarrow$ $\approx \mathbf{400}$ ovulations effectives.
    \item \textbf{Interprétation évolutive} : la stratégie « mâle » = grand nombre de gamètes, petite taille, mobilité, production continue. La stratégie « femelle » = petit nombre de gamètes, grande taille, immobilité, réserves cytoplasmiques massives, protection (zone pellucide). Les deux stratégies convergent vers le même but : la formation d'un zygote diploïde viable.
\end{itemize}
\end{exemplebox}

\begin{savaistu}[Âge maternel et risque de trisomie : le prix du blocage prolongé]
Les ovocytes I d'une femme de 40 ans sont restés bloqués en \textbf{prophase I (stade dictyotène)} depuis sa propre vie fœtale, soit plus de \textbf{40 ans}. Durant cette attente, les complexes de cohésine qui maintiennent ensemble les chromatides sœurs et les chromosomes homologues subissent une usure progressive. Lors de la reprise méiotique (ovulation), cette cohésion affaiblie augmente le risque de \textbf{non-disjonction} (mauvaise séparation des chromosomes) à l'anaphase I ou II. C'est la cause majeure de l'augmentation de l'incidence de la \textbf{trisomie 21} avec l'âge maternel : $\approx 1/1500$ naissances à 20 ans, $\approx 1/350$ à 35 ans, $\approx 1/30$ à 45 ans. Chez l'homme, la spermatogenèse continue renouvelle constamment les cohortes de spermatocytes I ; l'âge paternel a un effet beaucoup plus faible sur l'aneuploïdie (il est surtout associé à des mutations ponctuelles \emph{de novo}).
\end{savaistu}

\begin{popfigure}
\begin{tikzpicture}[scale=0.85, every node/.style={font=\small}]
\tikzset{
    cell/.style={draw=popdark, thick, rounded corners=3pt, minimum height=1.2cm, minimum width=2.2cm, align=center, fill=popcream},
    arrow/.style={-Stealth, thick, popdark},
    lab/.style={font=\tiny, text=popdark},
    phase/.style={font=\small\bfseries, text=popblue},
}
% --- SPERMATOGENÈSE (Gauche) ---
\node[phase] (sp_title) at (-5, 6.5) {SPERMATOGENÈSE (testicule)};
\node[cell, fill=popblue!10, align=center] (sg) at (-5, 5){Spermatogonies\\ (2n, 2Q)\\ \textit{mitoses de multiplication}};
\node[cell, fill=poporange!10, align=center] (sc1) at (-5, 3){Spermatocyte I\\ (2n, 4Q)\\ \textit{croissance + réplication}};
\node[cell, fill=popgreen!10, align=center] (sc2a) at (-7, 1){Spermatocyte II\\ (n, 2Q)};
\node[cell, fill=popgreen!10, align=center] (sc2b) at (-3, 1){Spermatocyte II\\ (n, 2Q)};
\node[cell, fill=poppurple!10, align=center] (st1) at (-8, -1){Spermatide\\ (n, Q)};
\node[cell, fill=poppurple!10, align=center] (st2) at (-6, -1){Spermatide\\ (n, Q)};
\node[cell, fill=poppurple!10, align=center] (st3) at (-4, -1){Spermatide\\ (n, Q)};
\node[cell, fill=poppurple!10, align=center] (st4) at (-2, -1){Spermatide\\ (n, Q)};
\node[cell, fill=poppink!10, align=center] (sz1) at (-8, -3){Spermatozoïde\\ (n, Q)};
\node[cell, fill=poppink!10, align=center] (sz2) at (-6, -3){Spermatozoïde\\ (n, Q)};
\node[cell, fill=poppink!10, align=center] (sz3) at (-4, -3){Spermatozoïde\\ (n, Q)};
\node[cell, fill=poppink!10, align=center] (sz4) at (-2, -3){Spermatozoïde\\ (n, Q)};

% Flèches Sperma
\draw[arrow] (sg) -- (sc1) node[midway, right, lab] {phase S};
\draw[arrow] (sc1) -- (sc2a) node[midway, left, lab] {Méiose I};
\draw[arrow] (sc1) -- (sc2b);
\draw[arrow] (sc2a) -- (st1) node[midway, left, lab] {Méiose II};
\draw[arrow] (sc2a) -- (st2);
\draw[arrow] (sc2b) -- (st3);
\draw[arrow] (sc2b) -- (st4);
\draw[arrow] (st1) -- (sz1) node[midway, left, lab] {Spermiogenèse};
\draw[arrow] (st2) -- (sz2);
\draw[arrow] (st3) -- (sz3);
\draw[arrow] (st4) -- (sz4);

% Légendes phases Sperma
\node[phase, rotate=90, anchor=east] at (-9.5, 4) {Multiplication};
\node[phase, rotate=90, anchor=east] at (-9.5, 2.5) {Croissance};
\node[phase, rotate=90, anchor=east] at (-9.5, 0) {Maturation (méiose)};
\node[phase, rotate=90, anchor=east] at (-9.5, -2) {Différenciation};

% --- OVOGENÈSE (Droite) ---
\node[phase] (ov_title) at (5, 6.5) {OVOGENÈSE (ovaire)};
\node[cell, fill=popblue!10, align=center] (og) at (5, 5){Ovogonie\\ (2n, 2Q)\\ \textit{mitoses (vie fœtale)}};
\node[cell, fill=poporange!10, align=center] (oc1) at (5, 3){Ovocyte I\\ (2n, 4Q)\\ \textbf{blocage en prophase I}\\ \textit{(dictyotène, années)}};
\node[cell, fill=popgreen!10, align=center] (oc2) at (3, 1){Ovocyte II\\ (n, 2Q)};
\node[cell, fill=popgreen!10, minimum width=1.5cm, align=center] (gp1) at (7, 1){Globule\\ polaire I\\ (n, 2Q)\\ \textit{(dégénère)}};
\node[cell, fill=poppurple!10, align=center] (ov) at (3, -1){Ovule\\ (n, Q)\\ \textit{(si fécondation)}};
\node[cell, fill=poppurple!10, minimum width=1.5cm, align=center] (gp2) at (7, -1){Globule\\ polaire II\\ (n, Q)\\ \textit{(dégénère)}};

% Flèches Ovo
\draw[arrow] (og) -- (oc1) node[midway, right, lab] {phase S};
\draw[arrow] (oc1) -- (oc2) node[midway, left, lab] {Méiose I (cycle ovarien)};
\draw[arrow] (oc1) -- (gp1) node[midway, right, lab] {division inégale};
\draw[arrow] (oc2) -- (ov) node[midway, left, lab] {Méiose II (si fécondation)};
\draw[arrow] (oc2) -- (gp2) node[midway, right, lab] {division inégale};

% Légendes phases Ovo
\node[phase, rotate=90, anchor=west, align=center] at (9.5, 4) {Multiplication\\ (vie fœtale seulement)};
\node[phase, rotate=90, anchor=west, align=center] at (9.5, 2.5) {Croissance + blocage long};
\node[phase, rotate=90, anchor=west, align=center] at (9.5, 0) {Maturation I\\ (par cycle)};
\node[phase, rotate=90, anchor=west, align=center] at (9.5, -1) {Maturation II\\ (fécondation)};

% Accolades temporelles
\draw[decorate, decoration={brace, amplitude=5pt}, thick, popdark] (-9.8, 5.6) -- (-9.8, -3.6) node[midway, left=8pt, rotate=90, font=\small\bfseries] {Continue dès la puberté};
\draw[decorate, decoration={brace, amplitude=5pt}, thick, popdark] (9.9, 5.6) -- (9.9, 3.6) node[midway, right=8pt, rotate=90, font=\small\bfseries] {Vie fœtale};
\draw[decorate, decoration={brace, amplitude=5pt}, thick, popdark] (9.9, 2.4) -- (9.9, -1.6) node[midway, right=8pt, rotate=90, font=\small\bfseries] {Puberté $\rightarrow$ ménopause (cyclique)};
\end{tikzpicture}
\legende{Schéma comparatif des gamétogenèses. À gauche, la spermatogenèse : production continue, divisions symétriques, 4 spermatozoïdes par spermatocyte I. À droite, l'ovogenèse : stock constitué \emph{in utero}, blocage prolongé en prophase I, divisions asymétriques (globules polaires), 1 ovule par ovocyte I. Les couleurs correspondent aux phases : \textcolor{popblue}{multiplication}, \textcolor{poporange}{croissance/réplication}, \textcolor{popgreen}{méiose I}, \textcolor{poppurple}{méiose II}, \textcolor{poppink}{différenciation}.}
\end{popfigure}

\begin{exercice}[Entraînement Bac : analyse de documents sur la gamétogenèse]
\textbf{Document 1 :} courbe de variation de la quantité d'ADN au cours de l'ovogenèse (schématique).
\textbf{Document 2 :} coupe histologique d'un follicule de De Graaf (ovocyte II + globule polaire I visibles).
\textbf{Document 3 :} tableau comparant le nombre de mitoses pré-méiotiques chez la souris (mâle : mitoses continues toute la vie ; femelle : aucune mitose germinale après la vie fœtale).

\textbf{Questions :}
\begin{enumerate}
    \item Sur le Document 1, identifier les étapes correspondant aux variations de la quantité d'ADN (2Q $\rightarrow$ 4Q $\rightarrow$ 2Q $\rightarrow$ Q) et les associer aux noms des cellules (ovocyte I, ovocyte II, ovule). Préciser à quel moment survient l'ovulation.
    \item À partir du Document 2, justifier que l'ovocyte II est bloqué en métaphase II.
    \item Expliquer, à l'aide du Document 3 et de vos connaissances, pourquoi le risque d'anomalies chromosomiques (aneuploïdies) augmente avec l'âge chez la femelle mais beaucoup moins chez le mâle.
    \item Calculer le nombre théorique de spermatozoïdes produits à partir d'une spermatogonie dont la descendance a subi 10 mitoses pré-méiotiques suivies de la méiose.
\end{enumerate}
\end{exercice}

\begin{corrige}[Exercice Entraînement Bac]
\textbf{1.} Courbe de l'ovogenèse :
\begin{itemize}
    \item Palier 2Q : ovogonies (mitoses fœtales) puis ovocyte I en début de croissance.
    \item Montée 2Q $\rightarrow$ 4Q : phase S (réplication) dans l'ovocyte I en croissance.
    \item Palier 4Q (très long) : ovocyte I bloqué en prophase I (dictyotène).
    \item Chute 4Q $\rightarrow$ 2Q : méiose I (réductionnelle) $\rightarrow$ ovocyte II (n, 2Q) + globule polaire I. \textbf{L'ovulation a lieu à ce stade} (libération de l'ovocyte II).
    \item Palier 2Q (court) : ovocyte II bloqué en métaphase II (dans la trompe utérine).
    \item Chute 2Q $\rightarrow$ Q : méiose II (équationnelle), \textbf{seulement si fécondation} $\rightarrow$ ovule (n, Q) + globule polaire II.
    \item Palier Q : ovule ; la fusion avec le noyau du spermatozoïde (n, Q) restaurera la diploïdie (2n, 2Q) dans le zygote.
\end{itemize}
\textbf{2.} Document 2 : l'ovocyte II présente un fuseau de méiose II avec les chromosomes (à deux chromatides) alignés sur la plaque équatoriale, et un globule polaire I adjacent. C'est la signature cytologique du blocage en métaphase II.
\textbf{3.} Chez le mâle : les mitoses pré-méiotiques se succèdent toute la vie (cellules souches spermatogoniales). Les spermatocytes I entrent en méiose quelques semaines après leur formation ; le temps passé en prophase I est court, la cohésion des chromosomes est « neuve » à chaque cycle. Chez la femelle : toutes les mitoses pré-méiotiques sont terminées \textbf{avant la naissance}. Les ovocytes I entrent en prophase I et y restent bloqués \textbf{pendant des décennies} (12 à 50 ans). La cohésion établie pendant la vie fœtale doit être maintenie pendant toute cette durée ; son usure entraîne des non-disjonctions, donc des aneuploïdies.
\textbf{4.} 1 spermatogonie $\xrightarrow{10 \text{ mitoses}}$ $2^{10} = 1024$ spermatocytes I. Chaque spermatocyte I $\xrightarrow{\text{méiose}}$ 4 spermatozoïdes. Total = $1024 \times 4 = \mathbf{4096}$ spermatozoïdes.
\end{corrige}

\coursec{La fécondation}

Dans la section précédente, tu as vu comment les gonades produisent des gamètes haploïdes : le testicule fabrique en continu des millions de spermatozoïdes mobiles, l'ovaire libère périodiquement un ovocyte II bloqué en métaphase de la méiose II. Mais ces cellules à $n$ chromosomes ne peuvent pas, seules, donner naissance à un nouvel individu. Il leur manque la moitié du patrimoine génétique. Comment la diploïdie est-elle restaurée ? C'est toute la question de la \cle{fécondation}, ultime acte de la reproduction sexuée, que nous allons maintenant étudier étape par étape.

\courssub{Définition et localisation de la fécondation}

\begin{definition}[Fécondation et zygote]
La \cle{fécondation} est l'\textbf{union d'un spermatozoïde} (gamète mâle, $n$ chromosomes) \textbf{et d'un ovule} (gamète femelle, $n$ chromosomes), aboutissant à la formation d'une cellule unique appelée \cle{cellule-œuf} ou \cle{zygote}, qui possède $2n$ chromosomes. Chez l'Homme, la fécondation a lieu dans le \textbf{tiers externe de la trompe de Fallope} (trompe utérine), au niveau de l'ampoule, à proximité de l'ovaire.
\end{definition}

\begin{attention}[Erreur fréquente au Bac]
Ne dis \textbf{jamais} que la fécondation a lieu dans l'utérus ! Chez l'Homme, elle se produit dans la \textbf{trompe de Fallope} (tiers externe, près du pavillon). L'utérus n'intervient que quelques jours plus tard, lors de la \emph{nidation} de l'embryon (stade blastocyste). Confondre les deux lieux est une faute classique qui coûte des points.
\end{attention}

\courssub{Un fait expérimental historique : les observations de Hertwig (1875)}

Longtemps, la nature exacte de la fécondation est restée mystérieuse : s'agit-il d'un simple « contact stimulant » ou d'une véritable fusion cellulaire ?

\begin{experience}[Les observations d'Oscar Hertwig chez l'oursin (1875)]
\textbf{Protocole.} Hertwig choisit l'oursin, animal marin dont la fécondation est \emph{externe} et dont les œufs sont transparents : on peut donc observer directement au microscope ce qui se passe. Il mélange des spermatozoïdes et des ovules d'oursin dans une goutte d'eau de mer et suit la scène en direct.

\textbf{Observations.}
\begin{itemize}
\item Un seul spermatozoïde pénètre dans l'ovule, malgré la présence de milliers de congénères ;
\item on observe ensuite \textbf{deux noyaux} dans le cytoplasme de l'œuf : le noyau du spermatozoïde et celui de l'ovule ;
\item ces deux noyaux se rapprochent puis \textbf{fusionnent} en un noyau unique ;
\item l'œuf se divise alors et donne un embryon.
\end{itemize}

\textbf{Interprétation.} La fécondation est un \textbf{événement cellulaire précis et unique} : la fusion de deux noyaux haploïdes ($n + n$) restaure un noyau diploïde ($2n$). Hertwig démontre ainsi que la transmission de l'hérédité passe par les \emph{noyaux} des gamètes, et non par un fluide ou un simple contact.

\textbf{Conclusion.} La fécondation est la \textbf{fusion de deux cellules haploïdes} aboutissant à une cellule diploïde ; c'est le complément obligatoire de la méiose.
\end{experience}

\courssub{Les étapes de la fécondation chez les Mammifères}

Chez les Mammifères, la fécondation est \emph{interne} et l'ovocyte II est entouré de barrières protectrices : la \textbf{couronne radiée} (cellules folliculaires externes riches en acide hyaluronique) et la \textbf{zone pellucide} (matrice extracellulaire glycoprotéique acellulaire, riche en ZP3). Le spermatozoïde doit franchir ces obstacles dans un ordre précis.

\textbf{Étape 1 : rencontre et reconnaissance.}\par
Propulsés par l'éjaculation puis guidés par la chimiotaxie et la thermotaxie dans les voies génitales femelles, les spermatozoïdes traversent d'abord la couronne radiée grâce aux mouvements de leur flagelle et à l'action de l'\textbf{hyaluronidase} (enzyme acrosomique qui dépolymérise l'acide hyaluronique). Le premier contact avec la \cle{zone pellucide} déclenche une \textbf{reconnaissance spécifique} : des récepteurs de la membrane plasmique du spermatozoïde (protéines de la famille Izumo) se fixent sur les glycoprotéines \textbf{ZP3} de la zone pellucide. Cette reconnaissance est \textbf{spécifique de l'espèce} : un spermatozoïde d'une autre espèce ne peut généralement pas se fixer, ce qui évite les fécondations inter-espèces.

\textbf{Étape 2 : la réaction acrosomique.}\par

\begin{definition}[Réaction acrosomique]
La \cle{réaction acrosomique} est l'exocytose de l'\textbf{acrosome} (vésicule dérivée de l'appareil de Golgi, située à l'extrémité antérieure de la tête du spermatozoïde) : sa membrane externe fusionne avec la membrane plasmique du spermatozoïde, provoquant la \textbf{libération d'enzymes protéolytiques} (principalement l'\textbf{acrosine} et d'autres protéases) qui digèrent localement la zone pellucide et ouvrent un passage au spermatozoïde.
\end{definition}

Grâce à cette digestion enzymatique localisée, le spermatozoïde traverse la zone pellucide et atteint la membrane plasmique de l'ovocyte II (appelée \emph{oolème}).

\textbf{Étape 3 : fusion des membranes et blocage de la polyspermie.}\par
La membrane plasmique du spermatozoïde (région post-acrosomique) fusionne avec l'oolème : le noyau mâle (déchromatiné) et le centriole proximal (qui organisera le fuseau de la première mitose) pénètrent dans le cytoplasme. Or, des centaines de spermatozoïdes peuvent être au contact de la zone pellucide au même instant. Si plusieurs noyaux mâles entraient, le zygote serait triploïde ($3n$) ou plus : un tel embryon n'est pas viable. L'entrée du \textbf{premier} spermatozoïde déclenche immédiatement une \textbf{réaction corticale} : les granules corticaux (vésicules sous-corticales) libèrent leur contenu (protéases, peroxydases) dans l'espace périvitellin, ce qui modifie les récepteurs ZP3 (perte de la capacité de fixation) et durcit la zone pellucide (réticulation des glycoprotéines). C'est le \textbf{blocage de la polyspermie}.

\begin{propriete}[Blocage de la polyspermie]
Le blocage de la polyspermie garantit qu'\textbf{un seul} noyau mâle fusionne avec le noyau femelle. Il assure ainsi la formation d'un zygote \textbf{exactement diploïde} ($2n$), condition indispensable au développement normal de l'embryon.
\end{propriete}

\textbf{Étape 4 : achèvement de la méiose II et amphimixie.}\par
Rappelle-toi : l'ovocyte II expulsé de l'ovaire est bloqué en \textbf{métaphase II} (activité MPF maintenue par le facteur cytostatique CSF). L'entrée du spermatozoïde (via un facteur soluble PLC$\zeta$) déclenche des oscillations calciques qui lèvent ce blocage : l'ovocyte II \textbf{achève sa méiose II} (ségrégation des chromatides sœurs) et émet un \textbf{second globule polaire}. On observe alors dans le cytoplasme deux noyaux haploïdes, appelés \cle{pronuclei} : le \textbf{pronucleus mâle} (issu du noyau spermatique, décondensé) et le \textbf{pronucleus femelle} (issu de la méiose II). Ces deux pronuclei se rapprochent au centre de la cellule, leurs enveloppes nucléaires se fragmentent et leurs chromosomes s'alignent sur un fuseau mitotique commun : c'est l'\textbf{amphimixie}.

\begin{definition}[Amphimixie]
L'\cle{amphimixie} est la \textbf{réunion des deux lots de chromosomes} (paternel et maternel) sur un même fuseau mitotique au cours de la fécondation, consécutive à la disparition des enveloppes des pronuclei. Elle réunit en un seul noyau les deux lots haploïdes de chromosomes ($n + n = 2n$) : le \textbf{zygote diploïde} est formé et peut entamer sa première mitose.
\end{definition}

\begin{popfigure}
\begin{center}
\begin{tikzpicture}[scale=0.85, every node/.style={font=\small}]
% Étape 1 : approche
\begin{scope}[shift={(0,0)}]
\draw[fill=poppinkL] (0,0) circle (1.1);
\draw[fill=poppurpleL,thick,draw=poppurple] (0,0) circle (0.85);
\draw[fill=white] (0,0) circle (0.65);
\fill[poppurple] (0,0) circle (0.15);
\foreach \a in {20,60,100,140,200,250,300,340}{\fill[popgold] (\a:1.25) circle (0.09);}
\draw[fill=popblue] (2.4,0.5) ellipse (0.22 and 0.13);
\draw[popblue,thick] (2.62,0.5) .. controls (3.2,0.9) and (3.4,0.1) .. (3.8,0.45);
\node[align=center] at (1.4,-1.9) {\textbf{1. Approche et}\\ \textbf{reconnaissance (ZP3)}};
\end{scope}
% Étape 2 : réaction acrosomique
\begin{scope}[shift={(6.5,0)}]
\draw[fill=poppinkL] (0,0) circle (1.1);
\draw[fill=poppurpleL,thick,draw=poppurple] (0,0) circle (0.85);
\draw[fill=white] (0,0) circle (0.65);
\fill[poppurple] (0,0) circle (0.15);
\draw[fill=poporange] (1.4,0) ellipse (0.22 and 0.13);
\draw[poporange,thick] (1.62,0) .. controls (2.15,0.4) and (2.35,-0.3) .. (2.75,0.1);
\foreach \p in {(1.1,0.25),(1.05,-0.2),(1.0,0.05)}{\fill[popgreen] \p circle (0.05);}
\node[align=center] at (1.2,-1.9) {\textbf{2. Réaction}\\ \textbf{acrosomique (acrosine)}};
\end{scope}
% Étape 3 : pénétration
\begin{scope}[shift={(13,0)}]
\draw[fill=poppinkL] (0,0) circle (1.1);
\draw[fill=poppurpleL,thick,draw=poppurple] (0,0) circle (0.85);
\draw[fill=white] (0,0) circle (0.65);
\fill[poppurple] (-0.3,0) circle (0.15);
\draw[fill=popblue] (0.5,0) ellipse (0.2 and 0.12);
\draw[popblue,thick] (0.7,0) .. controls (1.25,0.4) and (1.45,-0.3) .. (1.85,0.1);
\node[align=center] at (0.9,-1.9) {\textbf{3. Fusion membranaire}\\ \textbf{et blocage polyspermie}};
\end{scope}
% Étape 4 : pronuclei
\begin{scope}[shift={(3.2,-5)}]
\draw[fill=poppinkL] (0,0) circle (1.1);
\draw[fill=white] (0,0) circle (0.9);
\draw[fill=popblueL,draw=popblue,thick] (-0.35,0) circle (0.3);
\draw[fill=poppink,draw=poppurple,thick] (0.4,0) circle (0.3);
\fill[popgold] (0.8,0.7) circle (0.1);
\node[align=center] at (0,-1.9) {\textbf{4. Deux pronuclei}\\ \textbf{(+ 2\textsuperscript{e} globule polaire)}};
\end{scope}
% Étape 5 : zygote
\begin{scope}[shift={(9.5,-5)}]
\draw[fill=popgreenL,draw=popgreen,thick] (0,0) circle (1.1);
\draw[fill=white] (0,0) circle (0.9);
\draw[fill=poppurpleL,draw=poppurple,thick] (0,0) circle (0.45);
\foreach \p in {(-0.15,0.12),(0.15,0.12),(-0.15,-0.15),(0.15,-0.15)}{\draw[popdark,thick] \p -- ++(0.1,0.15);}
\node[align=center] at (0,-1.9) {\textbf{5. Amphimixie :}\\ \textbf{zygote $2n$}};
\end{scope}
% Flèches de transition
\draw[-{Stealth},thick,popmuted] (4.6,0) -- (5.1,0);
\draw[-{Stealth},thick,popmuted] (11.2,0) -- (11.7,0);
\draw[-{Stealth},thick,popmuted] (13.5,-1.5) -- (11.2,-3.8);
\draw[-{Stealth},thick,popmuted] (5.5,-5) -- (7.2,-5);
\end{tikzpicture}
\end{center}
\legende{Étapes séquentielles de la fécondation chez les Mammifères. 1 : le spermatozoïde (bleu) franchit la couronne radiée (cellules dorées) et se fixe sur la zone pellucide (anneau violet, récepteurs ZP3). 2 : la réaction acrosomique libère l'acrosine (points verts) qui digère localement la zone pellucide. 3 : la tête du spermatozoïde fusionne avec l'oolème ; la réaction corticale modifie la zone pellucide (blocage de la polyspermie). 4 : les deux pronuclei haploïdes (mâle en bleu clair, femelle en rose) côtoient le second globule polaire (doré). 5 : l'amphimixie réunit les chromosomes sur un fuseau commun : le zygote est diploïde ($2n$).}
\end{popfigure}

\begin{methode}[Identifier les étapes de la fécondation sur un document]
Face à une micrographie ou un schéma d'ovocyte en fécondation, procède par étapes :
\begin{enumerate}
\item \textbf{Repère l'acrosome} : s'il est intact (capuchon lisse à la pointe de la tête), la réaction acrosomique n'a pas encore eu lieu ; s'il a disparu ou est éclaté (vésicules résiduelles), elle a eu lieu.
\item \textbf{Observe la position du spermatozoïde} : hors de la zone pellucide (approche/reconnaissance), dans l'épaisseur de la zone pellucide (pénétration acrosomique), ou dans le cytoplasme (fusion membranaire réalisée).
\item \textbf{Compte les noyaux} : un seul noyau visible (noyau de l'ovocyte II en métaphase II) = ovocyte non fécondé ; \textbf{deux pronuclei} distincts (souvent centrés, enveloppe visible) = fécondation en cours (avant amphimixie) ; un seul noyau volumineux avec chromosomes condensés sur un fuseau = zygote en première mitose.
\item \textbf{Cherche les globules polaires} : la présence d'un \textbf{second globule polaire} indique que la méiose II s'est achevée, donc que la fécondation a été déclenchée.
\item \textbf{Conclus} en nommant l'étape précise et en justifiant par les indices observés.
\end{enumerate}
\end{methode}

\courssub{La nécessité de la fécondation : un couple indissociable avec la méiose}

Pourquoi la fécondation est-elle indispensable ? Deux raisons complémentaires :

\begin{itemize}
\item \textbf{Elle restaure la diploïdie propre à l'espèce.} La méiose divise par deux le nombre de chromosomes ($2n \rightarrow n$). Sans fécondation, chaque génération verrait ce nombre diminuer de moitié, ce qui est incompatible avec la stabilité du caryotype. La fécondation rétablit $2n$ : chez l'Homme, $23 + 23 = 46$ chromosomes.
\item \textbf{Elle réunit deux lots d'allèles différents.} Le zygote reçoit un lot paternel et un lot maternel : il est génétiquement \emph{nouveau}, différent de chacun de ses parents. Associée aux brassages de la méiose (croisements, répartition aléatoire), la fécondation est donc aussi une \textbf{source majeure de diversité génétique}.
\end{itemize}

\begin{aretenir}[Alternance méiose / fécondation]
La \textbf{constance du nombre de chromosomes} d'une génération à l'autre, au sein d'une espèce, résulte de l'\textbf{alternance de deux événements complémentaires} : la \textbf{méiose} (réduction : $2n \rightarrow n$, lors de la gamétogenèse) et la \textbf{fécondation} (restauration : $n + n \rightarrow 2n$). Ni l'une ni l'autre ne peut être supprimée sans rompre ce cycle.
\end{aretenir}

\begin{popfigure}
\begin{center}
\begin{tikzpicture}[scale=1.0, every node/.style={font=\small}]
% Adultes
\node[draw=popdark,thick,fill=popblueL,rounded corners,minimum width=2.8cm,minimum height=0.9cm] (pere) at (0,0) {Homme adulte $2n$};
\node[draw=popdark,thick,fill=poppinkL,rounded corners,minimum width=2.8cm,minimum height=0.9cm] (mere) at (7.5,0) {Femme adulte $2n$};
% Gamètes
\node[draw=popblue,thick,fill=white,rounded corners,minimum width=2.8cm,minimum height=0.9cm] (spz) at (0,-3) {Spermatozoïde $n$};
\node[draw=poppurple,thick,fill=white,rounded corners,minimum width=2.8cm,minimum height=0.9cm] (ov) at (7.5,-3) {Ovocyte II $n$};
% Zygote
\node[draw=popgreen,very thick,fill=popgreenL,rounded corners,minimum width=3.5cm,minimum height=1.1cm] (zyg) at (3.75,-6) {Zygote $2n$};
% Flèches méiose
\draw[-{Stealth},very thick,poporange] (pere) -- (spz) node[midway,left,align=center]{méiose\\ (réduction $2n \to n$)};
\draw[-{Stealth},very thick,poporange] (mere) -- (ov) node[midway,right,align=center]{méiose\\ (réduction $2n \to n$)};
% Flèches fécondation
\draw[-{Stealth},very thick,popgreen] (spz) -- (zyg);
\draw[-{Stealth},very thick,popgreen] (ov) -- (zyg);
\node[popgreen,align=center] at (3.75,-4.5) {\textbf{fécondation} (restauration $n+n=2n$)};
% Flèche développement
\draw[-{Stealth},very thick,popdark,dashed] (zyg.west) .. controls (-3,-6.5) and (-3,-1) .. (pere.south west);
\draw[-{Stealth},very thick,popdark,dashed] (zyg.east) .. controls (10.5,-6.5) and (10.5,-1) .. (mere.south east);
\node[popdark,align=center] at (-3.2,-3.6) {mitoses et\\ développement};
\node[popdark,align=center] at (10.7,-3.6) {mitoses et\\ développement};
\end{tikzpicture}
\end{center}
\legende{Cycle chromosomique de la reproduction sexuée. La méiose (flèches orange) réduit le stock chromosomique de moitié lors de la formation des gamètes ($2n \rightarrow n$) ; la fécondation (flèches vertes) restaure la diploïdie dans le zygote ($n + n = 2n$). Les mitoses (flèches noires en pointillés) assurent ensuite le développement de l'adulte, dont les gonades produiront à leur tour des gamètes : le cycle est bouclé.}
\end{popfigure}

\begin{exemplebox}[Une course contre la montre : les chiffres de la fécondation]
Lors d'un rapport sexuel, l'éjaculat humain contient environ \textbf{200 à 500 millions de spermatozoïdes} ($\sim 10^8$). La plupart périssent dans le vagin (milieu acide, pH $\approx$ 4--5) ou sont éliminés dans l'utérus (phagocytose) ; seuls \textbf{quelques centaines} (environ $10^2$--$10^3$) parviennent jusqu'au tiers externe de la trompe où se trouve l'ovocyte II. Et \textbf{un seul} le féconde, grâce au blocage de la polyspermie. La fécondation est donc un événement d'une sélectivité extraordinaire : probabilité de l'ordre de $1$ sur $10^6$ à $10^8$ pour chaque spermatozoïde éjaculé !
\end{exemplebox}

\begin{exoresolu}[Restituer le cycle chromosomique]
\textbf{Énoncé.} Chez le chimpanzé, la cellule-œuf possède 48 chromosomes.
\begin{enumerate}
\item Combien de chromosomes possèdent ses spermatozoïdes et ses ovules ? Justifie.
\item Nomme et décris les deux événements qui assurent la constance du nombre de chromosomes d'une génération à l'autre.
\end{enumerate}
\tcblower
\textbf{Solution détaillée.}
\begin{enumerate}
\item Les gamètes sont issus de la \textbf{méiose}, division réductionnelle qui divise par deux le nombre de chromosomes. Spermatozoïdes et ovules possèdent donc $48 / 2 = \textbf{24 chromosomes}$ ($n = 24$).
\item Les deux événements sont :
\begin{itemize}
\item la \textbf{méiose}, qui a lieu dans les gonades au cours de la gamétogenèse : elle fait passer les cellules sexuelles de $2n = 48$ à $n = 24$ chromosomes (\emph{réduction}) ;
\item la \textbf{fécondation}, union d'un spermatozoïde ($n = 24$) et d'un ovule ($n = 24$) dans la trompe : l'amphimixie restaure $2n = 48$ chromosomes dans le zygote (\emph{restauration}).
\end{itemize}
L'alternance de ces deux événements garantit la stabilité du caryotype de l'espèce au fil des générations.
\end{enumerate}
\end{exoresolu}

\begin{attention}[Pièges de rédaction au Bac]
\begin{itemize}
\item La fécondation n'est pas la simple « rencontre » des gamètes mais leur \textbf{fusion} (fusion des membranes plasmiques puis réunion des chromosomes) : le mot « union » ou « fusion » dans la définition est exigé.
\item Ne confonds pas \textbf{fusion des membranes} (entrée du noyau spermatique dans l'ooplasme, étape 3) et \textbf{amphimixie} (réunion des chromosomes sur un fuseau commun après disparition des enveloppes des pronuclei, étape 4) : ce sont deux événements successifs et distincts.
\item Le zygote n'est diploïde qu'\textbf{après} l'amphimixie ; tant que les deux pronuclei sont individualisés, on ne parle pas encore de noyau diploïde unique (le caryotype est $2n$ mais la structure nucléaire est encore double).
\item Précise toujours \textbf{ovocyte II} (et non ovule mature) pour la cellule ovulaire au moment de la fécondation chez les Mammifères ; l'ovule (gamète femelle achevé) n'existe qu'après émission du second globule polaire.
\end{itemize}
\end{attention}

\courssub{Compléments utiles pour le Bac (hors programme strict)}

\begin{savaistu}[La fécondation in vitro (FIV) et le sexe de l'enfant]
\textbf{La FIV.} Le 25 juillet 1978 naît en Angleterre \textbf{Louise Brown}, premier « bébé-éprouvette » : ses gamètes ont été réunis \emph{in vitro} (dans une boîte de Pétri, milieu de culture), puis l'embryon (stade 2 à 8 cellules ou blastocyste) a été replacé dans l'utérus de sa mère. Cette technique, qui applique directement les connaissances sur la réaction acrosomique (sélection des spermatozoïdes mobiles, parfois injection intracytoplasmique -- ICSI) et l'amphimixie, a depuis permis la naissance de \textbf{plusieurs millions d'enfants} dans le monde, y compris en Afrique et au Cameroun, où des centres de procréation médicalement assistée (PMA) aident les couples confrontés à l'infertilité (obstruction des trompes, oligospermie, endométriose).

\textbf{Le déterminisme chromosomique du sexe.} Chez l'Homme, l'ovocyte II porte toujours un chromosome \textbf{X}. Le spermatozoïde porte soit un chromosome \textbf{X}, soit un chromosome \textbf{Y} (hétérogamétie mâle). C'est donc le \textbf{spermatozoïde} qui détermine le sexe chromosomique de l'enfant : fécondation par un spermatozoïde $X \rightarrow$ zygote $XX$ (fille) ; fécondation par un spermatozoïde $Y \rightarrow$ zygote $XY$ (garçon). Chaque conception a donc une probabilité de $1/2$ d'être une fille et $1/2$ d'être un garçon, indépendamment des grossesses précédentes.
\end{savaistu}

\begin{exercice}[Pour t'entraîner]
Sur une micrographie d'ovocyte de Mammifère, on observe : une zone pellucide épaissie et modifiée, un second globule polaire visible dans l'espace périvitellin, et deux noyaux distincts de taille comparable, centrés dans le cytoplasme.
\begin{enumerate}
\item Identifie l'étape représentée en justifiant par trois indices morphologiques précis.
\item Que va-t-il se passer ensuite ? Nomme l'événement et précise son résultat chromosomique et structural.
\end{enumerate}
\end{exercice}

\begin{corrige}[Exercice]
\begin{enumerate}
\item Il s'agit de l'étape des \textbf{deux pronuclei} (fécondation en cours, avant amphimixie). Indices : (a) la présence du \textbf{second globule polaire} montre que la méiose II s'est achevée, donc qu'un spermatozoïde a pénétré et a déclenché les oscillations calciques ; (b) les \textbf{deux noyaux distincts} (pronucleus mâle et pronucleus femelle) avec enveloppe nucléaire visible sont caractéristiques de cette étape ; (c) la \textbf{zone pellucide épaissie/modifiée} traduit le blocage de la polyspermie déjà déclenché par la réaction corticale.
\item Les deux pronuclei vont migrer l'un vers l'autre, leurs enveloppes nucléaires vont se fragmenter (néoclase) et leurs chromosomes vont s'aligner sur un fuseau mitotique commun : c'est l'\textbf{amphimixie}. Les deux lots haploïdes se réunissent en un noyau unique : le zygote devient structurellement \textbf{diploïde} ($n + n = 2n$, un seul noyau) et peut entamer sa première mitose de segmentation.
\end{enumerate}
\end{corrige}

\coursec{Bilan : méiose et fécondation, un couple indissociable}

Nous venons de voir que la fécondation, ultime acte de la reproduction sexuée, réunit deux gamètes haploïdes pour former un zygote diploïde. Mais pourquoi les gamètes sont-ils haploïdes ? Et pourquoi la fécondation est-elle indispensable ? Il est temps de rassembler toutes les pièces du puzzle : la méiose et la fécondation forment un \cle{couple indissociable}, deux mécanismes complémentaires qui assurent à la fois le \cle{maintien de l'identité de l'espèce} et le \cle{renouvellement de l'information génétique} à chaque génération. C'est précisément le c\oe ur de notre famille de situations.

\courssub{Le problème initial : comment transmettre l'information génétique sans l'amplifier indéfiniment ?}

Posons le problème comme le ferait un chercheur. Chez les Mammifères, la reproduction sexuée met en jeu deux cellules reproductrices, les gamètes, qui fusionnent pour donner la cellule-\oe uf. Or chaque cellule somatique d'un individu possède un stock chromosomique diploïde ($2n$). Raisonnons par l'absurde :

\begin{itemize}
\item Si les gamètes étaient diploïdes ($2n$), leur fusion produirait un zygote à $4n$ chromosomes, puis $8n$ à la génération suivante : le nombre de chromosomes doublerait à chaque génération, ce qui est incompatible avec la \cle{constance du caryotype} observée dans toutes les espèces.
\item Si les gamètes étaient haploïdes ($n$) mais ne fusionnaient jamais, il n'y aurait pas de reproduction sexuée du tout.
\end{itemize}

La seule solution logique est donc un \textbf{double mécanisme} : une division réductionnelle qui divise par deux le stock chromosomique (la \cle{méiose} : $2n \rightarrow n$), suivie d'une fusion qui le restaure (la \cle{fécondation} : $n + n \rightarrow 2n$). Ni la méiose seule, ni la fécondation seule ne peut assurer la reproduction : la première sans la seconde conduirait à des individus haploïdes inviables chez les Mammifères ; la seconde sans la première provoquerait une explosion du nombre de chromosomes.

\begin{propriete}[Alternance obligatoire méiose/fécondation]
Chez toute espèce à reproduction sexuée, le cycle de vie comporte obligatoirement l'alternance de deux événements complémentaires :
\begin{itemize}
\item la \cle{méiose}, qui réduit de moitié le nombre de chromosomes ($2n \rightarrow n$) lors de la formation des gamètes ;
\item la \cle{fécondation}, qui restaure le nombre diploïde ($n + n \rightarrow 2n$) lors de la formation du zygote.
\end{itemize}
Cette alternance assure simultanément la \textbf{constance du nombre de chromosomes} (identité de l'espèce) et la \textbf{diversité génétique} (renouvellement de l'information) d'une génération à l'autre.
\end{propriete}

\begin{popfigure}
\begin{center}
\begin{tikzpicture}[scale=1, every node/.style={font=\small}]
% Nœuds du cycle
\node[draw=popblue, fill=popblueL, rounded corners, thick, minimum width=3.4cm, minimum height=1cm, align=center] (adulte) at (0,3.2) {\textbf{Adulte} ($2n$)\\ mâle et femelle};
\node[draw=poppurple, fill=poppurpleL, rounded corners, thick, minimum width=3.2cm, minimum height=1cm, align=center] (gametes) at (5.4,0) {\textbf{Gamètes} ($n$)\\ spermatozoïde + ovule};
\node[draw=popgreen, fill=popgreenL, rounded corners, thick, minimum width=3.2cm, minimum height=1cm, align=center] (zygote) at (-5.4,0) {\textbf{Zygote} ($2n$)\\ cellule-\oe uf};
% Flèches
\draw[-{Stealth[length=3mm]}, very thick, poppurple] (adulte.east) to[bend left=15] node[midway, above right, align=center, font=\small\itshape]{Méiose\\ (gamétogenèse)\\ $2n \rightarrow n$} (gametes.north);
\draw[-{Stealth[length=3mm]}, very thick, popgreen] (gametes.west) to[bend left=15] node[midway, below, align=center, font=\small\itshape]{Fécondation\\ $n + n \rightarrow 2n$} (zygote.east);
\draw[-{Stealth[length=3mm]}, very thick, popblue] (zygote.north) to[bend left=15] node[midway, above left, align=center, font=\small\itshape]{Développement\\ (mitoses)\\ $2n \rightarrow 2n$} (adulte.west);
\end{tikzpicture}
\end{center}
\legende{Schéma circulaire du cycle de vie d'un Mammifère à reproduction sexuée. La méiose (flèche violette) réduit la ploïdie de moitié lors de la gamétogenèse ; la fécondation (flèche verte) la restaure ; le développement embryonnaire (flèche bleue) s'effectue par mitoses, qui conservent le stock diploïde. Le cycle est bouclé : $2n \rightarrow n \rightarrow 2n$.}
\end{popfigure}

\courssub{La constance du caryotype : vérification chiffrée chez l'Homme}

Vérifions ce modèle avec les données du caryotype humain. Chez l'Homme, $2n = 46$ chromosomes, soit 23 paires de chromosomes homologues (22 paires d'autosomes + 1 paire de chromosomes sexuels).

\begin{center}
\begin{tabularx}{\linewidth}{|l|X|X|X|}
\hline
\rowcolor{popblueL} \textbf{Étape} & \textbf{Cellule concernée} & \textbf{Ploïdie} & \textbf{Chez l'Homme} \\
\hline
Avant gamétogenèse & Cellules germinales (spermatogonies, ovogonies) & $2n$ & 46 chromosomes \\
\hline
Après méiose & Gamètes (spermatozoïde, ovule) & $n$ & 23 chromosomes \\
\hline
Après fécondation & Zygote puis toutes les cellules de l'embryon & $2n$ & 46 chromosomes \\
\hline
\end{tabularx}
\end{center}

Le compte est bon : $46 \rightarrow 23 \rightarrow 46$. Le zygote reçoit 23 chromosomes d'origine paternelle et 23 chromosomes d'origine maternelle, qui s'apparient en 23 paires d'homologues. Ce résultat se \textbf{généralise à toute espèce à reproduction sexuée} : chimpanzé ($2n = 48 \rightarrow n = 24 \rightarrow 48$), chien ($2n = 78$), drosophile ($2n = 8$), maïs ($2n = 20$)\dots{} Le nombre $2n$ est une \cle{caractéristique de l'espèce}, maintenue constante de génération en génération grâce au couple méiose/fécondation.

\begin{methode}[Rédiger une réponse type Bac : « Expliquer la nécessité de la méiose / de la fécondation »]
Face à cette question classique, structure ta réponse en trois temps :
\begin{enumerate}
\item \textbf{Poser le problème} : rappeler que les cellules somatiques de l'organisme sont diploïdes ($2n$) et que le nombre de chromosomes doit rester constant d'une génération à l'autre (constance du caryotype, caractéristique de l'espèce).
\item \textbf{Décrire le mécanisme} : la méiose comporte deux divisions successives précédées d'une seule réplication de l'ADN, ce qui produit des gamètes haploïdes ($n$) ; la fécondation fusionne deux gamètes haploïdes et restaure la diploïdie ($n + n = 2n$).
\item \textbf{Conclure par une conséquence chiffrée} : chez l'Homme, $46 \rightarrow 23 \rightarrow 46$ ; sans méiose, le nombre de chromosomes doublerait à chaque génération ; sans fécondation, pas de zygote diploïde. Les deux mécanismes sont donc \emph{complémentaires et indissociables}.
\end{enumerate}
Pense à illustrer ta réponse d'un schéma simple $2n \rightarrow n \rightarrow 2n$ : il est toujours valorisé.
\end{methode}

\courssub{La diversité génétique : trois sources combinées}

Si la reproduction sexuée ne faisait que conserver le caryotype, tous les individus d'une même fratrie se ressembleraient parfaitement. Or l'observation quotidienne montre le contraire : hors vrais jumeaux, chaque individu est \cle{génétiquement unique}. D'où vient cette diversité ? De \textbf{trois sources qui se combinent}.

\textbf{Première source : le brassage interchromosomique.} En anaphase I de méiose, les paires de chromosomes homologues se séparent, mais l'orientation de chaque paire sur le fuseau est \emph{indépendante} des autres : le chromosome d'origine paternelle d'une paire peut migrer vers un pôle tandis que celui d'une autre paire migre vers l'autre pôle. Pour $n$ paires, le nombre de combinaisons possibles est $2^n$. Chez l'Homme, chaque individu peut donc produire $2^{23} \approx \num{8,4}$ millions de types de gamètes différents par ce seul mécanisme.

\textbf{Deuxième source : le brassage intrachromosomique (crossing-over).} En prophase I, les chromosomes homologues appariés (bivalents) échangent des fragments de chromatides : les \cle{crossing-over}. Les chromosomes transmis aux gamètes sont donc des chromosomes \emph{recombinés}, portant des mosaïques d'allèles d'origine paternelle et maternelle. Ce brassage, quasi infini dans ses résultats, multiplie encore la variété des gamètes.

\textbf{Troisième source : la rencontre aléatoire des gamètes à la fécondation.} Parmi les millions de spermatozoïdes émis, un seul féconde l'ovule, et cette rencontre est due au hasard. La diversité des zygotes possibles résulte donc du \emph{produit} des diversités des deux gamètes.

\begin{exemplebox}[Combien de zygotes humains différents ?]
Par le seul brassage interchromosomique, un homme peut produire $2^{23}$ spermatozoïdes génétiquement distincts et une femme $2^{23}$ ovules distincts. Le nombre de zygotes différents qu'un couple peut théoriquement engendrer est donc :
\[ 2^{23} \times 2^{23} = 2^{46} \approx 7 \times 10^{13} \]
soit environ \textbf{70 000 milliards} de combinaisons possibles, \emph{sans même compter} les crossing-over, qui rendent ce nombre pratiquement infini. C'est pourquoi, hors vrais jumeaux (issus d'un même zygote), deux individus génétiquement identiques n'existent pas.
\end{exemplebox}

\begin{attention}[Erreur fréquente au Bac]
Ne réduis pas la diversité génétique à la seule méiose ! Beaucoup d'élèves écrivent : « la méiose assure la diversité génétique ». C'est \textbf{incomplet} : la méiose (brassages inter- et intrachromosomiques) diversifie les \emph{gamètes}, mais la \cle{fécondation}, par la rencontre aléatoire de deux gamètes parmi des millions, amplifie considérablement cette diversité au niveau du \emph{zygote}. La formulation rigoureuse attendue est : « la diversité génétique résulte des brassages méiotiques \emph{et} de la rencontre aléatoire des gamètes lors de la fécondation ».
\end{attention}

\begin{popfigure}
\begin{center}
\begin{tikzpicture}
\begin{axis}[
    width=0.92\linewidth, height=6.2cm,
    xlabel={Étapes du cycle de reproduction},
    ylabel={Quantité d'ADN par cellule (unités arbitraires)},
    symbolic x coords={Spermatogonie, Spermatocyte I, Spermatocyte II, Spermatide, Spermatozoïde, Fécondation, Zygote},
    xtick=data,
    x tick label style={rotate=25, anchor=east, font=\footnotesize},
    ymin=0, ymax=5.2,
    ytick={1,2,4},
    yticklabels={$Q$, $2Q$, $4Q$},
    grid=major, grid style={popline!40},
    legend style={at={(0.02,0.97)}, anchor=north west, font=\footnotesize},
    clip=false
]
\addplot[popcurve, mark=*, mark size=1.8pt] coordinates {
    (Spermatogonie,2) (Spermatocyte I,4) (Spermatocyte II,2) (Spermatide,1) (Spermatozoïde,1) (Fécondation,2) (Zygote,2)
};
\addlegendentry{Quantité d'ADN par cellule}
% Annotations placées au milieu des pentes
\node[font=\footnotesize, align=center, popdark] at (axis cs:Spermatocyte I,3) {Réplication\\ (phase S)};
\node[font=\footnotesize, align=center, popdark] at (axis cs:Spermatocyte I,3) {Méiose I};
\node[font=\footnotesize, align=center, popdark] at (axis cs:Spermatocyte II,1.5) {Méiose II};
\node[font=\footnotesize, align=center, popgreen] at (axis cs:Fécondation,2.8) {$n+n \rightarrow 2n$};
\end{axis}
\end{tikzpicture}
\end{center}
\legende{Courbe synthétique de la quantité d'ADN par cellule au cours du cycle complet (spermatogenèse puis fécondation). La réplication double l'ADN ($2Q \rightarrow 4Q$) avant la méiose ; les deux divisions successives le ramènent à $Q$ dans les gamètes ; la fécondation restaure $2Q$ dans le zygote. La ploïdie suit le même rythme : $2n \rightarrow n \rightarrow 2n$.}
\end{popfigure}

\courssub{Lien avec la famille de situations : identité et renouvellement}

Notre famille de situations s'intitule « Maintien de l'identité cellulaire au cours de la transmission de l'information génétique ». Le bilan de cette leçon en éclaire pleinement le sens :

\begin{itemize}
\item \textbf{Maintien de l'identité} : grâce à l'alternance méiose/fécondation, le nombre de chromosomes $2n$, signature de l'espèce, est rigoureusement conservé à chaque génération. Le zygote reçoit un caryotype complet et équilibré, moitié paternel, moitié maternel.
\item \textbf{Renouvellement de l'information} : grâce aux brassages méiotiques et à la fécondation aléatoire, l'\emph{assemblage} des allèles est nouveau à chaque génération. L'espèce reste elle-même, mais chaque individu est une combinaison inédite.
\end{itemize}

La reproduction sexuée réalise donc un double exploit : \emph{fidélité} dans la quantité d'information transmise, \emph{innovation} dans son contenu. Cette diversité est d'ailleurs un atout évolutif majeur : elle fournit à la sélection naturelle une matière variée, ce qui explique le succès de la reproduction sexuée chez la quasi-totalité des espèces complexes.

\courssub{Conséquences pathologiques : quand la méiose se trompe}

La précision du mécanisme de méiose a son revers : une erreur, même rare, peut avoir des conséquences graves. Il arrive qu'une paire de chromosomes homologues (en méiose I) ou deux chromatides s\oe urs (en méiose II) ne se séparent pas : c'est la \cle{non-disjonction} chromosomique. Il en résulte des gamètes anormaux, portant soit un chromosome en excès ($n+1$), soit un chromosome en défaut ($n-1$). La fécondation d'un tel gamète par un gamète normal produit un zygote \cle{trisomique} ($2n+1$) ou \cle{monosomique} ($2n-1$), dont le caryotype déséquilibré perturbe le développement.

\begin{savaistu}[La trisomie 21 et le dépistage prénatal]
La \cle{trisomie 21} (présence de trois chromosomes 21 au lieu de deux) est l'anomalie chromosomique viable la plus fréquente chez l'Homme, avec une incidence d'environ 1 naissance sur 700. Elle résulte dans environ 95 \% des cas d'une non-disjonction du chromosome 21 lors de l'\textbf{ovogenèse} — rappelons que l'ovocyte I reste bloqué en prophase I parfois plus de 40 ans, ce qui fragilise l'appareil de division et explique que le risque augmente avec l'âge maternel (environ 1/100 après 40 ans). Le diagnostic peut être établi pendant la grossesse par analyse du caryotype de cellules f\oe tales (amniocentèse) ou, plus récemment, par dépistage non invasif de l'ADN f\oe tal circulant dans le sang maternel. Dans les maternités camerounaises comme dans le monde entier, le suivi prénatal permet d'informer et d'accompagner les familles.
\end{savaistu}

\begin{exoresolu}[Analyse d'un caryotype anormal]
\textbf{Énoncé.} Le caryotype d'un nouveau-né présente 47 chromosomes, avec trois exemplaires du chromosome 21. (1) Nommer cette anomalie et son mécanisme d'apparition. (2) Expliquer, en vous appuyant sur les valeurs de ploïdie, comment un tel caryotype peut résulter de la fécondation. (3) Cette observation contredit-elle le principe de constance du caryotype ?
\tcblower
\textbf{Solution détaillée.}
(1) Il s'agit d'une \textbf{trisomie 21}, conséquence d'une \textbf{non-disjonction} du chromosome 21 au cours de la méiose : la paire d'homologues (méiose I) ou les chromatides s\oe urs (méiose II) ne se sont pas séparés.

(2) La non-disjonction produit un gamète anormal à $n+1 = 24$ chromosomes (portant deux chromosomes 21). Sa fécondation par un gamète normal à $n = 23$ chromosomes donne un zygote à $24 + 23 = 47$ chromosomes, dont trois chromosomes 21 : $2n+1 = 47$.

(3) Non : cette observation \emph{confirme} le principe. La constance du caryotype repose précisément sur l'exactitude de la méiose et de la fécondation ; la trisomie 21 illustre ce qui se produit lorsque ce mécanisme faillit. Le déséquilibre du stock chromosomique entraîne des troubles du développement, ce qui montre à quel point la régularité du couple méiose/fécondation est indispensable.
\end{exoresolu}

\courssub{Tableau de synthèse final}

\begin{center}
\begin{tabularx}{\linewidth}{|l|X|X|X|X|}
\hline
\rowcolor{popblueL} \textbf{Étape} & \textbf{Lieu} & \textbf{Mécanisme} & \textbf{Ploïdie} & \textbf{ADN/cellule} \\
\hline
Cellules germinales & Gonades (testicules, ovaires) & Mitoses de multiplication & $2n$ & $2Q$ puis $4Q$ après réplication \\
\hline
Méiose I & Gonades & Séparation des homologues (brassages) & $2n \rightarrow n$ & $4Q \rightarrow 2Q$ \\
\hline
Méiose II & Gonades & Séparation des chromatides & $n$ & $2Q \rightarrow Q$ \\
\hline
Gamètes & Voies génitales & Différenciation (spermiogenèse) & $n$ & $Q$ \\
\hline
Fécondation & Trompe de Fallope (ampoule) & Fusion des noyaux (caryogamie) & $n + n \rightarrow 2n$ & $Q + Q \rightarrow 2Q$ \\
\hline
Zygote puis embryon & Utérus & Mitoses de développement & $2n$ & $2Q \leftrightarrow 4Q$ (cycle cellulaire) \\
\hline
\end{tabularx}
\end{center}

\begin{aretenir}[L'essentiel de la leçon]
\begin{itemize}
\item La reproduction sexuée repose sur l'\textbf{alternance obligatoire} de la méiose ($2n \rightarrow n$) et de la fécondation ($n + n \rightarrow 2n$) : ni l'une ni l'autre ne suffit seule.
\item Chez l'Homme : $46 \rightarrow 23 \rightarrow 46$ chromosomes ; la constance de $2n$ assure l'\textbf{identité de l'espèce}.
\item La \textbf{diversité génétique} a trois sources combinées : brassage interchromosomique ($2^{23}$ gamètes possibles), crossing-over, rencontre aléatoire des gamètes ($2^{46} \approx 7 \times 10^{13}$ zygotes possibles par couple).
\item Une \textbf{non-disjonction} méiotique produit des gamètes anormaux ($n+1$ ou $n-1$), source d'anomalies chromosomiques comme la trisomie 21 ($2n+1 = 47$).
\end{itemize}
\end{aretenir}

\begin{exercice}[Pour t'entraîner]
Chez le chimpanzé, $2n = 48$ chromosomes. (1) Indiquer le nombre de chromosomes de ses gamètes et de son zygote. (2) Calculer le nombre de gamètes différents qu'un chimpanzé peut produire par brassage interchromosomique seul. (3) Expliquer en quelques lignes pourquoi la fécondation est indispensable au maintien du caryotype de l'espèce.
\end{exercice}

\begin{corrige}[Exercice : le caryotype du chimpanzé]
(1) Les gamètes du chimpanzé sont haploïdes : $n = 24$ chromosomes ; le zygote est diploïde : $2n = 48$ chromosomes (24 d'origine paternelle + 24 d'origine maternelle).

(2) Le brassage interchromosomique produit $2^n = 2^{24} \approx \num{1,7} \times 10^{7}$, soit environ 17 millions de types de gamètes différents (sans compter les crossing-over).

(3) Si les gamètes étaient diploïdes, leur fusion doublerait le nombre de chromosomes à chaque génération ($48 \rightarrow 96 \rightarrow 192$\dots). La méiose ramène le stock à $n = 24$, et la fécondation restaure exactement $2n = 48$ dans le zygote : le caryotype de l'espèce est ainsi maintenu constant de génération en génération.
\end{corrige}

\newpage

\coursec{Activité d'intégration}

\courssub{Objectifs de l'activité}

À l'issue de cette activité d'intégration, tu devras être capable de :

\begin{itemize}
\item \textbf{ordonner} les différentes phases d'une méiose à partir de microphotographies, en justifiant chaque repère cytologique (savoir-faire n° 3 du programme) ;
\item \textbf{compléter, tracer et interpréter} la courbe de variation de la quantité d'ADN en fonction du temps au cours de la spermatogenèse (savoir-faire n° 4) ;
\item \textbf{identifier} une étape de la fécondation sur un document et la décrire (savoir-faire n° 5) ;
\item \textbf{expliquer} la nécessité du couple méiose--fécondation pour maintenir la constance du caryotype d'une génération à l'autre, tout en produisant la diversité génétique des individus (savoir-faire n° 1 et n° 2).
\end{itemize}

Cette activité se déroule en \textbf{groupes de 4 élèves} pendant \SI{55}{min}, suivie d'une restitution orale (\SI{5}{min} par groupe) et d'une correction collective avec une grille de notation de type Baccalauréat.

\courssub{Situation}

\textbf{Le mystère de la famille Ngo Bell : pourquoi trois enfants si différents ?}

M. et Mme Ngo Bell, habitants du quartier Nkolndongo à Yaoundé, se rendent au Centre Hospitalier Universitaire (CHU) de Yaoundé pour une consultation de génétique. Leur médecin traitant s'étonne, en plaisantant, que leurs trois enfants --- Chantal (17 ans), Junior (14 ans) et Sandrine (9 ans) --- soient « si différents alors qu'ils ont les mêmes parents » : teints légèrement différents, groupes sanguins différents, tailles et morphologies distinctes. Le biologiste du laboratoire d'anatomie pathologique réalise alors les caryotypes des trois enfants et constitue un dossier de quatre documents.

\textbf{Document 1 --- Les caryotypes des trois enfants.} Les trois caryotypes montrent chacun \textbf{$2n = 46$ chromosomes}, soit 22 paires d'autosomes et 1 paire de gonosomes (XX pour Chantal et Sandrine, XY pour Junior). Cependant, le biologiste souligne que les chromosomes homologues de chaque paire ne sont jamais identiques d'un enfant à l'autre : par exemple, pour la paire n° 1, Chantal possède les versions paternelle $P_1$ et maternelle $M_2$, Junior possède $P_2$ et $M_1$, et Sandrine $P_1$ et $M_1$. Les parents sont, eux aussi, à $2n = 46$ chromosomes.

\textbf{Document 2 --- Cellules de testicule de criquet en division.} Le laboratoire dispose d'une série de 6 microphotographies (légendées A, B, C, D, E, F) de cellules germinales mâles de criquet ($2n = 8$ chromosomes) observées au microscope optique, présentées dans le désordre :

\begin{center}
\begin{tabularx}{\linewidth}{|c|X|}\hline
\rowcolor{popblueL}\textbf{Photo} & \textbf{Description cytologique} \\ \hline
A & Deux lots de 4 chromosomes simples (à une chromatide) migrent vers les pôles opposés d'une cellule allongée. \\ \hline
B & 4 chromosomes doubles (à deux chromatides) s'alignent sur la plaque équatoriale ; les centromères des chromatides sœurs sont orientés vers les pôles opposés. \\ \hline
C & 4 paires de chromosomes homologues accolés (bivalents) se placent de part et d'autre du plan équatorial. \\ \hline
D & Les chromosomes sont condensés, les homologues sont appariés et des points de contact (chiasmas) sont visibles entre chromatides non sœurs. \\ \hline
E & Deux cellules-filles se forment, chacune contenant 4 chromosomes doubles ; la membrane nucléaire se reforme. \\ \hline
F & Les chromosomes homologues se séparent et migrent vers les pôles opposés ; chaque chromosome reste double. \\ \hline
\end{tabularx}
\end{center}

\textbf{Document 3 --- Courbe incomplète de la quantité d'ADN au cours de la spermatogenèse.} On mesure la quantité d'ADN par noyau (en picogrammes, pg) dans les cellules d'un tubule séminifère humain. La quantité de référence d'une cellule somatique en $G_1$ est $Q = \SI{6,6}{pg}$. La courbe fournie ne montre que le début du phénomène (phase de multiplication et phase d'accroissement, réplication incluse) ; il manque toute la partie correspondant aux deux divisions de méiose et à la spermiogenèse.

\begin{popfigure}
\begin{center}
\begin{tikzpicture}
\begin{axis}[width=0.85\linewidth,height=5.5cm,
xlabel={Temps (jours)},ylabel={Quantité d'ADN (pg)},
xmin=0,xmax=40,ymin=0,ymax=15,
xtick={0,5,10,15,20,25,30,35,40},ytick={0,3.3,6.6,9.9,13.2},
grid=both,grid style={popline!40}]
\addplot[popcurve,domain=0:8,samples=2]{6.6};
\addplot[popcurve,domain=8:14,samples=50]{6.6+(x-8)*1.1};
\addplot[popmuted,dashed,domain=14:40,samples=2]{13.2};
\node[font=\scriptsize,popdark] at (axis cs:4,7.6) {spermatogonies};
\node[font=\scriptsize,poporange] at (axis cs:11,9.2) {réplication};
\node[font=\scriptsize,popmuted] at (axis cs:28,14) {partie à compléter};
\end{axis}
\end{tikzpicture}
\end{center}
\legende{Courbe incomplète de la quantité d'ADN par noyau au cours de la spermatogenèse humaine. La partie en pointillés (à partir du jour 14) est à compléter. $Q = \SI{6,6}{pg}$ ; $2Q = \SI{13,2}{pg}$.}
\end{popfigure}

\textbf{Document 4 --- La fécondation chez l'oursin.} Une photographie au microscope électronique montre un spermatozoïde dont l'\textbf{acrosome} a émis un filament acrosomique au contact de la membrane vitelline de l'ovule ; un schéma annexe montre, quelques minutes plus tard, deux \textbf{pronucléus} (mâle et femelle, chacun à $n$ chromosomes) en voie de fusion au centre de l'œuf.

\textbf{Problème posé :} comment les trois enfants Ngo Bell peuvent-ils tous posséder exactement $2n = 46$ chromosomes comme leurs parents, tout en étant génétiquement différents les uns des autres ?

\courssub{Consignes}

\begin{enumerate}
\item \textbf{Ordonner la méiose.} Remets les six microphotographies du document 2 dans l'ordre chronologique de la méiose. Pour chaque photo, nomme la phase exacte (en précisant méiose I ou méiose II) et cite le repère cytologique qui justifie ton choix.
\item \textbf{Compléter la courbe d'ADN.} Reproduis et complète la courbe du document 3 jusqu'à la formation des spermatozoïdes. Annote chaque palier et chaque transition avec : le type cellulaire concerné (spermatogonie, spermatocyte I, spermatocyte II, spermatide, spermatozoïde), le nombre de chromosomes ($2n$ ou $n$) et l'événement responsable de chaque variation (réplication, première division, deuxième division).
\item \textbf{Identifier l'étape de fécondation.} Sur le document 4, identifie l'étape représentée sur la photographie et celle représentée sur le schéma. Explique pourquoi la fusion des pronucléus rétablit la diploïdie.
\item \textbf{Synthèse écrite.} En 15 lignes maximum, rédige une explication montrant comment les trois enfants Ngo Bell possèdent tous $2n = 46$ chromosomes tout en étant génétiquement différents. Ta réponse doit mobiliser obligatoirement : la réduction chromatique de la méiose, le brassage interchromosomique, le brassage intrachromosomique (crossing-over) et le hasard de la fécondation.
\end{enumerate}

\courssub{Ressources à mobiliser}

\begin{aretenir}[Les savoirs indispensables]
\begin{itemize}
\item \textbf{Méiose} : deux divisions successives précédées d'une \emph{seule} réplication de l'ADN. Méiose I (division réductionnelle) : séparation des chromosomes \emph{homologues} (prophase I avec appariement et crossing-over, métaphase I, anaphase I, télophase I). Méiose II (division équationnelle) : séparation des \emph{chromatides sœurs} (prophase II, métaphase II, anaphase II, télophase II).
\item \textbf{Quantités d'ADN} : spermatogonie $2n$, $Q$ ; après réplication, spermatocyte I $2n$, $2Q$ ; spermatocyte II $n$, $Q$ ; spermatide et spermatozoïde $n$, $Q/2$.
\item \textbf{Brassages} : interchromosomique (orientation aléatoire des bivalents en métaphase I : $2^n$ combinaisons, soit $2^{23} \approx \num{8,4}$ millions chez l'Homme) et intrachromosomique (crossing-over en prophase I).
\item \textbf{Fécondation} : reconnaissance gamétique, réaction acrosomique, fusion des membranes plasmiques, entrée du noyau mâle dans le cytoplasme de l'ovule, fusion des pronucléus ($n + n = 2n$), réactivation du métabolisme de l'œuf.
\end{itemize}
\end{aretenir}

\courssub{Démarche et corrigé commenté}

\begin{exoresolu}[Corrigé commenté de l'activité « Famille Ngo Bell »]
\textbf{Question 1.} Ordonner les six photos et justifier chaque classement.

\tcblower

\textbf{Question 1 --- Ordre chronologique : D $\rightarrow$ C $\rightarrow$ F $\rightarrow$ E $\rightarrow$ B $\rightarrow$ A.}

\begin{center}
\begin{tabularx}{\linewidth}{|c|l|X|}\hline
\rowcolor{popblueL}\textbf{Rang} & \textbf{Phase} & \textbf{Repère justificatif} \\ \hline
D & Prophase I & Chromosomes condensés, homologues appariés, \textbf{chiasmas} visibles (témoins des crossing-over). \\ \hline
C & Métaphase I & \textbf{Bivalents} (paires d'homologues) alignés de part et d'autre du plan équatorial. \\ \hline
F & Anaphase I & Séparation des \textbf{homologues} ; chaque chromosome qui migre est encore \textbf{double} (2 chromatides). \\ \hline
E & Télophase I & \textbf{Deux cellules-filles} à $n = 4$ chromosomes doubles ; membrane nucléaire reformée. \\ \hline
B & Métaphase II & 4 chromosomes doubles alignés sur la plaque équatoriale, centromères des chromatides sœurs orientés vers les pôles opposés. \\ \hline
A & Anaphase II & Migration de chromosomes \textbf{simples} (une chromatide) vers les pôles. \\ \hline
\end{tabularx}
\end{center}

\emph{Piège à éviter} : ne pas confondre anaphase I (F) et anaphase II (A) : en anaphase I, ce sont les homologues qui se séparent et les chromosomes restent doubles ; en anaphase II, ce sont les chromatides sœurs qui se séparent et les chromosomes deviennent simples. C'est d'ailleurs l'absence de prophase II et de télophase II dans la série qui doit t'alerter : la deuxième division suit immédiatement la première, sans réplication d'ADN intermédiaire.

\textbf{Question 2 --- Courbe complétée.} Les paliers et transitions sont :

\begin{itemize}
\item Jours 0 à 8 : \textbf{spermatogonies} ($2n$, $Q = \SI{6,6}{pg}$) --- mitoses de multiplication ;
\item Jours 8 à 14 : montée de $Q$ à $2Q = \SI{13,2}{pg}$ : \textbf{réplication de l'ADN} (phase S) pendant l'accroissement ; la cellule devient \textbf{spermatocyte I} ($2n$, $2Q$) ;
\item Vers le jour 20 : chute brutale de $2Q$ à $Q$ : \textbf{première division de méiose} $\rightarrow$ \textbf{spermatocytes II} ($n$, $Q$) ;
\item Vers le jour 27 : chute de $Q$ à $Q/2 = \SI{3,3}{pg}$ : \textbf{deuxième division de méiose} $\rightarrow$ \textbf{spermatides} ($n$, $Q/2$), qui se transforment ensuite en \textbf{spermatozoïdes} ($n$, $Q/2$) sans variation d'ADN (spermiogenèse).
\end{itemize}

\begin{popfigure}
\begin{center}
\begin{tikzpicture}
\begin{axis}[width=0.85\linewidth,height=5.5cm,
xlabel={Temps (jours)},ylabel={Quantité d'ADN (pg)},
xmin=0,xmax=40,ymin=0,ymax=15,
xtick={0,5,10,15,20,25,30,35,40},ytick={3.3,6.6,13.2},
yticklabels={$Q/2$,$Q$,$2Q$},
grid=both,grid style={popline!40}]
\addplot[popcurve,const plot] coordinates {(0,6.6)(8,6.6)};
\addplot[popcurve,domain=8:14,samples=50]{6.6+(x-8)*1.1};
\addplot[popcurve,const plot] coordinates {(14,13.2)(20,13.2)};
\addplot[popcurve,const plot] coordinates {(20,6.6)(27,6.6)};
\addplot[popcurve,const plot] coordinates {(27,3.3)(40,3.3)};
\node[font=\scriptsize,popdark] at (axis cs:4,7.6) {spermatogonies};
\node[font=\scriptsize,popdark] at (axis cs:17,14.2) {spermatocyte I};
\node[font=\scriptsize,popdark] at (axis cs:23.5,7.6) {spermatocytes II};
\node[font=\scriptsize,popdark] at (axis cs:34,4.3) {spermatides puis spermatozoïdes};
\node[font=\scriptsize,poporange] at (axis cs:11,9.5) {réplication};
\node[font=\scriptsize,poporange] at (axis cs:20,10.5) {méiose I};
\node[font=\scriptsize,poporange] at (axis cs:27,5.6) {méiose II};
\end{axis}
\end{tikzpicture}
\end{center}
\legende{Courbe complétée : variation de la quantité d'ADN par noyau au cours de la spermatogenèse. Une seule réplication ($Q \rightarrow 2Q$) suivie de deux divisions ($2Q \rightarrow Q \rightarrow Q/2$).}
\end{popfigure}

\emph{Lecture de la courbe} : chaque chute brutale correspond à une division (répartition de l'ADN dans deux cellules-filles) ; chaque montée progressive correspond à une réplication. On compte bien \textbf{une seule montée} pour \textbf{deux chutes} : c'est la signature de la méiose.

\textbf{Question 3 --- Fécondation.} La photographie montre la \textbf{réaction acrosomique} : au contact de la membrane vitelline, l'acrosome du spermatozoïde émet un filament acrosomique et libère des enzymes (dont la hyaluronidase et l'acrosine) qui dissolvent localement les enveloppes de l'ovule : c'est l'étape de \textbf{pénétration}. Le schéma montre une étape ultérieure : la \textbf{fusion des pronucléus} mâle ($n$) et femelle ($n$), c'est-à-dire la \textbf{caryogamie}. Cette fusion réunit dans un même noyau les deux lots haploïdes de chromosomes : $n + n = 2n$. La diploïdie est ainsi rétablie dans l'œuf (zygote), qui retrouve le nombre de chromosomes caractéristique de l'espèce et peut entamer sa segmentation.

\textbf{Question 4 --- Synthèse (modèle de réponse en 15 lignes).} Les parents Ngo Bell sont diploïdes ($2n = 46$). Lors de la gamétogenèse, la méiose comporte une seule réplication de l'ADN suivie de deux divisions : elle produit donc des gamètes haploïdes ($n = 23$), ce qui est la condition nécessaire pour que la fécondation ($n + n$) rétablisse $2n = 46$ chez chaque enfant, assurant la constance du caryotype de l'espèce humaine. Mais la méiose réalise aussi un double brassage génétique. Le brassage interchromosomique (orientation aléatoire des 23 bivalents en métaphase I) offre $2^{23} \approx \num{8,4}$ millions de combinaisons par parent ; le brassage intrachromosomique (crossing-over en prophase I) crée en outre des chromosomes recombinés, mosaïques de segments paternels et maternels. Chaque spermatozoïde et chaque ovule sont donc génétiquement uniques. Enfin, la fécondation associe au hasard un spermatozoïde et un ovule parmi ces millions de possibilités. Chantal, Junior et Sandrine reçoivent ainsi tous 46 chromosomes --- moitié du père, moitié de la mère --- mais des combinaisons d'allèles différentes (par exemple $P_1M_2$, $P_2M_1$, $P_1M_1$ pour la paire n° 1). Constance du nombre, diversité des combinaisons : voilà le mystère résolu.

\emph{Grille type Bac (sur 5 points)} : ordre correct des photos justifié (1,5 pt) ; courbe complétée avec types cellulaires et événements (1,5 pt) ; étapes de fécondation identifiées et diploïdie expliquée (1 pt) ; synthèse mobilisant les 4 notions exigées, rédigée en $\leq 15$ lignes (1 pt).
\end{exoresolu}

\courssub{Bilan de l'activité}

Cette activité t'a permis de mettre en évidence le \textbf{couple indissociable méiose--fécondation} :

\begin{itemize}
\item la \textbf{méiose} divise par deux le stock chromosomique ($2n \rightarrow n$) grâce à une seule réplication suivie de deux divisions ; sans elle, le nombre de chromosomes doublerait à chaque génération ;
\item la \textbf{fécondation} rétablit la diploïdie ($n + n = 2n$) par la fusion des pronucléus ; sans elle, les gamètes haploïdes resteraient sans descendance ;
\item ensemble, ces deux mécanismes assurent la \textbf{constance du caryotype} de l'espèce (ici $2n = 46$ chez l'Homme) tout en générant, par les brassages inter- et intrachromosomiques puis par le hasard de la rencontre des gamètes, une \textbf{diversité génétique quasi infinie} : c'est pourquoi Chantal, Junior et Sandrine Ngo Bell sont à la fois tous humains, tous enfants de leurs parents, et pourtant tous différents.
\end{itemize}

\begin{attention}[Erreurs fréquentes relevées pendant la correction]
Ne confonds jamais « quantité d'ADN » et « nombre de chromosomes » : après la réplication, le spermatocyte I possède toujours $2n$ chromosomes mais $2Q$ d'ADN (chaque chromosome est double, formé de deux chromatides sœurs). De même, la spermiogenèse (différenciation des spermatides en spermatozoïdes) ne modifie \emph{ni} le nombre de chromosomes \emph{ni} la quantité d'ADN : c'est une transformation morphologique (perte de cytoplasme, formation du flagelle et de l'acrosome), pas une division. Enfin, dans la rédaction, écris « pronucléus » (et non « pronuclei », qui est le pluriel latin) et précise toujours « méiose I » ou « méiose II » quand tu nommes une phase.
\end{attention}

\newpage

\coursec{Méthodes et exercices résolus}

\courssub{Méthode 1 : Expliquer la nécessité de la méiose dans la formation des gamètes}

\begin{methode}[Démontrer que la méiose est indispensable]
\begin{enumerate}
\item \textbf{Rappeler le point de départ} : les cellules de la lignée germinale sont \cle{diploïdes} ($2n$ chromosomes, soit $n$ paires de chromosomes homologues). Chez l'Homme, $2n = 46$.
\item \textbf{Raisonner par l'absurde} : si les gamètes étaient produits par mitose, ils resteraient à $2n$ ; à chaque fécondation, le nombre de chromosomes doublerait ($2n \rightarrow 4n \rightarrow 8n \ldots$), ce qui est incompatible avec la \cle{constance du caryotype} de l'espèce.
\item \textbf{Conclure sur le rôle de la méiose} : la méiose comporte \textbf{deux divisions successives précédées d'une seule réplication de l'ADN}. Elle produit 4 cellules \cle{haploïdes} ($n$) à partir d'une cellule diploïde ($2n$) : c'est une \textbf{division réductionnelle}.
\item \textbf{Ajouter le double bénéfice} : la méiose conserve le nombre de chromosomes d'une génération à l'autre (avec la fécondation) \emph{et} brasse l'information génétique (brassage interchromosomique en anaphase I, brassage intrachromosomique par crossing-over en prophase I), source de diversité.
\end{enumerate}
\textbf{Phrase type à rédiger} : « La méiose est nécessaire car elle ramène le stock chromosomique de $2n$ à $n$ dans les gamètes ; la fécondation rétablit ensuite la diploïdie, assurant la constance du caryotype de l'espèce. »
\end{methode}

\begin{exoresolu}[Pourquoi la méiose est-elle obligatoire ?]
Chez le Rat, la cellule-mère des spermatozoïdes possède $2n = 42$ chromosomes. Un élève affirme : « Les spermatozoïdes pourraient être produits par mitose, cela irait plus vite. » Montrer, par un raisonnement chiffré, que cette proposition est impossible, puis conclure sur le rôle de la méiose.
\tcblower
\textbf{Raisonnement par l'absurde (hypothèse de l'élève).}\\
Si les gamètes étaient formés par mitose, chaque spermatozoïde et chaque ovule conserverait $2n = 42$ chromosomes. À la fécondation, l'œuf posséderait :
\[ 42 + 42 = 84 \ \text{chromosomes} \ (4n). \]
À la génération suivante, on aurait $84 + 84 = 168$ chromosomes, puis $336$, etc. Le nombre de chromosomes doublerait à chaque génération : le caryotype de l'espèce ne serait plus constant, ce qui est contraire à l'observation (tous les Rats de l'espèce possèdent $2n = 42$ chromosomes).

\textbf{Rôle réel de la méiose.}\\
La méiose comporte une seule réplication de l'ADN suivie de deux divisions : une cellule à $2n = 42$ chromosomes donne 4 spermatides à $n = 21$ chromosomes. L'ovule est aussi à $n = 21$. À la fécondation :
\[ n + n = 21 + 21 = 42 = 2n \]
La diploïdie est rétablie et le caryotype de l'espèce est maintenu constant.

\textbf{Conclusion.} La méiose est une division réductionnelle \emph{indispensable} : elle produit des gamètes haploïdes ($n = 21$), ce qui, associé à la fécondation, garantit la constance du nombre de chromosomes ($2n = 42$) d'une génération à l'autre, tout en assurant le brassage génétique.
\end{exoresolu}

\courssub{Méthode 2 : Ordonner les phases d'une méiose}

\begin{methode}[Reconnaître et ordonner les phases sur des documents]
\begin{enumerate}
\item \textbf{Compter les divisions} : la méiose = 2 divisions (I : réductionnelle ; II : équationnelle). Ordre chronologique : \textbf{Prophase I $\rightarrow$ Métaphase I $\rightarrow$ Anaphase I $\rightarrow$ Télophase I $\rightarrow$ Prophase II $\rightarrow$ Métaphase II $\rightarrow$ Anaphase II $\rightarrow$ Télophase II}.
\item \textbf{Critères de reconnaissance de la division I} :
\begin{itemize}
\item Prophase I : chromosomes homologues appariés en \cle{bivalents}, éventuels chiasmas (crossing-over) ;
\item Métaphase I : bivalents alignés sur la plaque équatoriale ;
\item Anaphase I : séparation des \textbf{chromosomes homologues} (à 2 chromatides chacun) $\rightarrow$ c'est la phase réductionnelle.
\end{itemize}
\item \textbf{Critères de la division II} (semblable à une mitose) : chromosomes à 2 chromatides alignés seuls (métaphase II), puis séparation des \textbf{chromatides sœurs} (anaphase II).
\item \textbf{Vérifier le résultat} : 4 cellules haploïdes ($n$ chromosomes à 1 chromatide).
\end{enumerate}
\textbf{Astuce} : si l'on voit des chromosomes \emph{appariés par deux}, c'est forcément la division I ; si les chromosomes sont \emph{seuls mais doubles}, c'est la division II.
\end{methode}

\begin{exoresolu}[Remettre de l'ordre dans une méiose]
Sur une coupe de tubule séminifère de Rat ($2n = 6$ pour simplifier le document), on observe 4 cellules en division, photographiées au microscope :
\begin{itemize}
\item \textbf{Photo A} : 3 chromosomes doubles alignés seuls sur la plaque équatoriale ;
\item \textbf{Photo B} : 3 bivalents (chromosomes homologues appariés) sur la plaque équatoriale ;
\item \textbf{Photo C} : des chromosomes simples (à 1 chromatide) migrent vers les pôles ;
\item \textbf{Photo D} : des chromosomes doubles migrent vers les pôles, les homologues se séparant.
\end{itemize}
1. Identifier la phase représentée sur chaque photo. 2. Ordonner chronologiquement les photos.
\tcblower
\textbf{1. Identification.}\\
\textbf{Photo B} : les chromosomes homologues sont \emph{appariés en bivalents} sur la plaque équatoriale $\rightarrow$ \textbf{Métaphase I}.\\
\textbf{Photo D} : ce sont les \emph{chromosomes homologues} (encore doubles, à 2 chromatides) qui se séparent $\rightarrow$ \textbf{Anaphase I} (étape réductionnelle : passage de $2n$ à $n$).\\
\textbf{Photo A} : les chromosomes sont \emph{doubles mais seuls} (plus de bivalents) et alignés : la cellule est déjà haploïde ($n = 3$) $\rightarrow$ \textbf{Métaphase II}.\\
\textbf{Photo C} : ce sont les \emph{chromatides sœurs} qui se séparent ; les chromosomes migrants sont simples $\rightarrow$ \textbf{Anaphase II}.

\textbf{2. Ordre chronologique.}
\[ \boxed{B \ (\text{Métaphase I}) \rightarrow D \ (\text{Anaphase I}) \rightarrow A \ (\text{Métaphase II}) \rightarrow C \ (\text{Anaphase II})} \]

\textbf{Vérification.} En anaphase I, chaque pôle reçoit $n = 3$ chromosomes doubles ; en anaphase II, les chromatides sœurs se séparent : on obtient finalement 4 cellules à $n = 3$ chromosomes simples. Le compte est bon : la méiose est bien réductionnelle.
\end{exoresolu}

\courssub{Méthode 3 : Tracer et interpréter la courbe de la quantité d'ADN au cours de la spermatogenèse}

\begin{methode}[Construire la courbe $Q_{\text{ADN}} = f(t)$]
\begin{enumerate}
\item \textbf{Fixer la référence} : soit $q$ la quantité d'ADN d'une cellule haploïde (spermatozoïde). Une cellule diploïde non répliquée contient $2q$.
\item \textbf{Repérer les 4 étapes de la spermatogenèse} et les quantités d'ADN correspondantes :
\begin{itemize}
\item Phase de multiplication (spermatogonies, mitoses) : $2q$ ;
\item Phase de croissance (spermatocyte I) : \textbf{réplication} de l'ADN : $2q \rightarrow 4q$ ;
\item Méiose I (spermatocytes II) : $4q \rightarrow 2q$ ;
\item Méiose II (spermatides) : $2q \rightarrow q$ ;
\item Spermiogenèse (différenciation, sans division) : $q$ constant.
\end{itemize}
\item \textbf{Tracer} : temps en abscisse (étapes), quantité d'ADN en ordonnée ($q$, $2q$, $4q$). La courbe monte en rampe pendant la réplication, puis chute en deux paliers successifs ($4q \rightarrow 2q \rightarrow q$).
\item \textbf{Interpréter chaque saut} : la montée = réplication ; la première chute = méiose I ; la deuxième chute = méiose II.
\end{enumerate}
\end{methode}

\begin{exoresolu}[La courbe de l'ADN au cours de la spermatogenèse]
Chez un Mammifère, un spermatozoïde contient une quantité d'ADN égale à $q = \SI{3,2}{pg}$. 1. Donner la quantité d'ADN d'une spermatogonie, d'un spermatocyte I après croissance, d'un spermatocyte II et d'une spermatide. 2. Tracer l'allure de la courbe de variation de la quantité d'ADN au cours de la spermatogenèse et l'interpréter.
\tcblower
\textbf{1. Quantités d'ADN.} Avec $q = \SI{3,2}{pg}$ :
\begin{align*}
\text{Spermatogonie } (2n) &: 2q = \SI{6,4}{pg}\\
\text{Spermatocyte I après réplication} &: 4q = \SI{12,8}{pg}\\
\text{Spermatocyte II } (n, \text{ chromosomes doubles}) &: 2q = \SI{6,4}{pg}\\
\text{Spermatide puis spermatozoïde } (n) &: q = \SI{3,2}{pg}
\end{align*}

\textbf{2. Courbe et interprétation.}
\begin{popfigure}
\begin{center}
\begin{tikzpicture}
\begin{axis}[width=0.85\linewidth,height=6cm,
xlabel={Étapes de la spermatogenèse}, ylabel={Quantité d'ADN (pg)},
ymin=0,ymax=14, xmin=0, xmax=5,
xtick={0.5,1.5,2.5,3.5,4.5},
xticklabels={Multiplication,Croissance,Méiose I,Méiose II,Spermiogenèse},
x tick label style={font=\small,rotate=15,anchor=east},
ytick={3.2,6.4,12.8}, yticklabels={$q$,$2q$,$4q$}]
\addplot[popcurve,mark=*] coordinates {(0,6.4) (1,6.4) (2,12.8) (3,6.4) (4,3.2) (5,3.2)};
\end{axis}
\end{tikzpicture}
\legende{Variation de la quantité d'ADN par cellule au cours de la spermatogenèse : palier à $2q$ (spermatogonies), montée à $4q$ (réplication pendant la croissance), chute à $2q$ (méiose I), chute à $q$ (méiose II), palier à $q$ (spermiogenèse).}
\end{center}
\end{popfigure}

\textbf{Interprétation.} La montée de $2q$ à $4q$ traduit la \textbf{réplication de l'ADN} pendant la phase de croissance du spermatocyte I. La première chute ($4q \rightarrow 2q$) correspond à la \textbf{méiose I} (séparation des homologues) ; la seconde chute ($2q \rightarrow q$) correspond à la \textbf{méiose II} (séparation des chromatides). Pendant la spermiogenèse, la quantité d'ADN ne varie plus : il n'y a que différenciation. La spermatogenèse ramène donc bien la quantité d'ADN de $2q$ à $q$.
\end{exoresolu}

\courssub{Méthode 4 : Expliquer la nécessité de la fécondation dans la reproduction}

\begin{methode}[Justifier le rôle de la fécondation]
\begin{enumerate}
\item \textbf{Définir} : la fécondation est l'\textbf{union d'un gamète mâle ($n$) et d'un gamète femelle ($n$)} formant un œuf (zygote) \cle{diploïde} ($2n$).
\item \textbf{Justifier la restauration de la diploïdie} : seule la fusion des deux noyaux haploïdes (amphimixie) rétablit les $n$ paires de chromosomes homologues, condition du développement de l'embryon.
\item \textbf{Justifier la diversité} : la fécondation réunit deux lots génétiques issus de deux individus différents, eux-mêmes brassés par la méiose : chaque zygote est \textbf{génétiquement unique} (hors jumeaux monozygotes).
\item \textbf{Justifier l'activation} : l'entrée du spermatozoïde déclenche l'achèvement de la méiose de l'ovocyte II et la reprise des divisions (segmentation) : la fécondation \emph{déclenche} le développement.
\end{enumerate}
\textbf{Réflexe de rédaction} : toujours articuler méiose et fécondation : « la méiose sans fécondation appauvrirait le caryotype ; la fécondation sans méiose l'enrichirait indéfiniment ».
\end{methode}

\begin{exoresolu}[Pourquoi la fécondation est-elle indispensable ?]
Chez la Truie, $2n = 38$ chromosomes. L'ovocyte II bloqué en métaphase II contient $n = 19$ chromosomes. 1. Expliquer pourquoi l'ovocyte seul ne peut pas donner naissance à un porcelet viable dans la reproduction sexuée normale. 2. Donner deux rôles de la fécondation.
\tcblower
\textbf{1. Impossibilité d'un développement à partir de l'ovocyte seul.}\\
L'ovocyte II est \textbf{haploïde} ($n = 19$ chromosomes) : il ne possède qu'un seul chromosome de chaque paire d'homologues. Or le développement embryonnaire normal exige la \textbf{diploïdie} ($2n = 38$), c'est-à-dire la présence des deux lots d'allèles permettant l'expression équilibrée du génome. De plus, l'ovocyte II est \textbf{bloqué en métaphase II} : sans stimulus, sa méiose ne s'achève pas et aucune division de segmentation ne peut démarrer.

\textbf{2. Rôles de la fécondation.}\\
\emph{Rôle 1 — Restauration de la diploïdie} : la fusion du noyau du spermatozoïde ($n = 19$) avec celui de l'ovule ($n = 19$) donne un zygote à $19 + 19 = 38 = 2n$ chromosomes : le caryotype de l'espèce est rétabli.\\
\emph{Rôle 2 — Activation de l'œuf} : la pénétration du spermatozoïde déclenche l'achèvement de la méiose II (expulsion du 2\up{e} globule polaire) puis la segmentation.\\
\emph{Rôle 3 — Diversité génétique} : le zygote associe deux patrimoines génétiques différents, brassés par la méiose : il est unique.

\textbf{Conclusion.} La fécondation est nécessaire car elle rétablit la diploïdie ($2n = 38$), déclenche le développement de l'œuf et crée un individu génétiquement nouveau : elle est le complément indissociable de la méiose.
\end{exoresolu}

\courssub{Méthode 5 : Identifier les étapes de la fécondation sur des documents}

\begin{methode}[Reconnaître les étapes de la fécondation]
\begin{enumerate}
\item \textbf{Mémoriser la chronologie} : (1) contact et reconnaissance gamétique ; (2) \cle{réaction acrosomique} (enzymes de l'acrosome dissolvent la zone pellucide) ; (3) pénétration d'\textbf{un seul} spermatozoïde ; (4) \cle{réaction corticale} (la zone pellucide devient imperméable : blocage de la polyspermie) ; (5) achèvement de la méiose II de l'ovocyte et expulsion du 2\up{e} globule polaire ; (6) formation des deux \textbf{pronuclei} (mâle et femelle) ; (7) \cle{amphimixie} : fusion des patrimoines, zygote $2n$.
\item \textbf{Sur un document, repérer les indices} : spermatozoïdes nombreux autour de l'ovocyte = approche ; un spermatozoïde dans le cytoplasme = pénétration ; présence de deux globules polaires = méiose achevée ; deux pronuclei = fécondation en cours ; un seul noyau $2n$ = zygote formé.
\item \textbf{Rédiger dans l'ordre} en nommant précisément les structures (acrosome, zone pellucide, granules corticaux, pronuclei).
\end{enumerate}
\end{methode}

\begin{exoresolu}[Légender un document de fécondation]
Un document montre trois schémas d'un ovocyte de Lapine entouré de spermatozoïdes :
\begin{itemize}
\item \textbf{Schéma 1} : un spermatozoïde au contact de la zone pellucide, son acrosome libérant des enzymes ;
\item \textbf{Schéma 2} : un spermatozoïde (tête) dans le cytoplasme de l'ovocyte, la zone pellucide modifiée, les autres spermatozoïdes bloqués ;
\item \textbf{Schéma 3} : deux volumineux noyaux (pronuclei) dans le cytoplasme, un globule polaire expulsé.
\end{itemize}
1. Identifier l'étape représentée par chaque schéma. 2. Expliquer comment la polyspermie est évitée.
\tcblower
\textbf{1. Identification des étapes.}\\
\textbf{Schéma 1} : la libération d'enzymes par l'acrosome au contact de la zone pellucide correspond à la \textbf{réaction acrosomique} : les enzymes (hyaluronidase, acrosine) digèrent localement la zone pellucide pour permettre la pénétration.\\
\textbf{Schéma 2} : la tête d'un spermatozoïde est entrée dans l'ovocyte et la zone pellucide est devenue infranchissable pour les autres : c'est la \textbf{pénétration} suivie de la \textbf{réaction corticale} (les granules corticaux libèrent des enzymes qui durcissent la zone pellucide).\\
\textbf{Schéma 3} : la présence de deux pronuclei (mâle et femelle) et l'expulsion du 2\up{e} globule polaire montrent l'\textbf{achèvement de la méiose II} de l'ovocyte et la formation des pronuclei, juste avant l'\textbf{amphimixie} (fusion des deux patrimoines haploïdes donnant le noyau diploïde du zygote).

\textbf{2. Blocage de la polyspermie.}\\
Dès l'entrée du premier spermatozoïde, les \textbf{granules corticaux} de l'ovocyte déversent leur contenu enzymatique sous la membrane : la zone pellucide est \textbf{durcie et modifiée chimiquement}, devenant imperméable à tout autre spermatozoïde. Ainsi, un seul spermatozoïde féconde l'ovocyte, garantissant un zygote \textbf{diploïde} ($2n$) et non polyploïde.

\textbf{Conclusion.} La fécondation est une suite ordonnée d'événements : réaction acrosomique $\rightarrow$ pénétration $\rightarrow$ réaction corticale $\rightarrow$ achèvement de la méiose II $\rightarrow$ amphimixie, aboutissant au zygote $2n$.
\end{exoresolu}

\courssub{Méthode 6 : Comparer spermatogenèse et ovogenèse (lecture de documents)}

\begin{methode}[Construire une comparaison rigoureuse]
\begin{enumerate}
\item \textbf{Comparer selon des critères précis} : lieu (tubules séminifères du testicule / follicules de l'ovaire), période (continue de la puberté à la mort / cyclique, stock constitué avant la naissance), rendement (4 spermatozoïdes / 1 ovule + 2 à 3 globules polaires), durée (environ 74 jours chez l'Homme / plusieurs années avec blocages en prophase I puis métaphase II), nature des produits (cellules mobiles flagellées / cellule volumineuse riche en réserves).
\item \textbf{Dégager le point commun} : dans les deux cas, une méiose produit des gamètes haploïdes ($n$) à partir de cellules diploïdes ($2n$).
\item \textbf{Expliquer les différences} par leurs rôles : l'ovule, unique et riche en vitellus, nourrira le début du développement ; les spermatozoïdes, nombreux et mobiles, maximisent les chances de rencontre.
\end{enumerate}
\end{methode}

\begin{exoresolu}[Spermatogenèse et ovogenèse : le rendement]
Chez l'Homme, une spermatogonie en méiose donne 4 spermatozoïdes fonctionnels, alors qu'une ovogonie ne donne qu'un seul ovule. 1. Expliquer cette différence de rendement. 2. Justifier l'intérêt biologique de chaque modalité.
\tcblower
\textbf{1. Explication de la différence de rendement.}\\
Dans la \textbf{spermatogenèse}, les deux divisions de la méiose sont \textbf{équilibrées} : les cytokinèses partagent équitablement le cytoplasme ; un spermatocyte I donne 2 spermatocytes II, puis 4 spermatides qui se transforment toutes en 4 spermatozoïdes.\\
Dans l'\textbf{ovogenèse}, les cytokinèses sont \textbf{inégales} : à la méiose I, l'ovocyte I donne un gros ovocyte II et un petit 1\up{er} globule polaire ; à la méiose II, l'ovocyte II donne un ovule et un 2\up{e} globule polaire. Le cytoplasme et les réserves (vitellus) sont conservés dans une seule cellule : \textbf{1 seul ovule fonctionnel} ; les globules polaires dégénèrent.

\textbf{2. Intérêt biologique.}\\
\emph{Spermatogenèse prolifique} : produire des millions de spermatozoïdes mobiles compense la faible probabilité de rencontre avec l'ovocyte (un seul spermatozoïde sur des millions féconde).\\
\emph{Ovogenèse économe mais riche} : l'ovule unique concentre tout le cytoplasme et les réserves nutritives indispensables aux premières divisions de segmentation, avant l'implantation.

\textbf{Conclusion.} Les deux modalités illustrent une même méiose adaptée différemment : quantité et mobilité côté mâle, qualité et réserves côté femelle, les deux stratégies étant complémentaires pour le succès de la reproduction.
\end{exoresolu}

\begin{attention}[Les 5 erreurs qui coûtent des points au Bac]
\begin{enumerate}
\item \textbf{Confondre les deux anaphases} : en anaphase I, ce sont les \emph{chromosomes homologues} (à 2 chromatides) qui se séparent ; en anaphase II, ce sont les \emph{chromatides sœurs}. Écrire « séparation des chromatides en anaphase I » est une erreur grave.
\item \textbf{Oublier que la réplication de l'ADN a lieu une seule fois}, avant la méiose I (phase de croissance). Il n'y a \emph{pas} de réplication entre méiose I et méiose II : c'est ce qui explique le passage de $2n$ à $n$.
\item \textbf{Mal placer les paliers de la courbe d'ADN} : la spermatogonie est à $2q$ (pas à $q$ !), le spermatocyte I après croissance à $4q$, le spermatocyte II à $2q$, la spermatide à $q$. Vérifier que la courbe comporte exactement une montée et deux descentes.
\item \textbf{Dire que l'ovocyte II est un ovule} : l'ovocyte II n'achève sa méiose II (et ne devient ovule) qu'\emph{après} la pénétration du spermatozoïde. Chez la Femme, la « fécondation » a lieu sur un ovocyte II bloqué en métaphase II.
\item \textbf{Affirmer que plusieurs spermatozoïdes entrent dans l'ovocyte} : la réaction corticale bloque la polyspermie ; un seul spermatozoïde pénètre. Ne pas oublier non plus de citer l'\textbf{amphimixie} (fusion des pronuclei) comme acte final rétablissant la diploïdie.
\end{enumerate}
\end{attention}

\newpage

\coursec{Exercices}

\courssub{Je vérifie mes connaissances}

\begin{exercice}[QCM : les bons réflexes]
Pour chaque question, choisir la (ou les) bonne(s) réponse(s) en justifiant brièvement.

\textbf{1.} Une cellule humaine en fin de méiose I contient :
\begin{enumerate}[label=\alph*)]
\item 46 chromosomes à une chromatide ;
\item 23 chromosomes à deux chromatides ;
\item 23 chromosomes à une chromatide ;
\item 46 chromosomes à deux chromatides.
\end{enumerate}

\textbf{2.} Le crossing-over (enjambement) a lieu :
\begin{enumerate}[label=\alph*)]
\item en prophase I, entre chromosomes homologues ;
\item en prophase II, entre chromatides sœurs ;
\item en métaphase I, entre chromosomes non homologues ;
\item pendant la fécondation.
\end{enumerate}

\textbf{3.} L'acrosome du spermatozoïde :
\begin{enumerate}[label=\alph*)]
\item contient les mitochondries qui fournissent l'énergie ;
\item contient des enzymes qui dissolvent la zone pellucide ;
\item constitue le flagelle ;
\item contient le matériel génétique.
\end{enumerate}

\textbf{4.} Chez la femme, la fécondation a normalement lieu :
\begin{enumerate}[label=\alph*)]
\item dans l'utérus ;
\item dans l'ovaire ;
\item dans le tiers externe de la trompe de Fallope ;
\item dans le vagin.
\end{enumerate}

\textbf{5.} Une méiose complète produit, à partir d'une cellule mère :
\begin{enumerate}[label=\alph*)]
\item 2 cellules diploïdes ;
\item 4 cellules diploïdes ;
\item 4 cellules haploïdes génétiquement différentes ;
\item 4 cellules haploïdes génétiquement identiques.
\end{enumerate}
\end{exercice}

\begin{exercice}[Vrai ou faux, en justifiant]
Dire si chacune des affirmations suivantes est vraie ou fausse, puis justifier la réponse en une ou deux phrases.
\begin{enumerate}
\item « La réplication de l'ADN a lieu deux fois au cours de la méiose : avant la méiose I et avant la méiose II. »
\item « La spermatogenèse produit quatre spermatozoïdes fonctionnels à partir d'une spermatogonie, alors que l'ovogenèse ne produit qu'un seul ovocyte fonctionnel. »
\item « Chez la femme, l'ovogenèse débute à la puberté. »
\item « La fécondation rétablit la diploïdie dans la cellule-œuf. »
\item « Les globules polaires sont des gamètes mâles dégénérescents. »
\end{enumerate}
\end{exercice}

\begin{exercice}[Définitions éclair]
Définir en deux lignes maximum chacun des termes suivants, en employant un vocabulaire scientifique rigoureux : \cle{gonade}, \cle{méiose}, \cle{gamétogenèse}, \cle{fécondation}, \cle{cellule-œuf} (ou zygote), \cle{chromosomes homologues}.
\end{exercice}

\begin{exercice}[Calculs directs : ploïdie et ADN]
Chez le Rat, la cellule somatique contient $2n = 42$ chromosomes et une quantité d'ADN égale à $2Q$ en phase G1.
\begin{enumerate}
\item Donner le nombre de chromosomes et la quantité d'ADN (en fonction de $Q$) d'une spermatogonie en phase G1.
\item Mêmes questions pour : un spermatocyte I en fin de phase S ; un spermatocyte II ; une spermatide ; un spermatozoïde.
\item Combien de chromosomes contient la cellule-œuf issue de la fécondation chez le Rat ?
\item Une ovogonie de Rat subit la méiose. Combien d'ovocytes fonctionnels et combien de globules polaires sont théoriquement produits ?
\end{enumerate}
\end{exercice}

\courssub{Je m'entraîne}

\begin{exercice}[Ordonner les phases de la méiose]
On réalise des frottis de testicule de Rat et on observe au microscope photonique des cellules en division. Les descriptions suivantes correspondent à cinq clichés, dans le désordre :
\begin{itemize}
\item \textbf{Cliché A} : les chromosomes, à deux chromatides, se disposent par paires d'homologues sur le plan équatorial ;
\item \textbf{Cliché B} : les chromosomes homologues s'apparient intimement et des échanges de fragments sont visibles entre chromatides non sœurs ;
\item \textbf{Cliché C} : les chromosomes à une chromatide migrent vers les pôles opposés de la cellule ;
\item \textbf{Cliché D} : les chromosomes homologues se séparent et migrent chacun vers un pôle, chaque chromosome restant formé de deux chromatides ;
\item \textbf{Cliché E} : les chromosomes à deux chromatides s'alignent individuellement sur le plan équatorial.
\end{itemize}
\begin{enumerate}
\item Identifier la phase représentée par chaque cliché (nom de la phase et division concernée : méiose I ou méiose II).
\item Restituer l'ordre chronologique des cinq clichés en justifiant chaque choix.
\item Citer le cliché dont l'étape est à l'origine du brassage intrachromosomique, et préciser en quoi il consiste.
\end{enumerate}
\end{exercice}

\begin{exercice}[Courbe de l'ADN au cours de la spermatogenèse]
On mesure la quantité d'ADN par noyau dans différentes cellules prélevées dans un tube séminifère de Mammifère. On obtient trois valeurs : $Q$, $2Q$ et $4Q$.
\begin{enumerate}
\item Attribuer à chaque valeur le(s) type(s) cellulaire(s) correspondant(s) parmi : spermatogonie en G1, spermatocyte I après réplication, spermatocyte II, spermatide.
\item Tracer la courbe complète de variation de la quantité d'ADN en fonction du temps au cours de la spermatogenèse, depuis la spermatogonie en G1 jusqu'à la spermatide (axes légendés avec les grandeurs et les étapes).
\item Expliquer l'origine de chaque palier et de chaque chute de la courbe.
\end{enumerate}
\end{exercice}

\begin{exercice}[Spermatogenèse ou ovogenèse ?]
Recopier et compléter le tableau comparatif suivant :
\begin{center}
\begin{tabularx}{\linewidth}{|l|X|X|}
\hline
\rowcolor{popblueL} \textbf{Critère} & \textbf{Spermatogenèse} & \textbf{Ovogenèse} \\
\hline
Lieu & & \\
\hline
Début chronologique & & \\
\hline
Continuité dans le temps & & \\
\hline
Nombre de gamètes fonctionnels par cellule mère & & \\
\hline
Taille relative du gamète produit & & \\
\hline
\end{tabularx}
\end{center}
Expliquer ensuite, en trois ou quatre lignes, l'intérêt biologique de la formation des globules polaires au cours de l'ovogenèse.
\end{exercice}

\begin{exercice}[Les étapes de la fécondation]
Un document montre, chez l'oursin, quatre étapes schématisées : (1) un spermatozoïde au contact de la membrane vitelline, son acrosome libérant des enzymes ; (2) la traversée des enveloppes de l'ovule et la fusion des membranes des deux gamètes ; (3) l'entrée du noyau mâle dans le cytoplasme de l'ovule et le rapprochement des deux noyaux ; (4) la fusion des deux noyaux haploïdes.
\begin{enumerate}
\item Légender chaque étape à l'aide du vocabulaire exact (réaction acrosomique, pénétration, migration des pronuclei, caryogamie).
\item Restituer l'ordre chronologique des étapes.
\item Expliquer pourquoi la fécondation est dite « indispensable » à la reproduction sexuée : quelles sont ses deux conséquences fondamentales sur le plan chromosomique et sur le plan génétique ?
\end{enumerate}
\end{exercice}

\courssub{Je me prépare au Bac}

\begin{exercice}[Type Bac : microphotographies de divisions dans un testicule]
Au cours d'une séance de travaux pratiques au lycée, des élèves de Terminale C observent des coupes de testicule de Mammifère. Quatre microphotographies de cellules en division sont présentées dans le désordre :
\begin{itemize}
\item \textbf{Document 1} : cellule dont les chromosomes homologues appariés forment des figures en croix (chiasmas) bien visibles ;
\item \textbf{Document 2} : cellule dont les chromosomes à deux chromatides sont alignés par paires d'homologues de part et d'autre du plan équatorial ;
\item \textbf{Document 3} : cellule dont les chromosomes à une chromatide se séparent vers les deux pôles ;
\item \textbf{Document 4} : cellule dont les chromosomes à deux chromatides migrent vers les pôles, les homologues se séparant.
\end{itemize}
On précise que l'espèce étudiée possède $2n = 8$ chromosomes.

\textbf{Questions :}
\begin{enumerate}
\item Identifier la phase représentée dans chaque document et préciser s'il s'agit de la méiose I ou de la méiose II. (4 pts)
\item Restituer l'ordre chronologique de ces quatre étapes et justifier l'ordre choisi. (3 pts)
\item Indiquer le nombre de chromosomes et le nombre de chromatides par chromosome observés dans le document 2, puis dans le document 3. (2 pts)
\item Expliquer, à partir des documents, pourquoi la méiose est nécessaire à la formation des gamètes chez les Mammifères. (2 pts)
\item Schématiser le document 4 pour cette espèce ($2n = 8$) en représentant deux paires de chromosomes homologues par des couleurs différentes. (3 pts)
\end{enumerate}
\emph{Barème indicatif : 14 pts — durée conseillée : 45 min.}
\end{exercice}

\begin{exercice}[Type Bac : variation de la quantité d'ADN dans un tube séminifère]
Dans un laboratoire de l'hôpital universitaire de Yaoundé, on étudie la spermatogenèse chez un Mammifère. On mesure par cytophotométrie la quantité d'ADN contenue dans le noyau de cellules prélevées à différents niveaux d'un tube séminifère. Les résultats sont consignés dans le tableau ci-dessous :

\begin{center}
\begin{tabularx}{\linewidth}{|l|X|X|X|X|X|}
\hline
\rowcolor{popblueL} \textbf{Cellule} & A & B & C & D & E \\
\hline
\textbf{Quantité d'ADN} & $2Q$ & $4Q$ & $2Q$ & $Q$ & $Q$ \\
\hline
\textbf{Chromosomes} & $2n$ à 1 chromatide & $2n$ à 2 chromatides & $n$ à 2 chromatides & $n$ à 1 chromatide & $n$ à 1 chromatide, cellule allongée avec flagelle \\
\hline
\end{tabularx}
\end{center}

\textbf{Questions :}
\begin{enumerate}
\item Identifier chacune des cellules A, B, C, D et E en justifiant à partir des données du tableau. (5 pts)
\item Tracer la courbe de variation de la quantité d'ADN en fonction du temps au cours de la spermatogenèse, de la spermatogonie (phase G1) au spermatozoïde. On fera figurer sur l'axe des abscisses les étapes : phase S, méiose I, méiose II, spermiogenèse. (4 pts)
\item Expliquer l'origine du passage de $2Q$ à $4Q$, puis les deux chutes successives de $4Q$ à $2Q$ et de $2Q$ à $Q$. (3 pts)
\item La spermiogenèse s'accompagne-elle d'une variation de la quantité d'ADN ? Justifier. (2 pts)
\item Un produit bloque spécifiquement la réplication de l'ADN. Prévoir, en justifiant, la quantité d'ADN des cellules qui tenteraient alors d'entrer en méiose. (2 pts)
\end{enumerate}
\emph{Barème indicatif : 16 pts — durée conseillée : 60 min.}
\end{exercice}

\begin{exercice}[Type Bac : trisomie 21, méiose et fécondation]
Un couple consulte un centre de génétique médicale de Douala après la naissance d'un enfant atteint de trisomie 21. L'analyse des caryotypes donne les résultats suivants : le père présente 46 chromosomes, la mère présente 46 chromosomes, l'enfant présente 47 chromosomes avec trois exemplaires du chromosome 21. On rappelle que la méiose comporte deux divisions et que, normalement, les chromosomes homologues se séparent en anaphase I et les chromatides sœurs en anaphase II.

\textbf{Questions :}
\begin{enumerate}
\item Rappeler ce qu'est un caryotype et indiquer la formule chromosomique de l'enfant (on notera $2n + 1 = 47$). (2 pts)
\item Expliquer, à l'aide des connaissances sur la méiose, le mécanisme de la \cle{non-disjonction} chromosomique pouvant conduire à un gamète porteur de deux exemplaires du chromosome 21. (4 pts)
\item Préciser les deux étapes de la méiose où cette non-disjonction peut se produire, et indiquer, pour chaque cas, la proportion de gamètes anormaux produits. (4 pts)
\item Expliquer comment la fécondation d'un tel gamète anormal par un gamète normal conduit au caryotype de l'enfant. On s'appuiera sur un schéma simple. (3 pts)
\item Déduire de cette étude pourquoi la méiose et la fécondation forment un couple indissociable dans la reproduction sexuée, tant pour la stabilité du nombre de chromosomes que pour la diversité génétique. (3 pts)
\end{enumerate}
\emph{Barème indicatif : 16 pts — durée conseillée : 60 min.}
\end{exercice}

\newpage

\coursec{Corrigés des exercices}

\begin{corrige}[Exercice 1 — QCM : les bons réflexes]
\textbf{1. Réponse b : 23 chromosomes à deux chromatides.} La méiose I sépare les chromosomes \emph{homologues} (disjonction) mais pas les chromatides sœurs : chaque cellule fille reçoit donc un chromosome de chaque paire, encore formé de deux chromatides. Chez l'Homme : $n = 23$ chromosomes à 2 chromatides.

\textbf{2. Réponse a : en prophase I, entre chromosomes homologues.} Le crossing-over s'effectue au stade pachytène de la prophase I, entre chromatides \emph{non sœurs} de chromosomes homologues appariés en bivalents. Il n'existe ni en prophase II ni à la fécondation.

\textbf{3. Réponse b : contient des enzymes qui dissolvent la zone pellucide.} L'acrosome, vésicule dérivée de l'appareil de Golgi, libère des enzymes (hyaluronidase, acrosine) lors de la réaction acrosomique. Les mitochondries forment la pièce intermédiaire, le matériel génétique est dans le noyau, le flagelle est le filament terminal.

\textbf{4. Réponse c : dans le tiers externe de la trompe de Fallope.} C'est dans la portion ampullaire (tiers externe) de l'oviducte que spermatozoïdes et ovocyte II se rencontrent, environ 12 à \SI{24}{h} après l'ovulation.

\textbf{5. Réponse c : 4 cellules haploïdes génétiquement différentes.} Une méiose = deux divisions successives donnant 4 cellules à $n$ chromosomes. Elles sont génétiquement différentes à cause du brassage interchromosomique (disjonction aléatoire des homologues en anaphase I) et du brassage intrachromosomique (crossing-over en prophase I).
\end{corrige}

\begin{attention}[Piège classique]
Confondre « fin de méiose I » et « fin de méiose II » : à la fin de la méiose I, les chromosomes ont encore \textbf{deux} chromatides ; ce n'est qu'après l'anaphase II qu'ils n'en ont plus qu'une.
\end{attention}

\begin{corrige}[Exercice 2 — Vrai ou faux, en justifiant]
\begin{enumerate}
\item \textbf{Faux.} La réplication de l'ADN n'a lieu qu'\textbf{une seule fois}, en phase S précédant la méiose I (passage de $2Q$ à $4Q$). Il n'y a pas de phase S entre la méiose I et la méiose II : c'est précisément ce qui permet le passage de la diploïdie à l'haploïdie.
\item \textbf{Vrai.} Les divisions de la spermatogenèse sont équilibrées : une spermatogonie (via un spermatocyte I) donne 4 spermatides, toutes transformées en spermatozoïdes fonctionnels. Dans l'ovogenèse, les divisions sont inégales : l'ovocyte I donne un ovocyte II et un globule polaire, puis l'ovocyte II donne un seul ovule (ou ovocyte fécondé) et un globule polaire ; les globules polaires dégénèrent.
\item \textbf{Faux.} L'ovogenèse débute \emph{in utero}, pendant la vie fœtale : les ovogonies entrent en méiose et les ovocytes I restent bloqués en prophase I (stade dictyotène) jusqu'à la puberté. La puberté marque seulement la \emph{reprise} cyclique de la méiose, pas son début.
\item \textbf{Vrai.} La caryogamie (fusion des deux noyaux haploïdes, $n + n$) rétablit la diploïdie ($2n$) dans la cellule-œuf : le nombre de chromosomes caractéristique de l'espèce est ainsi maintenu de génération en génération.
\item \textbf{Faux.} Les globules polaires sont de petites cellules \emph{femelles}, produites par les divisions inégales de l'ovogenèse ; haploïdes mais quasi dépourvues de cytoplasme, elles ne sont pas fonctionnelles et dégénèrent.
\end{enumerate}
\end{corrige}

\begin{corrige}[Exercice 3 — Définitions éclair]
\begin{itemize}
\item \cle{Gonade} : glande génitale (testicule chez le mâle, ovaire chez la femelle) qui produit les gamètes (fonction reproductrice) et sécrète les hormones sexuelles (fonction endocrine).
\item \cle{Méiose} : succession de deux divisions cellulaires précédées d'une seule réplication de l'ADN, qui fait passer une cellule diploïde ($2n$) à quatre cellules haploïdes ($n$) génétiquement différentes.
\item \cle{Gamétogenèse} : ensemble des processus de formation des gamètes dans les gonades, comprenant la méiose puis la différenciation (spermatogenèse chez le mâle, ovogenèse chez la femelle).
\item \cle{Fécondation} : union d'un gamète mâle et d'un gamète femelle, comportant la reconnaissance, la pénétration puis la fusion des noyaux haploïdes (caryogamie), aboutissant à la cellule-œuf diploïde.
\item \cle{Cellule-œuf} (zygote) : cellule diploïde ($2n$) issue de la fécondation, première cellule du nouvel individu, dont toutes les cellules de l'organisme dériveront par mitoses.
\item \cle{Chromosomes homologues} : paire de chromosomes de même taille, même forme et portant les mêmes gènes (mais pas forcément les mêmes allèles), l'un d'origine paternelle, l'autre d'origine maternelle.
\end{itemize}
\end{corrige}

\begin{corrige}[Exercice 4 — Calculs directs : ploïdie et ADN]
Données : $2n = 42$ chromosomes, donc $n = 21$ ; quantité d'ADN en G1 : $2Q$.

\begin{enumerate}
\item \textbf{Spermatogonie en G1} : c'est une cellule diploïde non répliquée : $2n = 42$ chromosomes à 1 chromatide, quantité d'ADN $= 2Q$.
\item \textbf{Autres étapes} :
\begin{itemize}
\item \textbf{Spermatocyte I en fin de phase S} : l'ADN a été répliqué mais aucune division n'a eu lieu : $2n = 42$ chromosomes à 2 chromatides, ADN $= 4Q$.
\item \textbf{Spermatocyte II} : issu de la méiose I (disjonction des homologues) : $n = 21$ chromosomes à 2 chromatides, ADN $= 2Q$.
\item \textbf{Spermatide} : issue de la méiose II (séparation des chromatides) : $n = 21$ chromosomes à 1 chromatide, ADN $= Q$.
\item \textbf{Spermatozoïde} : la spermiogenèse est une différenciation sans division : $n = 21$ chromosomes à 1 chromatide, ADN $= Q$.
\end{itemize}
\item \textbf{Cellule-œuf} : la fécondation réunit deux gamètes haploïdes : $n + n = 2n = 42$ chromosomes (ADN $= 2Q$ en G1).
\item \textbf{Ovogenèse} : une ovogonie (via un ovocyte I) produit théoriquement \textbf{1 seul ovocyte fonctionnel} et \textbf{3 globules polaires} (2 issus de l'ovocyte I — en réalité le premier globule polaire peut se diviser — et 1 issu de l'ovocyte II), qui dégénèrent.
\end{enumerate}
\end{corrige}

\begin{attention}[Vérification des calculs]
Contrôle de cohérence : $4Q \xrightarrow{\text{méiose I}} 2 \times 2Q \xrightarrow{\text{méiose II}} 4 \times Q$. La quantité totale d'ADN est conservée à chaque division : $4Q = 2 \times 2Q = 4 \times Q$. Si vos valeurs ne vérifient pas cette conservation, il y a une erreur.
\end{attention}

\begin{corrige}[Exercice 5 — Ordonner les phases de la méiose]
\textbf{1. Identification des clichés :}
\begin{itemize}
\item \textbf{Cliché A} : paires d'homologues (bivalents) sur le plan équatorial $\Rightarrow$ \textbf{métaphase I}.
\item \textbf{Cliché B} : appariement intime des homologues (synapsis) et échanges entre chromatides non sœurs $\Rightarrow$ \textbf{prophase I}.
\item \textbf{Cliché C} : chromosomes à \emph{une} chromatide migrant vers les pôles $\Rightarrow$ \textbf{anaphase II} (séparation des chromatides sœurs).
\item \textbf{Cliché D} : homologues à \emph{deux} chromatides se séparant $\Rightarrow$ \textbf{anaphase I} (disjonction).
\item \textbf{Cliché E} : chromosomes à deux chromatides alignés \emph{individuellement} (et non par paires) $\Rightarrow$ \textbf{métaphase II}.
\end{itemize}

\textbf{2. Ordre chronologique :}
\[ \text{B (prophase I)} \rightarrow \text{A (métaphase I)} \rightarrow \text{D (anaphase I)} \rightarrow \text{E (métaphase II)} \rightarrow \text{C (anaphase II)} \]
Justification : la méiose I précède toujours la méiose II ; au sein de chaque division, l'ordre est prophase $\rightarrow$ métaphase $\rightarrow$ anaphase. Le critère discriminant entre I et II est le nombre de chromatides des chromosomes qui se séparent (2 en anaphase I, 1 en anaphase II) et la disposition équatoriale (par paires en métaphase I, individuelle en métaphase II).

\textbf{3. Brassage intrachromosomique} : il est provoqué par l'étape du \textbf{cliché B} (prophase I). Il consiste en des \cle{crossing-over} : échanges réciproques de fragments de chromatides non sœurs entre chromosomes homologues au niveau des chiasmas. Il crée des chromosomes recombinés, porteurs de combinaisons alléliques nouvelles (mélange d'allèles paternels et maternels sur un même chromosome).
\end{corrige}

\begin{corrige}[Exercice 6 — Courbe de l'ADN au cours de la spermatogenèse]
\textbf{1. Attribution des valeurs :}
\begin{center}
\begin{tabularx}{\linewidth}{|l|X|}
\hline
\rowcolor{popblueL} \textbf{Quantité d'ADN} & \textbf{Type(s) cellulaire(s)} \\
\hline
$2Q$ & spermatogonie en G1 ; spermatocyte II (après méiose I) \\
\hline
$4Q$ & spermatocyte I après réplication (fin de phase S jusqu'à la fin de la méiose I) \\
\hline
$Q$ & spermatide (et spermatozoïde) \\
\hline
\end{tabularx}
\end{center}

\textbf{2. Courbe complète :}
\begin{popfigure}
\begin{tikzpicture}
\begin{axis}[
  width=\linewidth, height=6.5cm,
  xlabel={Étapes de la spermatogenèse}, ylabel={Quantité d'ADN par noyau},
  xmin=0, xmax=10, ymin=0, ymax=5,
  xtick={1,3,5,7,9},
  xticklabels={G1, phase S, méiose I, méiose II, spermatide},
  ytick={1,2,4}, yticklabels={$Q$, $2Q$, $4Q$},
  grid=major, grid style={popline!40},
]
\addplot[popcurve, very thick] coordinates {(0,2) (2,2) (2,2) (4,4) (4,4) (6,4) (6,2) (8,2) (8,1) (10,1)};
\addplot[popcurve, dashed] coordinates {(2,2) (4,4)};
\end{axis}
\end{tikzpicture}
\legende{Variation de la quantité d'ADN par noyau au cours de la spermatogenèse : palier à $2Q$ (G1), montée progressive à $4Q$ (phase S), palier à $4Q$ (spermatocyte I), chute à $2Q$ (fin de méiose I), palier à $2Q$ (spermatocyte II), chute à $Q$ (fin de méiose II), palier à $Q$ (spermatide).}
\end{popfigure}

\textbf{3. Origine des paliers et des chutes :}
\begin{itemize}
\item \textbf{Palier $2Q$} : spermatogonie en G1, chromosomes à une chromatide.
\item \textbf{Montée $2Q \rightarrow 4Q$} : réplication semi-conservative de l'ADN pendant la phase S ; chaque chromosome devient formé de deux chromatides sœurs identiques.
\item \textbf{Palier $4Q$} : spermatocyte I en méiose I (chromosomes à deux chromatides, cellule encore unique).
\item \textbf{Chute $4Q \rightarrow 2Q$} : première division (méiose I) : répartition des homologues dans deux cellules filles ; chaque spermatocyte II reçoit la moitié de l'ADN.
\item \textbf{Palier $2Q$} : spermatocyte II (pas de phase S entre les deux divisions).
\item \textbf{Chute $2Q \rightarrow Q$} : deuxième division (méiose II) : séparation des chromatides sœurs ; chaque spermatide reçoit $n$ chromosomes à une chromatide.
\end{itemize}
\end{corrige}

\begin{corrige}[Exercice 7 — Spermatogenèse ou ovogenèse ?]
\textbf{Tableau complété :}
\begin{center}
\begin{tabularx}{\linewidth}{|l|X|X|}
\hline
\rowcolor{popblueL} \textbf{Critère} & \textbf{Spermatogenèse} & \textbf{Ovogenèse} \\
\hline
Lieu & Tubes séminifères des testicules & Follicules ovariens (ovaire) \\
\hline
Début chronologique & À la puberté & Pendant la vie fœtale (blocage en prophase I), reprise à la puberté \\
\hline
Continuité dans le temps & Continue, de la puberté à la mort & Cyclique (un ovocyte par cycle) et interrompue à la ménopause ; stock limité d'ovocytes \\
\hline
Nombre de gamètes fonctionnels par cellule mère & 4 spermatozoïdes fonctionnels & 1 seul ovocyte fonctionnel (+ 3 globules polaires dégénérescents) \\
\hline
Taille relative du gamète produit & Très petit, mobile (flagelle), peu de cytoplasme & Très gros, immobile, cytoplasme abondant riche en réserves (vitellus) \\
\hline
\end{tabularx}
\end{center}

\textbf{Intérêt biologique des globules polaires :} les divisions inégales de l'ovogenèse concentrent la quasi-totalité du cytoplasme et des réserves nutritives dans une seule cellule, l'ovocyte, au détriment des globules polaires qui restent minuscules et dégénèrent. Ainsi, tout en accomplissant la méiose (réduction chromosomique indispensable), la femelle produit un gamète volumineux, doté des réserves (vitellus) et du matériel cytoplasmique (ribosomes, ARNm, mitochondries) nécessaires aux premières divisions de l'embryon, avant que celui-ci ne puisse s'alimenter par lui-même.
\end{corrige}

\begin{corrige}[Exercice 8 — Les étapes de la fécondation]
\textbf{1. Légendes :}
\begin{itemize}
\item \textbf{(1)} \cle{Réaction acrosomique} : l'acrosome libère ses enzymes (hyaluronidase, acrosine) qui digèrent localement les enveloppes de l'ovule.
\item \textbf{(2)} \cle{Pénétration} : traversée des enveloppes (membrane vitelline chez l'oursin, zone pellucide chez les Mammifères) et fusion des membranes plasmiques des deux gamètes.
\item \textbf{(3)} \cle{Migration des pronuclei} : le noyau mâle (pronucleus mâle) pénètre dans le cytoplasme et se rapproche du pronucleus femelle.
\item \textbf{(4)} \cle{Caryogamie} : fusion des deux noyaux haploïdes ($n + n$) en un noyau diploïde ($2n$) : la cellule-œuf est formée.
\end{itemize}

\textbf{2. Ordre chronologique :} (1) réaction acrosomique $\rightarrow$ (2) pénétration $\rightarrow$ (3) migration des pronuclei $\rightarrow$ (4) caryogamie.

\textbf{3. Pourquoi la fécondation est-elle indispensable ?}
\begin{itemize}
\item \textbf{Sur le plan chromosomique} : la caryogamie rétablit la diploïdie ($n + n = 2n$). Associée à la méiose qui avait réduit le nombre de chromosomes de moitié, elle assure la \emph{stabilité du nombre de chromosomes} de l'espèce au fil des générations (sans elle, ce nombre doublerait à chaque génération).
\item \textbf{Sur le plan génétique} : la fécondation réunit dans une même cellule deux lots chromosomiques d'origines différentes (paternel et maternel). Combinée aux brassages méiotiques, elle crée un \emph{patrimoine génétique nouveau et unique} : c'est une source majeure de la diversité génétique intra-spécifique.
\end{itemize}
\end{corrige}

\begin{corrige}[Exercice 9 — Type Bac : microphotographies de divisions dans un testicule]
\textbf{1. Identification (4 pts — 1 pt par document) :}
\begin{itemize}
\item \textbf{Document 1} : chiasmas et homologues appariés $\Rightarrow$ \textbf{prophase I} (stade diplotène/diacinèse).
\item \textbf{Document 2} : bivalents alignés de part et d'autre du plan équatorial $\Rightarrow$ \textbf{métaphase I}.
\item \textbf{Document 3} : chromosomes à une chromatide se séparant $\Rightarrow$ \textbf{anaphase II}.
\item \textbf{Document 4} : homologues à deux chromatides se séparant $\Rightarrow$ \textbf{anaphase I}.
\end{itemize}

\textbf{2. Ordre chronologique (3 pts) :} Document 1 $\rightarrow$ Document 2 $\rightarrow$ Document 4 $\rightarrow$ Document 3. Justification : au sein de la méiose I, la prophase (appariement, chiasmas) précède la métaphase (alignement des bivalents), elle-même suivie de l'anaphase (disjonction des homologues) ; la méiose II (anaphase II, séparation des chromatides) ne peut avoir lieu qu'après l'achèvement de la méiose I.

\textbf{3. Nombres (2 pts) :}
\begin{itemize}
\item Document 2 (métaphase I) : $2n = 8$ chromosomes, chacun à \textbf{2 chromatides} (soit 16 chromatides, 4 bivalents).
\item Document 3 (anaphase II) : dans chaque cellule, $n = 4$ chromosomes se séparant, chacun passant à \textbf{1 chromatide} (8 chromosomes à 1 chromatide si l'on compte les deux lots en migration).
\end{itemize}

\textbf{4. Nécessité de la méiose (2 pts) :} les documents montrent qu'une seule réplication est suivie de deux divisions : la méiose I sépare les homologues (doc. 4), la méiose II sépare les chromatides (doc. 3). On passe ainsi d'une cellule diploïde ($2n = 8$) à des cellules haploïdes ($n = 4$). Cette réduction de moitié du stock chromosomique est indispensable : sans elle, la fécondation doublerait le nombre de chromosomes à chaque génération. De plus, les chiasmas (doc. 1) et la disjonction aléatoire des homologues (doc. 4) brassent les allèles, garantissant la diversité génétique des gamètes.

\textbf{5. Schéma du document 4 (3 pts) :}
\begin{popfigure}
\begin{tikzpicture}[scale=0.9]
\draw[popline, dashed] (0,-2.6) -- (0,2.6);
\fill[popmuted] (0,2.6) circle (0.08) node[above]{\small pôle};
\fill[popmuted] (0,-2.6) circle (0.08) node[below]{\small pôle};
% Homologues paire 1 (grands) : bleu vers le haut, orange vers le bas
\draw[very thick, popblue] (-1.6,0.5) -- (-1.2,0.9) -- (-0.8,0.5);
\draw[very thick, popblue] (-1.6,0.9) -- (-1.2,0.5) -- (-0.8,0.9);
\draw[very thick, poporange] (-1.6,-0.5) -- (-1.2,-0.9) -- (-0.8,-0.5);
\draw[very thick, poporange] (-1.6,-0.9) -- (-1.2,-0.5) -- (-0.8,-0.9);
% Homologues paire 2 (petits) : bleu vers le bas, orange vers le haut (disjonction indépendante)
\draw[very thick, popblue] (0.9,-0.45) -- (1.2,-0.7) -- (1.5,-0.45);
\draw[very thick, popblue] (0.9,-0.7) -- (1.2,-0.45) -- (1.5,-0.7);
\draw[very thick, poporange] (0.9,0.45) -- (1.2,0.7) -- (1.5,0.45);
\draw[very thick, poporange] (0.9,0.7) -- (1.2,0.45) -- (1.5,0.7);
\draw[->, popdark] (-1.2,1.1) -- (-1.2,2.1);
\draw[->, popdark] (-1.2,-1.1) -- (-1.2,-2.1);
\draw[->, popdark] (1.2,1.0) -- (1.2,2.1);
\draw[->, popdark] (1.2,-1.0) -- (1.2,-2.1);
\end{tikzpicture}
\legende{Anaphase I chez une espèce à $2n = 8$ (2 paires représentées) : les chromosomes homologues (bleu = origine paternelle, orange = origine maternelle), encore formés de deux chromatides, migrent vers des pôles opposés. La répartition des couleurs est aléatoire (brassage interchromosomique).}
\end{popfigure}
\end{corrige}

\begin{attention}[Points de vigilance — barème]
Le correcteur attend : le \textbf{nombre de chromatides} comme critère d'identification (obligatoire pour distinguer anaphase I de anaphase II) ; la justification de l'ordre par la \textbf{logique interne de chaque division} ; un schéma \textbf{légendé} (couleurs = origine parentale, flèches = sens de migration, chromosomes en V ou en X selon la phase). Oublier de préciser « méiose I » ou « méiose II » fait perdre la moitié des points de la question 1.
\end{attention}

\begin{corrige}[Exercice 10 — Type Bac : variation de la quantité d'ADN dans un tube séminifère]
\textbf{1. Identification des cellules (5 pts — 1 pt par cellule) :}
\begin{itemize}
\item \textbf{A} : $2n$ chromosomes à 1 chromatide, $2Q$ $\Rightarrow$ \textbf{spermatogonie en G1} (cellule diploïde avant réplication).
\item \textbf{B} : $2n$ chromosomes à 2 chromatides, $4Q$ $\Rightarrow$ \textbf{spermatocyte I} (après phase S, en méiose I).
\item \textbf{C} : $n$ chromosomes à 2 chromatides, $2Q$ $\Rightarrow$ \textbf{spermatocyte II} (après méiose I).
\item \textbf{D} : $n$ chromosomes à 1 chromatide, $Q$ $\Rightarrow$ \textbf{spermatide} (après méiose II).
\item \textbf{E} : $n$, $Q$, cellule allongée à flagelle $\Rightarrow$ \textbf{spermatozoïde} (spermatide différenciée par spermiogenèse).
\end{itemize}

\textbf{2. Courbe (4 pts) :}
\begin{popfigure}
\begin{tikzpicture}
\begin{axis}[
  width=\linewidth, height=6.5cm,
  xlabel={Étapes}, ylabel={Quantité d'ADN par noyau},
  xmin=0, xmax=11, ymin=0, ymax=5,
  xtick={1,3,5,7,9,10.4},
  xticklabels={G1, phase S, méiose I, méiose II, spermio\-genèse, spermatozoïde},
  x tick label style={font=\small},
  ytick={1,2,4}, yticklabels={$Q$, $2Q$, $4Q$},
  grid=major, grid style={popline!40},
]
\addplot[popcurve, very thick] coordinates {(0,2) (2,2) (4,4) (6,4) (6,2) (8,2) (8,1) (11,1)};
\addplot[popcurve, dashed] coordinates {(2,2) (4,4)};
\end{axis}
\end{tikzpicture}
\legende{Quantité d'ADN par noyau au cours de la spermatogenèse : $2Q$ (spermatogonie G1) $\rightarrow$ $4Q$ (fin de phase S) $\rightarrow$ $2Q$ (fin de méiose I) $\rightarrow$ $Q$ (fin de méiose II) ; la spermiogenèse laisse la quantité d'ADN inchangée.}
\end{popfigure}

\textbf{3. Origine des variations (3 pts) :}
\begin{itemize}
\item $2Q \rightarrow 4Q$ : \textbf{réplication semi-conservative de l'ADN} en phase S ; chaque chromosome passe de 1 à 2 chromatides sœurs, sans changement du nombre de chromosomes.
\item $4Q \rightarrow 2Q$ : \textbf{méiose I} : la disjonction des homologues répartit le stock d'ADN en deux spermatocytes II ($n$ chromosomes à 2 chromatides).
\item $2Q \rightarrow Q$ : \textbf{méiose II} : la séparation des chromatides sœurs répartit à nouveau l'ADN en deux spermatides ($n$ chromosomes à 1 chromatide).
\end{itemize}

\textbf{4. Spermiogenèse et ADN (2 pts) :} \textbf{Non}, la spermiogenèse ne modifie pas la quantité d'ADN (palier à $Q$ entre D et E). C'est une \emph{différenciation} sans division ni réplication : condensation du noyau, formation de l'acrosome (à partir de l'appareil de Golgi), du flagelle et de la gaine de mitochondries, perte de cytoplasme.

\textbf{5. Blocage de la réplication (2 pts) :} si la phase S est bloquée, les cellules entreraient en méiose avec seulement $2Q$ d'ADN et des chromosomes à \textbf{une seule chromatide}. La méiose II deviendrait impossible (pas de chromatides sœurs à séparer) et les gamètes obtenus seraient anormaux (voire la méiose serait bloquée dès la prophase I, les homologues ne pouvant former de bivalents fonctionnels). Quantité attendue à l'entrée en méiose : $2Q$ au lieu de $4Q$.
\end{corrige}

\begin{attention}[Points de vigilance — barème]
Exigences du correcteur : axes \textbf{légendés avec les grandeurs} (quantité d'ADN en fonction du temps/des étapes) ; justification de chaque identification par \textbf{deux} indices du tableau (ploidie ET nombre de chromatides) ; citer explicitement « pas de phase S entre méiose I et méiose II » pour expliquer le palier à $2Q$.
\end{attention}

\begin{corrige}[Exercice 11 — Type Bac : trisomie 21, méiose et fécondation]
\textbf{1. Caryotype et formule chromosomique (2 pts) :} le \cle{caryotype} est la représentation ordonnée (par taille et forme décroissantes) de l'ensemble des chromosomes d'une cellule, réalisée en métaphase. L'enfant présente trois exemplaires du chromosome 21 : sa formule est $2n + 1 = 47$ chromosomes (trisomie 21).

\textbf{2. Mécanisme de la non-disjonction (4 pts) :} normalement, en anaphase I, les deux chromosomes homologues 21 (l'un paternel, l'autre maternel) se séparent et migrent vers des pôles opposés ; en anaphase II, les deux chromatides sœurs de chaque chromosome se séparent. La \cle{non-disjonction} est le défaut de cette séparation : soit les deux homologues 21 migrent ensemble vers le même pôle en anaphase I, soit les deux chromatides sœurs du chromosome 21 ne se séparent pas en anaphase II. Il en résulte un gamète à $n + 1 = 24$ chromosomes, porteur de \textbf{deux exemplaires} du chromosome 21 (et un gamète à $n - 1 = 22$, dépourvu de chromosome 21).

\textbf{3. Étapes possibles et proportions (4 pts) :}
\begin{itemize}
\item \textbf{Non-disjonction en anaphase I} : les deux homologues partent dans la même cellule. Après méiose II, \textbf{les 4 gamètes sont anormaux} : 2 gamètes à $n + 1$ (deux chromosomes 21) et 2 gamètes à $n - 1$ (aucun chromosome 21). Proportion : $100\,\%$ de gamètes anormaux, dont $50\,\%$ à $n+1$.
\item \textbf{Non-disjonction en anaphase II} : la méiose I est normale ; un seul des deux spermatocytes II (ou l'ovocyte II) est concerné. On obtient \textbf{2 gamètes normaux} ($n$), 1 gamète à $n + 1$ et 1 gamète à $n - 1$. Proportion : $50\,\%$ de gamètes anormaux, dont $25\,\%$ à $n+1$.
\end{itemize}

\textbf{4. Fécondation et caryotype de l'enfant (3 pts) :}
\begin{popfigure}
\begin{tikzpicture}[scale=0.95]
\node[draw, popblue, fill=popblueL, rounded corners, minimum width=3.6cm, minimum height=1cm] (g1) at (0,2) {gamète anormal : $n+1 = 24$};
\node[draw, poporange, fill=poporangeL, rounded corners, minimum width=3.6cm, minimum height=1cm] (g2) at (6.5,2) {gamète normal : $n = 23$};
\node[draw, popgreen, fill=popgreenL, rounded corners, minimum width=4.4cm, minimum height=1cm] (z) at (3.25,0) {cellule-œuf : $2n+1 = 47$};
\draw[->, very thick, popdark] (g1) -- (z);
\draw[->, very thick, popdark] (g2) -- (z);
\node[popdark] at (3.25,-1.1) {\small 2 exemplaires du chr. 21 (un parent) $+$ 1 exemplaire (autre parent) $=$ 3 exemplaires};
\end{tikzpicture}
\legende{La fécondation d'un gamète à $n+1$ (deux chromosomes 21) par un gamète normal à $n$ (un chromosome 21) additionne les stocks chromosomiques : $24 + 23 = 47$ chromosomes, dont trois chromosomes 21.}
\end{popfigure}

La caryogamie additionne les deux lots haploïdes : $(n+1) + n = 2n + 1 = 47$. L'enfant possède donc trois chromosomes 21 dans toutes ses cellules, d'où le phénotype de trisomie 21.

\textbf{5. Méiose et fécondation : un couple indissociable (3 pts) :}
\begin{itemize}
\item \textbf{Stabilité du nombre de chromosomes} : la méiose divise par deux le stock chromosomique ($2n \rightarrow n$) ; la fécondation le rétablit ($n + n = 2n$). L'une sans l'autre entraînerait soit un doublement, soit un appauvrissement du caryotype à chaque génération — cette étude montre d'ailleurs qu'une simple erreur d'une unité ($n+1$) a des conséquences majeures.
\item \textbf{Diversité génétique} : la méiose brasse les allèles (crossing-over, disjonction aléatoire : $2^{23}$ combinaisons chez l'Homme) ; la fécondation, aléatoire, réunit deux gamètes parmi des millions de combinaisons possibles. Ensemble, ils garantissent à la fois la \emph{constance de l'espèce} et l'\emph{unicité de chaque individu}.
\end{itemize}
\end{corrige}

\begin{attention}[Points de vigilance — barème]
Ne pas écrire « la trisomie 21 est due à un chromosome surnuméraire apparu à la fécondation » : l'anomalie est \textbf{méiotique} (non-disjonction), la fécondation ne fait que l'additionner. Exigé : distinguer clairement les proportions selon l'étape de non-disjonction ($100\,\%$ de gamètes anormaux si anaphase I, $50\,\%$ si anaphase II) et schématiser la fécondation avec les formules $n+1$ et $n$ explicitées.
\end{attention}

\newpage

\coursec{Fiche bilan}

\begin{popfigure}
\begin{tikzpicture}[mindmap, grow cyclic, every node/.style={concept}, text width=2.6cm, align=center,
  root concept/.append style={concept color=popblue, font=\bfseries, text=white, minimum size=3.2cm},
  level 1/.append style={level distance=4.6cm, sibling angle=72, font=\small\bfseries, text width=2.4cm},
  level 2/.append style={level distance=2.6cm, font=\scriptsize, text width=2.0cm}]
\node[root concept]{Reproduction sexuée chez les Mammifères}
  child[concept color=poporange]{ node {Gonades : structure et rôles}
    child[concept color=poporangeL]{ node {Testicules : tubes séminifères, cellules de Leydig et de Sertoli} }
    child[concept color=poporangeL]{ node {Ovaires : follicules, corps jaune} }
  }
  child[concept color=popgreen]{ node {Méiose}
    child[concept color=popgreenL]{ node {1\iere{} division : réduction ($2n \to n$)} }
    child[concept color=popgreenL]{ node {2\ieme{} division : équationnelle} }
    child[concept color=popgreenL]{ node {Brassage inter- et intrachromosomique} }
  }
  child[concept color=poppurple]{ node {Spermatogenèse}
    child[concept color=poppurpleL]{ node {Continue, dès la puberté} }
    child[concept color=poppurpleL]{ node {1 spermatocyte I $\to$ 4 spermatozoïdes} }
  }
  child[concept color=poppink]{ node {Ovogenèse}
    child[concept color=poppinkL]{ node {Cyclique, bloquée en prophase I puis métaphase II} }
    child[concept color=poppinkL]{ node {1 ovocyte I $\to$ 1 ovule + 3 globules polaires} }
  }
  child[concept color=popgold]{ node {Fécondation}
    child[concept color=popgoldL]{ node {Réaction acrosomique, fusion des membranes} }
    child[concept color=popgoldL]{ node {Caryogamie : restauration de $2n$} }
  };
\end{tikzpicture}
\legende{Carte mentale de la leçon : les gonades produisent les gamètes par gamétogenèse (qui comporte la méiose) ; la fécondation restaure la diploïdie et associe deux lots parentaux de chromosomes.}
\end{popfigure}

\begin{bilan}
\textbf{Définitions clés}
\begin{itemize}
  \item \cle{Gonade} : glande génitale (testicule ou ovaire) à double fonction : \textbf{exocrine} (production des gamètes) et \textbf{endocrine} (sécrétion des hormones sexuelles).
  \item \cle{Méiose} : succession de deux divisions cellulaires précédées d'une seule réplication de l'ADN ; elle fait passer une cellule de $2n$ chromosomes à $n$ chromosomes (division \textbf{réductionnelle} puis division \textbf{équationnelle}).
  \item \cle{Gamétogenèse} : formation des gamètes dans les gonades : \textbf{spermatogenèse} (testicules) et \textbf{ovogenèse} (ovaires).
  \item \cle{Fécondation} : fusion d'un spermatozoïde ($n$) et d'un ovule ($n$) donnant une cellule-\oe uf ou \textbf{zygote} ($2n$).
  \item \cle{Brassage génétique} : \textbf{interchromosomique} (migration aléatoire des chromosomes à l'anaphase I) et \textbf{intrachromosomique} (crossing-over en prophase I) ; sources de la diversité génétique.
\end{itemize}

\textbf{À connaître par c\oe ur}
\begin{center}
\begin{tabularx}{\linewidth}{|l|X|X|}
\hline
\rowcolor{popblueL} \textbf{Critère} & \textbf{Spermatogenèse} & \textbf{Ovogenèse} \\
\hline
Lieu & Tubes séminifères du testicule & Follicules de l'ovaire \\
\hline
Déroulement & Continue, de la puberté à la mort & Cyclique ; débutée avant la naissance, bloquée en prophase I puis en métaphase II \\
\hline
Bilan par cellule I & 1 spermatocyte I $\to$ \textbf{4 spermatozoïdes} & 1 ovocyte I $\to$ \textbf{1 ovule} + 3 globules polaires \\
\hline
Quantité d'ADN & $Q \to 2Q \to Q \to Q/2$ & $Q \to 2Q \to Q \to Q/2$ \\
\hline
\end{tabularx}
\end{center}

\textbf{Étapes de la fécondation (ordre exigible)}
\begin{enumerate}
  \item Approche et contact gamète--ovule (réaction acrosomique : enzymes de l'acrosome).
  \item Franchissement des enveloppes de l'ovule et fusion des membranes.
  \item Pénétration du contenu spermatique ; blocage de la polyspermie.
  \item Achèvement de la méiose II de l'ovocyte et formation des deux \textbf{pronuclei}.
  \item \textbf{Caryogamie} : fusion des pronuclei $\to$ zygote $2n$.
\end{enumerate}

\textbf{Méthodes express}
\begin{itemize}
  \item Courbe d'ADN : repérer la phase S ($Q \to 2Q$), la méiose I ($2Q \to Q$), la méiose II ($Q \to Q/2$). Une seule réplication pour deux divisions.
  \item Ordonner la méiose : prophase I (chromosomes bivalents, crossing-over), métaphase I (plaques de bivalents), anaphase I (séparation des chromosomes \emph{à deux chromatides}), puis division II (séparation des chromatides).
  \item Justifier la nécessité du couple méiose--fécondation : la méiose divise par deux le stock chromosomique, la fécondation le restaure ; la diploïdie est ainsi constante d'une génération à l'autre.
\end{itemize}
\end{bilan}

\begin{aretenir}[Checklist avant le Bac]
\begin{itemize}
  \item[$\square$] Je sais définir gonade, méiose, gamétogenèse, fécondation, caryogamie.
  \item[$\square$] Je connais la structure du testicule (tubes séminifères, cellules de Sertoli et de Leydig) et de l'ovaire (follicules, corps jaune).
  \item[$\square$] Je sais que la méiose comporte deux divisions précédées d'une seule réplication de l'ADN.
  \item[$\square$] Je sais ordonner et décrire les phases de la méiose (prophase I avec crossing-over, anaphase I réductionnelle).
  \item[$\square$] Je distingue brassage interchromosomique et intrachromosomique et je sais citer leur étape.
  \item[$\square$] Je sais tracer et interpréter la courbe $Q \to 2Q \to Q \to Q/2$ de la quantité d'ADN au cours de la spermatogenèse.
  \item[$\square$] Je connais le bilan comparé : 4 spermatozoïdes contre 1 ovule + 3 globules polaires.
  \item[$\square$] Je sais que l'ovogenèse est bloquée en prophase I (naissance) puis en métaphase II (ovulation).
  \item[$\square$] Je sais identifier sur document les étapes de la fécondation (réaction acrosomique, pronuclei, caryogamie).
  \item[$\square$] Je sais expliquer pourquoi méiose et fécondation maintiennent la constance de $2n$ d'une génération à l'autre.
  \item[$\square$] Je n'écris jamais que la méiose « divise les chromosomes en deux » : je précise qu'elle réduit le \emph{nombre de lots} chromosomiques de $2n$ à $n$.
  \item[$\square$] Je légende mes schémas avec un vocabulaire exact (bivalent, chromatide, acrosome, pronucleus).
\end{itemize}
\end{aretenir}

\findecours

\end{document}
